% Rédiger un résumé de texte — Français Tle Tech — Cours signé OnBuch+
% Programme officiel MINESEC. Compiler avec : tectonic N22-rediger-resume-texte.tex
\newcommand{\DOCMATIERE}{Français}
\newcommand{\DOCNIVEAU}{Tle Tech}
\newcommand{\DOCMODULE}{Français – classe de Terminale technique}
\newcommand{\DOCLECON}{N22}
\newcommand{\DOCTITRE}{Rédiger un résumé de texte}
% OnBuch+ — préambule « cours signés » ENRICHI (compilé avec tectonic / XeLaTeX)
% Système visuel « Pop » d'OnBuch+ : fond crème (jamais blanc), bordures encre
% épaisses, ombres « dures » décalées, titres Archivo Black en capitales.
% Version MATHÉMATIQUES : graphes (pgfplots/tikz), boîtes pédagogiques
% (définition, propriété, méthode, attention, le savais-tu, exercice résolu,
% exercice, TP, bilan), page de garde et signature OnBuch+ ancrée en bas.
%
% Macros attendues dans main.tex AVANT \input{preamble} :
%   \DOCMATIERE  \DOCNIVEAU  \DOCMODULE  \DOCLECON (numéro)  \DOCTITRE
\documentclass[11pt]{article}

\usepackage{fontspec}
\usepackage[french]{babel}
\usepackage[a4paper,margin=2cm,top=3.4cm,bottom=2.8cm,headheight=1.4cm,footskip=1.3cm]{geometry}
\usepackage{amsmath,amssymb,amsfonts}
\usepackage{mathtools} % \xmapsto, \coloneqq... souvent utilisés par les modèles
\usepackage{cancel} % simplifications barrées (bilans de forces)
% \abs{z} (module, valeur absolue) et \norm{u} (norme), utilisés par les modèles
\providecommand{\abs}{}\let\abs\relax\DeclarePairedDelimiter{\abs}{\lvert}{\rvert}
\providecommand{\norm}{}\let\norm\relax\DeclarePairedDelimiter{\norm}{\lVert}{\rVert}
\usepackage[table]{xcolor} % [table] : fournit aussi \rowcolors (alternance de lignes)
\usepackage{pagecolor}
\usepackage{tikz}
\usetikzlibrary{arrows.meta,positioning,calc,shapes.geometric,shapes.misc,shapes.arrows,decorations.pathmorphing,decorations.pathreplacing,decorations.markings,patterns,shadows,mindmap,trees,fit,backgrounds,matrix,intersections,angles,quotes}
\usepackage{pgfplots}
\usepackage[version=4]{mhchem} % équations nucléaires \ce{^{235}_{92}U}
\usepackage[european,siunitx]{circuitikz} % schémas électriques
\pgfplotsset{compat=1.18}
\usepgfplotslibrary{fillbetween,groupplots} % aires entre courbes (calcul intégral)
\usepackage{enumitem}
\usepackage{fancyhdr}
% [most] charge toutes les bibliothèques tcolorbox (listings, minted, raster,
% documentation...) et gonfle inutilement la mémoire TeX sur un document long
% (~300 boîtes) : on ne charge que ce qui est réellement utilisé.
\usepackage[skins,breakable]{tcolorbox}
\usepackage{graphicx}
\usepackage{colortbl}
\usepackage{booktabs}
\usepackage{tabularx}
\usepackage{multirow}
\usepackage{diagbox} % case d'en-tête diagonale des tableaux à double entrée
% Colonnes extensibles centrée / gauche / droite (C, L, R), souvent utilisées par les modèles
\newcolumntype{C}{>{\centering\arraybackslash}X}
\newcolumntype{L}{>{\raggedright\arraybackslash}X}
\newcolumntype{R}{>{\raggedleft\arraybackslash}X}
\usepackage{array}
\usepackage{multicol}
\usepackage{siunitx}
% parse-numbers=false : accepte aussi \SI{2,5 \times 10^{-3}}{...}
\sisetup{locale=FR,output-decimal-marker={,},group-minimum-digits=4,parse-numbers=false}
\DeclareSIUnit\ohm{\ensuremath{\symup{\Omega}}} % U+2126 absent de la police
\DeclareSIUnit\division{div} % réglages d'oscilloscope (ms/div, V/div)
\DeclareSIUnit\tr{tr} % tours (vitesse de rotation, ex. \SI{1500}{\tr\per\minute})
\DeclareSIUnit\molar{M} % concentration molaire (ex. \SI{3}{\molar}, \SI{1}{\micro\molar})
\DeclareSIUnit\an{an} % année (le modèle confond parfois avec \year, un compteur TeX) — ex. \SI{5}{\kilo\meter\per\an}
\DeclareSIUnit\FCFA{FCFA} % franc CFA (ex. \SI{500}{\FCFA\per\kilo\gram})
\DeclareSIUnit\degreeBrix{\textdegree Brix} % teneur en sucre d'un sirop/jus (réfractométrie)
\DeclareSIUnit\month{mois} % le modèle utilise parfois \month, un compteur TeX, comme unité de durée
\DeclareSIUnit\microliter{\micro\liter} % le modèle écrit parfois \microliter en un seul token
\DeclareSIUnit\million{millions} % dénombrement (ex. \SI{15}{\million\per\mL} spermatozoïdes)
\usepackage{qrcode}
% \degree : commande fréquemment utilisée par les modèles pour un angle ou une
% température en dehors de \SI{}{} (non fournie par les paquets chargés).
\providecommand{\degree}{^{\circ}}
% \qed (fin de démonstration), utilisé par les modèles sans amsthm
\providecommand{\qed}{\hfill$\square$}
% \barre : double barre des valeurs interdites dans un tableau de variations (tkz-tab)
\providecommand{\barre}{\ensuremath{\|}}
% environnement proof (amsthm non chargé) utilisé par les modèles
\ifdefined\proof\else\newenvironment{proof}[1][Démonstration]{\par\noindent\textit{#1.}\ }{\qed\par}\fi
\usepackage{lastpage}
\usepackage{hyperref}
\hypersetup{hidelinks}

% ============================================================
% Palette « Pop » OnBuch+ — mêmes hex que la classe P (plus_ui.dart)
% ============================================================
\definecolor{poporange}{HTML}{F59321}
\definecolor{popdark}{HTML}{E07A0C}
\definecolor{popcream}{HTML}{FFF6E8}
\definecolor{popink}{HTML}{1C1714}
\definecolor{poppurple}{HTML}{7A5AE0}
\definecolor{popgreen}{HTML}{2E9E6A}
\definecolor{poppink}{HTML}{E0457B}
\definecolor{popblue}{HTML}{2F7DE1}
\definecolor{popgold}{HTML}{C98A12}
\definecolor{popmuted}{HTML}{8A7F76}
\definecolor{popline}{HTML}{EADFD2}
\definecolor{popgray}{HTML}{8A7F76}
\colorlet{popgrey}{popgray}
% Alias tolérants (le modèle nomme parfois les couleurs différemment)
\colorlet{popwhite}{white}
\colorlet{popblack}{popink}
\colorlet{popred}{poppink}
\colorlet{popyellow}{popgold}
% Teintes claires (fonds de tableaux/encadrés) — le modèle nomme parfois
% les couleurs avec un suffixe L (ex. popblueL) sans les avoir définies.
\colorlet{poporangeL}{poporange!20}
\colorlet{popdarkL}{popdark!20}
\colorlet{poppurpleL}{poppurple!20}
\colorlet{popgreenL}{popgreen!20}
\colorlet{poppinkL}{poppink!20}
\colorlet{popblueL}{popblue!20}
\colorlet{popgoldL}{popgold!20}
\colorlet{popredL}{popred!20}
\colorlet{popgrayL}{popgray!20}
% Teintes claires pour fonds de tableaux / zones
\colorlet{popblueL}{popblue!10!white}
\colorlet{poporangeL}{poporange!14!white}
\colorlet{popgreenL}{popgreen!12!white}
\colorlet{poppurpleL}{poppurple!12!white}
\colorlet{poppinkL}{poppink!12!white}
\colorlet{popgoldL}{popgold!14!white}

\pagecolor{popcream}

% ============================================================
% Typographie
% ============================================================
\defaultfontfeatures{Ligatures=TeX}
\setmainfont{PlusJakartaSans-Medium.ttf}[Path=../../../fonts/,BoldFont=PlusJakartaSans-Bold.ttf]
% Maths en sans-serif (Fira Math) avec chiffres et lettres droites de Plus
% Jakarta, pour que les formules se fondent dans le texte.
\usepackage[math-style=ISO,bold-style=ISO]{unicode-math}
\setmathfont{FiraMath-Regular.otf}[Path=../../../fonts/]
\setmathfont{PlusJakartaSans-Medium.ttf}[Path=../../../fonts/,range={up/num,up/{Latin,latin},\mathup}]
\usepackage{icomma} % virgule décimale sans espace en mode math
% Caractères mathématiques Unicode absents des polices de texte (Plus Jakarta,
% Archivo Black) : rendus par leur équivalent mathématique, en texte comme en
% formule — sinon ils s'affichent en carré vide (ex. « Arithmétique dans ℤ »).
\usepackage{newunicodechar}
\newunicodechar{ℝ}{\ensuremath{\mathbb{R}}}
\newunicodechar{ℤ}{\ensuremath{\mathbb{Z}}}
\newunicodechar{ℂ}{\ensuremath{\mathbb{C}}}
\newunicodechar{ℕ}{\ensuremath{\mathbb{N}}}
\newunicodechar{ℚ}{\ensuremath{\mathbb{Q}}}
\newunicodechar{𝔻}{\ensuremath{\mathbb{D}}}
\newunicodechar{α}{\ensuremath{\alpha}}
\newunicodechar{β}{\ensuremath{\beta}}
\newunicodechar{γ}{\ensuremath{\gamma}}
\newunicodechar{δ}{\ensuremath{\delta}}
\newunicodechar{ε}{\ensuremath{\varepsilon}}
\newunicodechar{θ}{\ensuremath{\theta}}
\newunicodechar{λ}{\ensuremath{\lambda}}
\newunicodechar{μ}{\ensuremath{\mu}}
\newunicodechar{σ}{\ensuremath{\sigma}}
\newunicodechar{τ}{\ensuremath{\tau}}
\newunicodechar{φ}{\ensuremath{\varphi}}
\newunicodechar{ψ}{\ensuremath{\psi}}
\newunicodechar{ω}{\ensuremath{\omega}}
\newunicodechar{Σ}{\ensuremath{\Sigma}}
% majuscules produites par \MakeUppercase dans les titres (α-aminés -> Α)
\newunicodechar{Α}{\ensuremath{\alpha}}
\newunicodechar{Β}{\ensuremath{\beta}}
\newunicodechar{Γ}{\ensuremath{\Gamma}}
\newunicodechar{Δ}{\ensuremath{\Delta}}
\newunicodechar{Ε}{\ensuremath{\varepsilon}}
\newunicodechar{Θ}{\ensuremath{\Theta}}
\newunicodechar{Λ}{\ensuremath{\Lambda}}
\newunicodechar{Μ}{\ensuremath{\mu}}
\newunicodechar{Τ}{\ensuremath{\tau}}
\newunicodechar{Φ}{\ensuremath{\Phi}}
\newunicodechar{Ψ}{\ensuremath{\Psi}}
\newunicodechar{Ω}{\ensuremath{\Omega}}
\newunicodechar{↦}{\ensuremath{\mapsto}}
\newunicodechar{∈}{\ensuremath{\in}}
\newunicodechar{∉}{\ensuremath{\notin}}
\newunicodechar{∧}{\ensuremath{\wedge}}
\newunicodechar{⇔}{\ensuremath{\Leftrightarrow}}
\newunicodechar{⟺}{\ensuremath{\Longleftrightarrow}}
\newunicodechar{⇒}{\ensuremath{\Rightarrow}}
\newunicodechar{⟹}{\ensuremath{\Longrightarrow}}
\newunicodechar{∘}{\ensuremath{\circ}}
\newunicodechar{∪}{\ensuremath{\cup}}
\newunicodechar{∩}{\ensuremath{\cap}}
\newunicodechar{≡}{\ensuremath{\equiv}}
\newunicodechar{⊂}{\ensuremath{\subset}}
\newunicodechar{⊥}{\ensuremath{\perp}}
\newunicodechar{∃}{\ensuremath{\exists}}
\newunicodechar{∀}{\ensuremath{\forall}}
\newunicodechar{∛}{\ensuremath{\sqrt[3]{\,}}}
\newunicodechar{∜}{\ensuremath{\sqrt[4]{\,}}}
\newunicodechar{⃗}{\ensuremath{{}^{\to}}}
\newunicodechar{ⁿ}{\ifmmode^{n}\else\textsuperscript{\ensuremath{n}}\fi}
\newunicodechar{ˣ}{\ifmmode^{x}\else\textsuperscript{\ensuremath{x}}\fi}
\newunicodechar{ᵇ}{\ifmmode^{b}\else\textsuperscript{\ensuremath{b}}\fi}
\newunicodechar{ʳ}{\ifmmode^{r}\else\textsuperscript{\ensuremath{r}}\fi}
\newunicodechar{ᵘ}{\ifmmode^{u}\else\textsuperscript{\ensuremath{u}}\fi}
\newunicodechar{ʸ}{\ifmmode^{y}\else\textsuperscript{\ensuremath{y}}\fi}
\newunicodechar{ᵃ}{\ifmmode^{a}\else\textsuperscript{\ensuremath{a}}\fi}
\newunicodechar{ᵉ}{\ifmmode^{e}\else\textsuperscript{\ensuremath{e}}\fi}
\newunicodechar{ᵏ}{\ifmmode^{k}\else\textsuperscript{\ensuremath{k}}\fi}
\newunicodechar{ᵖ}{\ifmmode^{p}\else\textsuperscript{\ensuremath{p}}\fi}
\newunicodechar{ᴺ}{\ifmmode^{N}\else\textsuperscript{\ensuremath{N}}\fi}
\newunicodechar{ᶜ}{\ifmmode^{c}\else\textsuperscript{\ensuremath{c}}\fi}
\newunicodechar{ʰ}{\ifmmode^{h}\else\textsuperscript{\ensuremath{h}}\fi}
\newunicodechar{ᵠ}{\ifmmode^{q}\else\textsuperscript{\ensuremath{q}}\fi}
\newunicodechar{ₙ}{\ifmmode_{n}\else\textsubscript{\ensuremath{n}}\fi}
\newunicodechar{ₐ}{\ifmmode_{a}\else\textsubscript{\ensuremath{a}}\fi}
\newunicodechar{ᵢ}{\ifmmode_{i}\else\textsubscript{\ensuremath{i}}\fi}
\newunicodechar{ᵣ}{\ifmmode_{r}\else\textsubscript{\ensuremath{r}}\fi}
\newunicodechar{ₖ}{\ifmmode_{k}\else\textsubscript{\ensuremath{k}}\fi}
\newunicodechar{ᵦ}{\ifmmode_{\beta}\else\textsubscript{\ensuremath{\beta}}\fi}
\newunicodechar{ₓ}{\ifmmode_{x}\else\textsubscript{\ensuremath{x}}\fi}
\newfontfamily\popdisplay{ArchivoBlack-Regular.ttf}[Path=../../../fonts/]
\newfontfamily\poplogo{SpaceGrotesk-Bold.ttf}[Path=../../../fonts/]
\newfontfamily\popsemi{PlusJakartaSans-SemiBold.ttf}[Path=../../../fonts/]
\newfontfamily\popxbold{PlusJakartaSans-ExtraBold.ttf}[Path=../../../fonts/]
\newfontfamily\popxboldls{PlusJakartaSans-ExtraBold.ttf}[Path=../../../fonts/,LetterSpace=3.0]

\setlist[enumerate]{leftmargin=1.7em,itemsep=5pt,topsep=4pt}
\setlist[itemize]{leftmargin=1.7em,itemsep=4pt,topsep=4pt,label={\color{poporange}\textbullet}}
\setlength{\parindent}{0pt}
\setlength{\parskip}{5pt}
\renewcommand{\arraystretch}{1.25}

% pgfplots : style Pop par défaut
\pgfplotsset{
  every axis/.append style={axis line style={popink,line width=1pt},tick style={popink},
    grid style={popline,line width=0.5pt},label style={font=\small},tick label style={font=\footnotesize}},
  popcurve/.style={poporange,line width=1.8pt,no markers,smooth},
}
% Même style disponible dans un \draw TikZ simple (hors pgfplots)
\tikzset{popcurve/.style={poporange,line width=1.8pt}}
\tikzset{popgrid/.style={popline,very thin}}
\colorlet{popcurve}{poporange} % le nom du style est parfois utilisé comme couleur

% ============================================================
% Pill / squiggle
% ============================================================
\newcommand{\poppill}[3]{%
  \tikz[baseline=-0.55ex]\node[draw=popink,line width=1.2pt,rounded corners=2.6pt,%
    fill=#2,inner xsep=6.5pt,inner ysep=3pt]{%
    \popxboldls\fontsize{7}{7}\selectfont\color{#3}\MakeUppercase{#1}};%
}
\newcommand{\popsquiggle}[1]{%
  \tikz[baseline=-0.4ex]{
    \draw[#1,line width=2.6pt,line cap=round]
      plot [smooth] coordinates {(0,0.09) (0.34,0.01) (0.68,0.21) (1.02,0.01) (1.36,0.10)};
  }%
}
% Mot-clé surligné (marqueur orange sous le texte)
\newcommand{\cle}[1]{{\popxbold\color{popdark}#1}}

% ============================================================
% En-tête / pied
% ============================================================
\pagestyle{fancy}
\fancyhf{}
\renewcommand{\headrulewidth}{0pt}
\renewcommand{\footrulewidth}{0pt}
\lhead{\raisebox{-0.15ex}{\poplogo\fontsize{13}{13}\selectfont\color{popink}OnBuch\color{poporange}+}}
\rhead{\poppill{\DOCMATIERE\ \textbullet\ \DOCNIVEAU\ \textbullet\ Leçon \DOCLECON}{white}{popink}}
\lfoot{\popxbold\fontsize{7.6}{7.6}\selectfont\color{popmuted}\MakeUppercase{Cours signé OnBuch+}\ \textbullet\ {\popsemi\DOCTITRE}}
\rfoot{\tikz[baseline=-0.6ex]\node[rounded corners=8.5pt,fill=popink,minimum height=17pt,minimum width=17pt,inner xsep=5pt,inner sep=0]{\popxbold\fontsize{8}{8}\selectfont\color{popcream}\thepage\,/\,\pageref*{LastPage}};}

% ============================================================
% Titres
% ============================================================
\newcommand{\chaptitre}[2]{%
  \begin{tcolorbox}[enhanced,colback=poporange,colframe=popink,boxrule=2.6pt,arc=11pt,
      drop shadow=popink,left=15pt,right=15pt,top=12pt,bottom=11pt,before skip=3pt,after skip=18pt]
    \poppill{#2}{popink}{popcream}\par\vspace{9pt}
    {\raggedright\hyphenpenalty=10000\exhyphenpenalty=10000\popdisplay\fontsize{22}{24}\selectfont\color{popcream}\MakeUppercase{#1}\par}
  \end{tcolorbox}
}
% Section numérotée : pastille orange avec le numéro + titre Archivo
\newcounter{popsec}
\newcommand{\coursec}[1]{%
  \stepcounter{popsec}\setcounter{popsub}{0}%
  \par\vspace{16pt}\noindent
  \begin{minipage}{\linewidth}
  \tikz[baseline=-0.7ex]\node[draw=popink,line width=1.6pt,fill=poporange,rounded corners=4pt,minimum size=22pt,inner sep=0]{\popdisplay\fontsize{12}{12}\selectfont\color{popcream}\thepopsec};%
  \hspace{8pt}{\popdisplay\fontsize{15}{17}\selectfont\color{popink}\MakeUppercase{#1}}\par
  \vspace{4pt}\noindent{\color{poporange}\rule{36pt}{3.6pt}}
  \end{minipage}\par\nopagebreak\vspace{8pt}
}
\newcounter{popsub}
\newcommand{\courssub}[1]{%
  \stepcounter{popsub}%
  \par\vspace{9pt}\noindent{\popxbold\fontsize{11.5}{13}\selectfont\color{popink}{\color{poporange}\thepopsec.\thepopsub}\hspace{6pt}#1}\par\nopagebreak\vspace{4pt}
}

% ============================================================
% Boîtes pédagogiques — même recette : fond blanc, bordure épaisse colorée,
% ombre dure encre, badge plein. Titre optionnel [#1].
% ============================================================
\tcbset{popbase/.style={breakable,enhanced,colback=white,boxrule=2.3pt,arc=9pt,
  shadow={3pt}{-3pt}{0pt}{popink},left=14pt,right=14pt,top=11pt,bottom=11pt,before skip=11pt,after skip=15pt,
  fontupper=\popsemi\fontsize{10.6}{14.5}\selectfont\color{popink},
  fontlower=\popsemi\fontsize{10.6}{14.5}\selectfont\color{popink}}}
% Coche dessinée (le glyphe ✓ n'existe pas dans les polices du document)
\newcommand{\popcheck}[1]{\tikz[baseline=-0.5ex]\draw[#1,line width=1.6pt,line cap=round,line join=round](-0.13,0.0)--(-0.04,-0.09)--(0.13,0.1);}
\providecommand{\coche}{\popcheck{popgreen}}
\providecommand{\resultat}[1]{\par\smallskip\noindent\fcolorbox{popgreen}{popgreenL}{\parbox{\dimexpr\linewidth-2\fboxsep-2\fboxrule\relax}{\textbf{Résultat.} #1}}\par\smallskip}
\newcommand{\popboxhead}[3]{\poppill{#1}{#2}{white}\ifx\relax#3\relax\else\hspace{6pt}{\popxbold\fontsize{10.6}{12}\selectfont\color{popink}#3}\fi\par\vspace{7pt}}

\newtcolorbox{aretenir}[1][]{popbase,colframe=popgreen,before upper={\popboxhead{À retenir}{popgreen}{#1}}}
\newtcolorbox{exemplebox}[1][]{popbase,colframe=poporange,before upper={\popboxhead{Exemple}{poporange}{#1}}}
\newtcolorbox{definition}[1][]{popbase,colframe=popblue,before upper={\popboxhead{Définition}{popblue}{#1}}}
\newtcolorbox{propriete}[1][]{popbase,colframe=poppurple,before upper={\popboxhead{Propriété}{poppurple}{#1}}}
\newtcolorbox{methode}[1][]{popbase,colframe=popgold,before upper={\popboxhead{Méthode}{popgold}{#1}}}
\newtcolorbox{attention}[1][]{popbase,colframe=poppink,before upper={\popboxhead{Attention}{poppink}{#1}}}
\newtcolorbox{experience}[1][]{popbase,colframe=popblue,colback=popblueL,before upper={\popboxhead{Expérience}{popblue}{#1}}}
% « Le savais-tu ? » : carte violette pleine (contexte, culture, Cameroun)
\newtcolorbox{savaistu}[1][]{popbase,colback=poppurple,colframe=popink,
  fontupper=\popsemi\fontsize{10.4}{14.2}\selectfont\color{white},
  before upper={\poppill{Le savais-tu\,?}{white}{poppurple}\ifx\relax#1\relax\else\hspace{6pt}{\popxbold\color{white}#1}\fi\par\vspace{7pt}}}
% Exercice résolu : énoncé en haut, \tcblower puis la solution sur fond crème
\newcounter{popexo}\newcounter{popexr}
\newtcolorbox{exoresolu}[1][]{popbase,colframe=poporange,colbacklower=popcream,
  segmentation style={popink,line width=1.2pt,dashed},
  before upper={\stepcounter{popexr}\popboxhead{Exercice résolu \thepopexr}{poporange}{#1}},
  before lower={\poppill{Solution}{popink}{popcream}\par\vspace{6pt}}}
% Exercice (non résolu en place) — la correction est dans la partie Corrigés
\newtcolorbox{exercice}[1][]{popbase,colframe=popink,
  before upper={\stepcounter{popexo}\popboxhead{Exercice \thepopexo}{popink}{#1}}}
% Corrigé
\newtcolorbox{corrige}[1][]{popbase,colframe=popgreen,colback=popgreenL,
  before upper={\popboxhead{Corrigé}{popgreen}{#1}}}
% Objectifs / prérequis / bilan
\newtcolorbox{objectifs}{popbase,colframe=popink,colback=poporangeL,
  before upper={\popboxhead{Objectifs de la leçon}{popink}{}}}
\newtcolorbox{prerequis}{popbase,colframe=popink,colback=popblueL,
  before upper={\popboxhead{Prérequis}{popblue}{}}}
\newtcolorbox{bilan}{popbase,colframe=popink,colback=white,boxrule=2.8pt,
  before upper={\popboxhead{Fiche bilan}{poporange}{L'essentiel à connaître pour l'examen}}}
% Figure encadrée avec légende
\newtcolorbox{popfigure}[1][]{enhanced,breakable=false,colback=white,colframe=popline,boxrule=1.4pt,arc=8pt,
  left=10pt,right=10pt,top=10pt,bottom=8pt,before skip=10pt,after skip=12pt,halign=center,
  lower separated=false,halign lower=center,
  fontlower=\popsemi\fontsize{9}{11.5}\selectfont\color{popmuted},
  #1}
\newcommand{\legende}[1]{\par\vspace{6pt}{\popsemi\fontsize{9}{11.5}\selectfont\color{popmuted}#1}}

% ============================================================
% Signature OnBuch+
% ============================================================
\newcommand{\poplogoinline}[1]{{\poplogo\fontsize{#1}{#1}\selectfont\color{popink}OnBuch\color{poporange}+}}
\newcommand{\signatureonbuch}{%
  \begin{tcolorbox}[enhanced,colback=popink,colframe=popink,boxrule=2.2pt,arc=10pt,
      drop shadow=poporange,left=14pt,right=14pt,top=10pt,bottom=10pt,before skip=0pt,after skip=0pt]
    \tikz[baseline=-0.7ex]\node[circle,fill=popgreen,draw=popcream,line width=1.3pt,minimum size=22pt,inner sep=0]{\popcheck{white}};%
    \hspace{9pt}\begin{minipage}[c]{0.62\linewidth}
      {\popxbold\fontsize{12}{13}\selectfont\color{popcream}Signé\ \textbullet\ {\poplogo OnBuch\color{poporange}+}}\par\vspace{2pt}
      {\raggedright\popsemi\fontsize{8}{10.5}\selectfont\color{popline}\MakeUppercase{Équipe pédagogique OnBuch+}\\\MakeUppercase{Conforme au programme officiel MINESEC}\par}
    \end{minipage}\hfill
    {\poplogo\fontsize{22}{22}\selectfont\color{popcream}OnBuch\color{poporange}+}
  \end{tcolorbox}%
}

% ============================================================
% Page de garde
% #1 titre · #2 matière · #3 niveau · #4 module · #5 n° leçon · #6 durée ·
% #7 description / objectifs (texte ou itemize)
% ============================================================
\newcommand{\pagegarde}[7]{%
  \thispagestyle{empty}
  \newgeometry{a4paper,margin=1.7cm,top=1.6cm,bottom=1.4cm}%
  \begin{tikzpicture}[remember picture,overlay]
    \fill[poporange!18] ([xshift=-3.2cm,yshift=-3.6cm]current page.north east) circle (4.6cm);
    \fill[poppurple!14] ([xshift=2.4cm,yshift=7.6cm]current page.south west) circle (3.8cm);
  \end{tikzpicture}%
  {\poplogo\fontsize{30}{30}\selectfont\color{popink}OnBuch\color{poporange}+}\hfill
  \raisebox{0.6ex}{\poppill{Cours signé \textbullet\ Premium}{popink}{popcream}}\par
  \vspace{1pt}\popsquiggle{poppurple}\par
  \vspace{8pt}
  \begin{tcolorbox}[enhanced,colback=poporange,colframe=popink,boxrule=2.8pt,arc=13pt,
      drop shadow=popink,left=18pt,right=18pt,top=12pt,bottom=13pt,after skip=10pt]
    \poppill{#2 \textbullet\ #3}{popink}{popcream}\hspace{5pt}\poppill{#4}{white}{popink}\par\vspace{8pt}
    {\popdisplay\fontsize{14}{14}\selectfont\color{popink}LEÇON #5}\par\vspace{4pt}
    {\raggedright\hyphenpenalty=10000\exhyphenpenalty=10000\popdisplay\fontsize{25}{27}\selectfont\color{popcream}\MakeUppercase{#1}\par}
    \vspace{8pt}
    {\popxbold\fontsize{9.5}{9.5}\selectfont\color{popink}\MakeUppercase{Durée conseillée : #6}}
  \end{tcolorbox}
  \begin{tcolorbox}[enhanced,colback=white,colframe=popink,boxrule=2.2pt,arc=10pt,
      drop shadow=popink,left=16pt,right=16pt,top=10pt,bottom=10pt,after skip=8pt]
    \poppill{Dans cette leçon}{poppurple}{white}\par\vspace{5pt}
    \popsemi\fontsize{9.3}{12.6}\selectfont\color{popink}#7
  \end{tcolorbox}
  \noindent\begin{minipage}[t]{0.487\linewidth}
    \begin{tcolorbox}[enhanced,colback=popgreen,colframe=popink,boxrule=2.2pt,arc=10pt,
        drop shadow=popink,left=11pt,right=11pt,top=10pt,bottom=10pt,height=3.1cm,valign=center]
      \tcbox[colback=white,colframe=popink,boxrule=1.3pt,arc=5pt,boxsep=0pt,nobeforeafter,
        left=3pt,right=3pt,top=3pt,bottom=3pt,valign=center]{\qrcode[height=1.9cm]{https://play.google.com/store/apps/details?id=com.luvvix.onbuch&hl=fr}}%
      \hspace{8pt}\begin{minipage}[c]{0.5\linewidth}
        {\popdisplay\fontsize{11}{12.5}\selectfont\color{white}\MakeUppercase{Télécharge OnBuch}}\par\vspace{4pt}
        {\raggedright\popsemi\fontsize{8.4}{11}\selectfont\color{white}Gratuit \textbullet\ Google Play\\Tuteur IA, annales, concours, résultats\par}
      \end{minipage}
    \end{tcolorbox}
  \end{minipage}\hfill
  \begin{minipage}[t]{0.487\linewidth}
    \begin{tcolorbox}[enhanced,colback=poppurple,colframe=popink,boxrule=2.2pt,arc=10pt,
        drop shadow=popink,left=11pt,right=11pt,top=10pt,bottom=10pt,height=3.1cm,valign=center]
      \tikz[baseline=-0.7ex]\node[circle,fill=white,draw=popink,line width=1.3pt,minimum size=1.6cm,inner sep=0]{\popdisplay\fontsize{24}{24}\selectfont\color{poppurple}+};%
      \hspace{8pt}\begin{minipage}[c]{0.6\linewidth}
        {\popdisplay\fontsize{11}{12.5}\selectfont\color{white}\MakeUppercase{Passe à OnBuch+}}\par\vspace{4pt}
        {\raggedright\popsemi\fontsize{8.4}{11}\selectfont\color{white}Tous les cours signés, Tuteur IA illimité, fiches PDF — onglet OnBuch+ de l'app\par}
      \end{minipage}
    \end{tcolorbox}
  \end{minipage}
  \vspace{8pt}
  \popcontenu
  \vfill
  \signatureonbuch
  \restoregeometry
  \newpage
}

% Rangée « Ce cours contient » (page de garde)
\newcommand{\poptile}[3]{% #1 couleur #2 gros texte #3 légende
  \begin{tcolorbox}[enhanced,colback=#1,colframe=popink,boxrule=2pt,arc=9pt,shadow={2.5pt}{-2.5pt}{0pt}{popink},
      left=6pt,right=6pt,top=9pt,bottom=9pt,halign=center,valign=center,height=1.75cm,width=\linewidth,nobeforeafter]
    {\popdisplay\fontsize{12.5}{13}\selectfont\color{white}\MakeUppercase{#2}}\par\vspace{3pt}
    {\centering\popsemi\fontsize{7.8}{9.5}\selectfont\color{white}#3\par}
  \end{tcolorbox}}
\newcommand{\popcontenu}{%
  \par\noindent{\popxboldls\fontsize{7.5}{7.5}\selectfont\color{popmuted}CE COURS SIGNÉ CONTIENT}\par\vspace{7pt}
  \noindent\begin{minipage}[t]{0.235\linewidth}\poptile{popblue}{Cours}{complet, illustré, pas à pas}\end{minipage}\hfill
  \begin{minipage}[t]{0.235\linewidth}\poptile{poporange}{Méthodes}{et exercices résolus}\end{minipage}\hfill
  \begin{minipage}[t]{0.235\linewidth}\poptile{poppink}{Exercices}{type examen + corrigés}\end{minipage}\hfill
  \begin{minipage}[t]{0.235\linewidth}\poptile{popgreen}{Fiche bilan}{l'essentiel à retenir}\end{minipage}\par}

% Sommaire
\newtcolorbox{sommaire}{enhanced,colback=white,colframe=popink,boxrule=2.2pt,arc=10pt,
  drop shadow=popink,left=18pt,right=18pt,top=13pt,bottom=13pt,before skip=6pt,after skip=20pt,
  before upper={\poppill{Sommaire}{poppurple}{white}\par\vspace{9pt}\popsemi\fontsize{10.6}{14.5}\selectfont\color{popink}}}

% Signature de fin de cours — poussée tout en bas de la dernière page
\newcommand{\findecours}{%
  \par\vfill\nopagebreak
  \begin{center}\popsquiggle{poporange}\end{center}\vspace{4pt}
  \signatureonbuch
}
\AtBeginDocument{\def\llbracket{\mathopen{[\mkern-3.5mu[}}\def\rrbracket{\mathclose{]\mkern-3.5mu]}}} % intervalles d'entiers (glyphe ⟦ absent de la police)
% Glyphes absents de Fira Math : redéfinis à partir de glyphes disponibles
\AtBeginDocument{%
  \def\setminus{\mathbin{\backslash}}\let\smallsetminus\setminus
  \def\vdots{\mathord{\vbox{\baselineskip3.2pt\lineskiplimit0pt\kern4pt\hbox{.}\hbox{.}\hbox{.}}}}%
  \def\ddots{\mathinner{\mkern1mu\raise6.5pt\hbox{.}\mkern2mu\raise3.6pt\hbox{.}\mkern2mu\raise0.7pt\hbox{.}\mkern1mu}}%
  \def\triangle{\mathord{\scalebox{0.9}{$\symup{\Delta}$}}}%
  \def\nabla{\mathord{\raisebox{\depth}{\scalebox{1}[-1]{$\symup{\Delta}$}}}}%
  \def\checkmark{\popcheck{popdark}}%
  \def\star{\mathbin{\ast}}\def\bigstar{\mathord{\ast}}%
  \def\diamond{\mathbin{\scalebox{0.6}{$\Diamond$}}}%
  \def\bigcup{\mathop{\mathchoice{\raisebox{-0.3em}{\scalebox{1.7}{$\cup$}}}{\scalebox{1.25}{$\cup$}}{\cup}{\cup}}\displaylimits}%
  \def\bigcap{\mathop{\mathchoice{\raisebox{-0.3em}{\scalebox{1.7}{$\cap$}}}{\scalebox{1.25}{$\cap$}}{\cap}{\cap}}\displaylimits}%
  \def\lesseqqgtr{\mathrel{\lessgtr}}\def\bigoplus{\mathop{\oplus}\displaylimits}%
}
\providecommand{\ssi}{si et seulement si} % abréviation fréquente
\AtBeginDocument{\providecommand{\wideparen}{\overparen}} % arc de cercle

%% FIGFIX-BEGIN v3 — correctifs globaux des figures (docs/onbuch-generation/tools/figfix)
\usepackage{adjustbox}\usepackage{etoolbox}
% toute figure TikZ/pgfplots plus large que la ligne est réduite à la largeur disponible
\BeforeBeginEnvironment{tikzpicture}{\begin{adjustbox}{max width=\linewidth}}
\AfterEndEnvironment{tikzpicture}{\end{adjustbox}}
% légendes pgfplots lisibles par-dessus les courbes
\pgfplotsset{every axis/.append style={legend style={fill=white,fill opacity=0.92,text opacity=1,draw=black!15,inner sep=3pt}}}
% étiquettes d'arêtes (syntaxe edge["..."]) avec fond blanc
\tikzset{every edge quotes/.append style={fill=white,inner sep=1.2pt}}
% bulles de cartes mentales : texte à la ligne dans la bulle, bulle un peu plus grande
\tikzset{every concept/.append style={text width=2.0cm,align=center,minimum size=2.3cm,inner sep=2pt}}
%% FIGFIX-END
\begin{document}

\pagegarde{Rédiger un résumé de texte}{Français}{Tle Tech}{Français – classe de Terminale technique}{N22}{10 séances (fiche de progression harmonisée)}{Cette leçon outille l'élève de Terminale technique pour rédiger un résumé de texte fidèle, cohérent et correctement formulé, à travers l'étude du Vieux nègre et la médaille de Ferdinand Oyono. Elle mobilise la lecture analytique, la reformulation, la grammaire de l'énonciation et la stylistique des tonalités satirique et polémique.
\vspace{4pt}
\begin{multicols}{2}\small
\begin{itemize}
\item À la fin de cette leçon, je sais lire activement un texte narratif et en dégager le thème et la thèse
\item À la fin de cette leçon, je sais identifier et hiérarchiser les idées essentielles d'un texte
\item À la fin de cette leçon, je sais reformuler des idées sans recopier les formulations de l'auteur
\item À la fin de cette leçon, je sais distinguer un énoncé constatif d'un énoncé performatif et les utiliser à bon escient
\item À la fin de cette leçon, je sais repérer les tonalités satirique et polémique d'un texte et en rendre compte dans mon résumé
\item À la fin de cette leçon, je sais rédiger un résumé respectant la contrainte de longueur (1/4 du texte)
\end{itemize}\end{multicols}}

\begin{sommaire}
\begin{multicols}{2}
\begin{enumerate}
\item Lecture et compréhension du Vieux nègre et la médaille (Exposé N°1)
\item La contraction de texte : reformuler les idées essentielles
\item Langue : énoncés constatifs et performatifs
\item Stylistique : les tonalités satirique et polémique
\item Production et amélioration du résumé
\item Contexte de l'œuvre, bilan et intégration
\item Activité d'intégration
\item Méthodes et exercices résolus
\item Exercices
\item Corrigés des exercices
\item Fiche bilan
\end{enumerate}
\end{multicols}
\end{sommaire}

\begin{objectifs}
\begin{itemize}
\item À la fin de cette leçon, je sais lire activement un texte narratif et en dégager le thème et la thèse
\item À la fin de cette leçon, je sais identifier et hiérarchiser les idées essentielles d'un texte
\item À la fin de cette leçon, je sais reformuler des idées sans recopier les formulations de l'auteur
\item À la fin de cette leçon, je sais distinguer un énoncé constatif d'un énoncé performatif et les utiliser à bon escient
\item À la fin de cette leçon, je sais repérer les tonalités satirique et polémique d'un texte et en rendre compte dans mon résumé
\item À la fin de cette leçon, je sais rédiger un résumé respectant la contrainte de longueur (1/4 du texte)
\item À la fin de cette leçon, je sais situer Le Vieux nègre et la médaille dans son contexte colonial et littéraire
\item À la fin de cette leçon, je sais corriger et améliorer mon propre résumé à l'aide d'une grille d'évaluation
\item À la fin de cette leçon, je sais justifier ma démarche de contraction dans une copie d'examen
\end{itemize}
\end{objectifs}

\begin{prerequis}
\begin{itemize}
\item Notions de résumé vues en classe de Première (dégager l'essentiel, respecter la longueur)
\item Types de textes et registres littéraires (narratif, descriptif, argumentatif)
\item Figures de style de base : ironie, antiphrase, hyperbole, comparaison
\item Grammaire de la phrase : types et formes de phrases (4e-3e)
\item Connaissance générale de la colonisation du Cameroun et de la littérature africaine francophone
\end{itemize}
\end{prerequis}

\newpage

\coursec{Lecture et compréhension du Vieux nègre et la médaille (Exposé N°1)}

À l'approche du Baccalauréat technique, un élève de Yaoundé doit lire un texte de 40 lignes en 30 minutes et en restituer l'essentiel en 10 lignes, sans trahir l'auteur. Au ministère, un secrétaire résume chaque jour de longs rapports pour son chef de service pressé. Sur les réseaux sociaux, les jeunes Camerounais condensent des articles entiers en quelques phrases pour leurs abonnés. Comment faire tenir l'essentiel d'un texte dans un quart de sa longueur, sans déformation ni copie ? C'est tout l'art du \cle{résumé}, compétence clé du Bac et de la vie professionnelle.

Mais avant de résumer un texte, il faut d'abord savoir le \emph{lire}. Et pour s'entraîner à lire activement, rien de tel qu'une œuvre complète du programme : \emph{Le Vieux nègre et la médaille} de Ferdinand Oyono. Cette première séance pose les fondations : connaître l'auteur, comprendre l'intrigue et les personnages, et maîtriser la technique de lecture qui servira ensuite à contracter n'importe quel texte.

\textbf{Problématique :} comment lire un texte de manière active et méthodique pour en dégager le thème, la thèse et les idées essentielles, conditions indispensables d'un bon résumé ?

\courssub{Ferdinand Oyono : vie, œuvres et place dans la littérature africaine}

\begin{definition}[Qui est Ferdinand Oyono ?]
Ferdinand Léopold Oyono (1929-2010) est un écrivain camerounais, né à Ngoulemakong, dans la région du Sud. Romancier de langue française, il est l'une des grandes voix de la \cle{littérature africaine} francophone d'expression anticoloniale. Il est aussi un homme d'État : diplomate (ambassadeur du Cameroun dans plusieurs pays) et ministre.
\end{definition}

Son œuvre romanesque est brève mais capitale :

\begin{itemize}
\item \emph{Une vie de boy} (1956) : journal d'un domestique africain au service des colons ;
\item \emph{Le Vieux nègre et la médaille} (1956) : notre œuvre au programme ;
\item \emph{Chemin d'Europe} (1960) : un jeune Africain rêve de partir en Europe.
\end{itemize}

Ces trois romans dénoncent la même réalité : la \cle{colonisation} et ses mensonges. Oyono écrit avant les indépendances, à une époque où prendre la plume contre l'administration coloniale est un acte de courage. Avec Mongo Beti (autre grand écrivain camerounais), il incarne la génération qui a utilisé le roman comme arme de combat. Sa place est donc double : \textbf{littéraire} (un classique étudié dans toute l'Afrique francophone) et \textbf{historique} (un témoin de la lutte pour la dignité africaine).

\begin{savaistu}[Un écrivain engagé... et un ministre]
Ferdinand Oyono, écrivain camerounais, a aussi été diplomate et ministre : ambassadeur, puis ministre des Relations extérieures et ministre de la Culture. \emph{Le Vieux nègre et la médaille} paraît en \textbf{1956}, soit quatre ans \emph{avant} l'indépendance du Cameroun (1960). Dénoncer la colonisation par écrit, à cette époque, exposait l'auteur à la censure et aux poursuites.
\end{savaistu}

\courssub{Lecture analytique du roman : les personnages}

L'intrigue repose sur un petit nombre de personnages, qu'il faut savoir situer avec précision.

\textbf{Meka}, le « vieux nègre » du titre, est le héros. Vieillard du village de Doum, il a tout donné à la France coloniale : ses deux fils sont morts à la guerre « pour la France », et ses terres ont été cédées à la mission catholique. En récompense, l'administration décide de lui remettre une \cle{médaille} de l'amitié franco-africaine. Meka est naïf, fier, crédule : il croit sincèrement que les Blancs sont ses amis.

\textbf{Kelara}, sa femme, est plus lucide. Elle se méfie des Blancs et ne comprend pas l'enthousiasme de son mari. Elle représente la sagesse populaire et le bon sens du village.

\textbf{Le commandant} du cercle représente l'autorité coloniale. C'est lui qui organise la cérémonie de remise de la médaille. Il incarne l'administration : paternaliste, orgueilleuse, convaincue de sa supériorité.

\textbf{Le père Vandermayer} est le missionnaire. Il profite des terres que Meka a données à la mission. Il représente la religion chrétienne utilisée comme instrument de la colonisation : il prêche l'amour tout en exploitant les villageois.

\textbf{Les Gonthier}, couple de colons commerçants, incarnent le racisme ordinaire des petits Blancs installés en Afrique. Madame Gonthier, en particulier, méprise ouvertement les Africains.

\begin{popfigure}
\begin{center}
\begin{tikzpicture}[
  every node/.style={font=\small},
  hero/.style={circle, draw=popink, fill=poporangeL, minimum size=1.9cm, align=center, font=\small\bfseries},
  perso/.style={rectangle, rounded corners, draw=popblue, fill=popblueL, align=center, minimum width=2.6cm, minimum height=0.9cm, font=\small},
  rel/.style={-{Stealth[length=2.5mm]}, thick, popdark}
]
\node[hero, align=center] (meka) at (0,0){Meka\\ (le vieux nègre)};
\node[perso, align=center] (kelara) at (-4.6,1.8){Kelara\\ (son épouse)};
\node[perso, align=center] (commandant) at (4.6,1.8){Le commandant\\ (administration)};
\node[perso, align=center] (vander) at (-4.6,-1.8){Père Vandermayer\\ (mission)};
\node[perso, align=center] (gonthier) at (4.6,-1.8){Les Gonthier\\ (colons)};
\draw[rel] (kelara) -- node[above, sloped, font=\footnotesize]{méfiance} (meka);
\draw[rel] (commandant) -- node[above, sloped, font=\footnotesize]{décore} (meka);
\draw[rel] (vander) -- node[above, sloped, font=\footnotesize]{exploite} (meka);
\draw[rel] (gonthier) -- node[above, sloped, font=\footnotesize]{méprisent} (meka);
\draw[rel, poppurple, dashed] (commandant) -- node[right, font=\footnotesize]{complice} (vander);
\end{tikzpicture}
\end{center}
\legende{Carte des relations entre les personnages du roman. Meka est au centre : toutes les forces coloniales (administration, mission, colons) convergent vers lui, tandis que Kelara, son épouse, reste la voix de la lucidité.}
\end{popfigure}

\courssub{L'intrigue : de la médaille à la désillusion (Exposé N°1 et N°2)}

Le roman se construit en cinq grandes étapes, qu'il faut pouvoir restituer dans l'ordre. L'Exposé N°1 couvre généralement les trois premières étapes (montée de l'illusion), l'Exposé N°2 les deux dernières (chute et lucidité).

\begin{enumerate}
\item \textbf{L'annonce :} Meka apprend qu'il va recevoir la médaille de l'amitié franco-africaine. Le village de Doum est en émoi ; le vieillard est ivre de fierté.
\item \textbf{La préparation :} on lui confectionne un costume, on le coiffe, on le « civilise » pour la cérémonie. Le village fête l'événement.
\item \textbf{La cérémonie :} le commandant épingle la médaille sur la poitrine de Meka, devant les Blancs et les villageois rassemblés. Discours pompeux sur l'amitié franco-africaine.
\item \textbf{L'arrestation :} la nuit suivante, Meka, perdu sous la pluie après la fête, est arrêté par les gardes de cercle qui ne le reconnaissent pas. Il est \emph{battu et emprisonné}, sa médaille ne le protège de rien.
\item \textbf{La désillusion :} Meka comprend que la médaille n'était qu'un morceau de métal. L'amitié des Blancs était un mensonge. Il retrouve son village et sa dignité, mais ses illusions sont mortes.
\end{enumerate}

\begin{popfigure}
\begin{center}
\begin{tikzpicture}[
  etape/.style={circle, draw=popink, fill=popgoldL, minimum size=0.7cm, font=\small\bfseries},
  txt/.style={below, align=center, font=\footnotesize, text width=2.5cm}
]
\draw[-{Stealth[length=3mm]}, very thick, popline] (0,0) -- (13.6,0);
\node[etape] (e1) at (0.6,0) {1};
\node[etape] (e2) at (3.8,0) {2};
\node[etape] (e3) at (7.0,0) {3};
\node[etape] (e4) at (10.2,0) {4};
\node[etape] (e5) at (13.4,0) {5};
\node[txt] at (e1.south) {Annonce de la médaille};
\node[txt] at (e2.south) {Préparation de la cérémonie};
\node[txt] at (e3.south) {Remise de la médaille};
\node[txt] at (e4.south) {Arrestation et nuit en prison};
\node[txt] at (e5.south) {Désillusion et retour au village};
\draw[popgreen, thick, -{Stealth}] (0.6,0.55) .. controls (4,1.5) .. (7.0,0.55);
\node[font=\footnotesize, popgreen] at (3.8,1.55) {illusion croissante (Exposé 1)};
\draw[popink, thick, -{Stealth}] (7.0,0.55) .. controls (10,1.5) .. (13.4,0.55);
\node[font=\footnotesize, popink] at (10.2,1.55) {chute et lucidité (Exposé 2)};
\end{tikzpicture}
\end{center}
\legende{Frise chronologique de l'intrigue en cinq étapes. La courbe verte montre la montée de l'illusion (de l'annonce à la cérémonie), la courbe rouge la chute brutale (de l'arrestation à la désillusion finale).}
\end{popfigure}

\begin{exemplebox}[La médaille, symbole de l'hypocrisie coloniale]
La remise de la médaille à Meka illustre l'hypocrisie de la « mission civilisatrice » coloniale. Les colonisateurs prétendent apporter la civilisation et l'amitié, mais leurs actes contredisent leurs discours : le même État qui décore Meka le jour le fait battre la nuit. La médaille n'honore pas l'homme, elle récompense sa \emph{soumission} : Meka a donné ses fils et ses terres sans révolte.
\end{exemplebox}

\courssub{Les thèmes majeurs du roman}

Quatre thèmes traversent l'œuvre et doivent être repérés dès la première lecture :

\begin{itemize}
\item \textbf{L'aliénation coloniale :} Meka a intériorisé les valeurs du colonisateur au point d'oublier les siennes. Il est fier d'être « l'ami des Blancs » et méprise inconsciemment sa propre culture.
\item \textbf{L'illusion de l'assimilation :} la colonisation promettait aux Africains de devenir « presque des Français ». La médaille est le symbole de cette promesse mensongère : jamais Meka ne sera traité en égal.
\item \textbf{La révolte :} après l'humiliation de la prison, Meka ouvre les yeux. Sa révolte est intérieure mais réelle : il rejette le monde des Blancs et retrouve les siens.
\item \textbf{La solidarité villageoise :} le village de Doum partage la joie de Meka, puis sa honte. C'est la communauté qui l'accueille à son retour, quand les « amis » blancs l'ont abandonné.
\end{itemize}

\courssub{La lecture active : une technique professionnelle}

Pourquoi un élève qui lit « normalement » oublie-t-il la moitié du texte ? Parce que la lecture passive laisse filer l'information. L'\emph{expérience} est facile à faire en classe : lisez un paragraphe sans stylo, puis essayez d'en restituer les idées ; recommencez en soulignant et en annotant. La seconde restitution est toujours plus complète et plus fidèle. La lecture active transforme le lecteur en \emph{travailleur} du texte.

\begin{definition}[La lecture active]
La \cle{lecture active} est une lecture méthodique avec prise de notes, visant à dégager la structure et les idées essentielles d'un texte. Elle repose sur trois gestes : \textbf{souligner} les mots-clés, \textbf{annoter} les marges, \textbf{repérer} les articulations du texte.
\end{definition}

Détaillons ces trois gestes :

\begin{enumerate}
\item \textbf{Souligner les mots-clés :} un mot-clé est un mot porteur de sens (nom, verbe d'action, notion abstraite), non un mot-outil. Dans « Meka reçut la médaille de l'amitié franco-africaine », on souligne \emph{Meka}, \emph{médaille}, \emph{amitié franco-africaine}.
\item \textbf{Annoter les marges :} en face de chaque paragraphe, écrire en trois ou quatre mots l'idée principale (« fierté de Meka », « discours du commandant », « arrestation »). Ces annotations formeront le squelette du futur résumé.
\item \textbf{Repérer les articulations :} les connecteurs logiques (\emph{d'abord, mais, cependant, en effet, ainsi, finalement}) signalent la progression de la pensée. Ils indiquent où le texte change de direction.
\end{enumerate}

\begin{methode}[Dégager le thème et la thèse en trois questions]
Pour identifier ce dont parle un texte et ce que l'auteur en dit, pose successivement trois questions :
\begin{enumerate}
\item \textbf{De quoi parle le texte ?} $\rightarrow$ c'est le \cle{thème} (le sujet général). Exemple : la colonisation et ses mensonges.
\item \textbf{Que dit l'auteur sur ce thème ?} $\rightarrow$ c'est la \cle{thèse} (le point de vue défendu). Exemple : l'amitié franco-africaine est une illusion qui cache l'exploitation.
\item \textbf{Comment le dit-il ?} $\rightarrow$ ce sont les \emph{arguments} et les \emph{procédés} (récit, ironie, exemples). Exemple : par le contraste entre la cérémonie glorieuse et la nuit de prison.
\end{enumerate}
\end{methode}

\begin{attention}[Ne pas confondre résumé de l'intrigue et analyse du texte]
Erreur fréquente au Bac : raconter l'histoire (« Meka reçoit une médaille, puis il est arrêté, puis il rentre au village ») en croyant analyser le texte. Raconter ne suffit pas. L'analyse dégage le \emph{sens} : pourquoi l'auteur raconte-t-il cela ? Que dénonce-t-il ? Un bon résumé restitue les \textbf{idées} et la \textbf{thèse}, pas seulement la chronologie des événements.
\end{attention}

\begin{exoresolu}[Appliquer la méthode des trois questions à un passage]
Énoncé : « Le commandant épingle la médaille sur la poitrine de Meka et parle d'amitié franco-africaine. La nuit suivante, les gardes de ce même commandant frappent Meka et le jettent en prison. » Dégagez le thème et la thèse de ce passage.
\tcblower
Solution détaillée :
\begin{enumerate}
\item \emph{De quoi parle le passage ?} Thème : le rapport entre les colonisateurs et les Africains, symbolisé par la médaille.
\item \emph{Que dit l'auteur ?} Thèse : l'amitié proclamée par les colons est un mensonge ; leurs actes (violence, mépris) démentent leurs paroles.
\item \emph{Comment le dit-il ?} Par le \textbf{contraste} brutal entre la cérémonie (jour, honneur, discours) et l'arrestation (nuit, coups, prison) : l'ironie de la situation fait éclater l'hypocrisie.
\end{enumerate}
Réponse rédigée : « Ce passage traite de l'hypocrisie coloniale. Oyono montre que l'amitié franco-africaine n'est qu'un discours mensonger, démenti par la violence réelle de l'administration. Il le démontre par le contraste ironique entre la remise solennelle de la médaille et l'arrestation brutale du décoré. »
\end{exoresolu}

\courssub{Pont vers l'étude de la langue : énoncés constatifs et performatifs}

\textbf{Repère programme :} cette séquence prépare l'étude des \cle{énoncés constatifs} (qui décrivent un état de fait, vérifiables : « Meka a reçu une médaille ») et des \cle{énoncés performatifs} (qui réalisent une action par le fait de l'énoncer : « Je vous remets cette médaille », « Je vous arrête »). Le discours du commandant lors de la cérémonie est performatif : il \emph{crée} l'ami franco-africain par la parole. L'acte d'arrestation par les gardes est aussi un performatif (exercice de la force publique). La contradiction entre le performatif officiel (décorer) et le performatif réel (arrêter) structure l'ironie du roman.

\begin{aretenir}[L'essentiel de la séance]
\begin{itemize}
\item Ferdinand Oyono (1929-2010), écrivain et homme d'État camerounais, publie \emph{Le Vieux nègre et la médaille} en 1956, avant les indépendances.
\item Le roman raconte la médaille offerte à Meka, son arrestation la nuit même et sa désillusion finale : cinq étapes à connaître (Exposé 1 : 1 à 3 ; Exposé 2 : 4 et 5).
\item Thèmes majeurs : aliénation coloniale, illusion de l'assimilation, révolte, solidarité villageoise.
\item La lecture active = souligner les mots-clés + annoter les marges + repérer les articulations.
\item Thème (de quoi ?), thèse (que dit l'auteur ?), procédés (comment ?) : les trois questions fondatrices du résumé.
\item Distinction cruciale : résumé $\neq$ récit de l'intrigue. Le résumé restitue les idées et la thèse.
\end{itemize}
\end{aretenir}

\begin{exercice}[Pour s'entraîner]
1. Répondez aux trois questions de la méthode pour le passage où Meka donne ses terres à la mission.\par
2. Relevez dans un extrait fourni par le professeur cinq mots-clés et deux connecteurs logiques, puis rédigez une annotation de marge pour chaque paragraphe.\par
3. Rédigez en quatre lignes la différence entre le thème et la thèse du roman, en citant un exemple précis.\par
4. Identifiez dans le discours du commandant (chapitre de la cérémonie) un énoncé performatif et un énoncé constatif. Justifiez.
\end{exercice}

\coursec{La contraction de texte : reformuler les idées essentielles}

Après avoir lu et compris \emph{Le Vieux nègre et la médaille} (Exposé N°1), nous abordons maintenant le cœur technique de l'épreuve : la contraction de texte. Au Baccalauréat technique, le résumé n'est pas un simple exercice de raccourcissement ; c'est une \textbf{reconstruction intellectuelle}. Vous devez prouver que vous avez saisi la substantifique moelle du texte pour la restituer avec vos propres mots, dans une langue correcte et une structure autonome.

\courssub{Définition, enjeux et critères de réussite}

\begin{definition}[La contraction de texte (résumé)]
Opération de réécriture qui consiste à réduire un texte source à environ \textbf{un quart de sa longueur} tout en conservant \textbf{uniquement les idées essentielles}, exprimées \textbf{avec ses propres mots} (reformulation), dans le \textbf{respect de la pensée de l'auteur} (fidélité) et selon une \textbf{organisation logique} (cohérence).
\end{definition}

\begin{propriete}[Les quatre piliers d'un bon résumé au Bac Tech]
Un résumé réussi repose sur quatre critères indissociables, notés par le correcteur :
\begin{enumerate}
    \item \textbf{Fidélité :} Respect strict du sens, de la thèse et de la progression logique de l'auteur. Aucune interprétation personnelle, aucun contresens.
    \item \textbf{Concision (Règle du quart) :} Longueur cible = 25\% du texte original (marge de tolérance habituelle : $\pm 10\%$).
    \item \textbf{Reformulation (Appropriation) :} Zéro copier-coller de phrases entières. Changement de syntaxe, de vocabulaire (synonymie), de voix (actif/passif), nominalisation.
    \item \textbf{Cohérence \& Autonomie :} Le résumé se lit comme un texte à part entière. Connecteurs logiques, reprises anaphoriques (pronoms, synonymes), progression thématique claire.
\end{enumerate}
\end{propriete}

\begin{aretenir}[La règle du quart : mode de calcul]
\begin{itemize}
    \item Comptez le nombre de \textbf{mots} du texte original (ou de lignes $\times$ mots par ligne).
    \item Calculez le quart : $N_{\text{objectif}} = \frac{N_{\text{original}}}{4}$.
    \item Intervalle acceptable : $[ \num{0,9} \times N_{\text{objectif}} \ ; \ \num{1,1} \times N_{\text{objectif}} ]$.
\end{itemize}
\textbf{Exemple chiffré :} Un texte de 400 mots impose un résumé d'environ \textbf{100 mots} (entre 90 et 110 mots).
\end{aretenir}

\begin{popfigure}
\begin{tikzpicture}[scale=0.9, every node/.style={font=\small}]
    % Entonnoir
    \filldraw[popblue!20, draw=popblue, thick] (0,0) -- (6,0) -- (5,-4) -- (1,-4) -- cycle;
    \filldraw[poporange!30, draw=poporange, thick] (1,-4) -- (5,-4) -- (4.5,-5.5) -- (1.5,-5.5) -- cycle;
    \filldraw[popgreen!30, draw=popgreen, thick] (1.5,-5.5) -- (4.5,-5.5) -- (4,-6.5) -- (2,-6.5) -- cycle;
    
    % Texte 100%
    \node[align=center, text width=5cm] at (3,1) {\textbf{Texte Original (100\%)}\\\textit{Ex: 400 mots}\\\textit{Toutes les phrases, exemples, détails, connecteurs}};
    \draw[->, thick, popblue] (3,0.2) -- (3,-0.5);
    
    % Étape 1
    \node[align=center, text width=3.5cm] at (3,-2.5) {\textbf{Analyse \& Sélection}\\\textit{Repérage des idées forces}\\\textit{Suppression exemples, répétitions, digressions}};
    \draw[->, thick, poporange] (3,-3.7) -- (3,-4.3);
    
    % Étape 2
    \node[align=center, text width=2.5cm] at (3,-5) {\textbf{Reformulation}\\\textit{Synonymie, nominalisation}\\\textit{Changement de structure}};
    \draw[->, thick, popgreen] (3,-5.7) -- (3,-6.3);
    
    % Résultat 25%
    \node[align=center, text width=2cm, fill=popgreen!10, draw=popgreen, thick, rounded corners] at (3,-7.2) {\textbf{Résumé (25\%)}\\\textit{Ex: 100 mots}\\\textit{Fidèle, Bref, Clair, Impersonnel}};
    
    % Annotations latérales
    \node[left, align=left, text width=3cm, text=popmuted] at (-2.5, -1) {\textit{Étape 1 : Comprendre \& Découper}\\(Unités de sens)};
    \node[left, align=left, text width=3cm, text=popmuted] at (-2.5, -4) {\textit{Étape 2 : Filtrer}\\(Hiérarchiser)};
    \node[left, align=left, text width=3cm, text=popmuted] at (-2.5, -6) {\textit{Étape 3 : Réécrire}\\(Appropriation)};
\end{tikzpicture}
\legende{Le processus de contraction : de la matière brute (100\%) à l'essence restructurée (25\%).}
\end{popfigure}

\courssub{Étape 1 : Repérer et hiérarchiser les idées essentielles}

La première erreur est de croire que toutes les phrases se valent. Un texte argumentatif ou narratif suit une architecture : une thèse, des arguments, des illustrations. Votre mission : distinguer le \textbf{squelette} (idées principales) de la \textbf{chair} (exemples, précisions, répétitions, digressions).

\begin{methode}[Repérer l'idée essentielle par paragraphe (ou unité de sens)]
\begin{enumerate}
    \item \textbf{Découper} le texte en unités de sens (paragraphe = 1 idée force généralement ; parfois 2 paragraphes = 1 idée si développement long).
    \item \textbf{Poser la question magique :} « \emph{De quoi l'auteur veut-il me convaincre / m'informer ICI et MAINTENANT ?} »
    \item \textbf{Repérer le noyau prédicatif :} Identifiez le \textbf{verbe d'action principal} (souvent un verbe d'opinion : \emph{affirmer, démontrer, critiquer, montrer, expliquer}) et son \textbf{sujet} (l'auteur, « il », « l'écrivain »).
    \item \textbf{Éliminer les expansions :} Supprimez les compléments de circonstance (temps, lieu, manière), les relatives explicatives, les parenthèses, les énumérations d'exemples.
    \item \textbf{Formuler en une phrase nominale ou verbale courte} l'idée retenue. Notez-la en marge (brouillon).
\end{enumerate}
\end{methode}

\begin{exemplebox}[Application sur un extrait type « Vieux nègre »]
\textbf{Texte source (simulé, $\approx$ 60 mots) :}
\begin{quote}
\textit{« Le vieil homme, courbé sous le poids des ans et des médailles qu'il arbore fièrement sur sa poitrine usée par le soleil, avance lentement vers la tribune officielle. La foule, massée derrière les barrières, observe ce spectacle avec un mélange de respect pour l'ancien combattant et de curiosité pour les décorations qui brillent. »}
\end{quote}
\textbf{Analyse :}
\begin{itemize}
    \item Sujet/Verbe noyau : \textbf{Le vieil homme avance} vers la tribune.
    \item Expansions à supprimer : « courbé sous... », « qu'il arbore... », « usée par... », « massée derrière... », « avec un mélange... ».
    \item \textbf{Idée essentielle (marge) :} \textit{Le vieux médaillé rejoint la tribune sous le regard de la foule.}
\end{itemize}
\end{exemplebox}

\begin{attention}[Piège : Confondre « idée principale » et « détail illustratif »]
Dans \emph{Le Vieux nègre et la médaille}, l'auteur décrit souvent des détails sensoriels (la chaleur, la poussière, l'éclat des médailles) pour servir une thèse (l'ironie de la gloire coloniale, la solitude du vétéran).
\begin{itemize}
    \item \textbf{Interdit au résumé :} « Le soleil brille sur les médailles » (détail descriptif).
    \item \textbf{Requis au résumé :} « L'éclat des médailles contraste avec la misère du vétéran » (idée interprétée/essentielle).
\end{itemize}
Ne résumez pas l'image, résumez \textbf{ce que l'image prouve}.
\end{attention}

\courssub{Étape 2 : Techniques de reformulation (L'appropriation langagière)}

C'est ici que se joue la note de « Langue » et de « Reformulation ». Recopier une phrase en changeant deux mots (\emph{patchwork}) est sanctionné. Il faut \textbf{reconstruire la phrase}.

\begin{center}
\begin{tabularx}{\linewidth}{|l|X|X|}\hline
\rowcolor{popblueL}
\textbf{Phrase Originale} & \textbf{Reformulation Acceptable (Vos mots)} & \textbf{Copie Interdite (Patchwork)} \\ \hline
Le vieux nègre, humilié par l'indifférence des officiels, refuse de tendre la main. & L'ancien combattant, méprisé par les autorités, rejette l'aumône. & Le vieux nègre, humilié par l'indifférence des officiels, refuse de tendre la main. \\ \hline
C'est une médaille qui ne vaut rien, car elle a été donnée par des colons ingrats. & La décoration, offerte par des colons ingrats, est dénuée de valeur. & C'est une médaille qui ne vaut rien, car elle a été donnée par des colons ingrats. \\ \hline
Il y a une foule immense qui regarde sans comprendre le drame intérieur du héros. & Une foule nombreuse observe, incompréhensive, la tragédie intime du vétéran. & Il y a une foule immense qui regarde sans comprendre le drame intérieur du héros. \\ \hline
\end{tabularx}
\end{center}

\begin{propriete}[Boîte à outils de la reformulation (4 leviers)]
Pour transformer la structure sans trahir le sens, combinez ces techniques :
\begin{enumerate}
    \item \textbf{Synonymie \& Hyperonymie :} Remplacer un mot précis par un terme plus général ou un synonyme. \textit{Ex: « médailles, croix, décorations » $\rightarrow$ « distinctions honorifiques » / « insignes ».}
    \item \textbf{Nominalisation (Verbe $\rightarrow$ Nom) :} Transforme une action en concept. \textit{Ex: « Il refuse de tendre la main » $\rightarrow$ « Son refus de la mendicité » / « Le rejet de l'aumône ».} \textbf{Gain de place énorme.}
    \item \textbf{Changement de voix / diathèse :} Actif $\leftrightarrow$ Passif / Impersonnel. \textit{Ex: « Les officiels l'ignorent » $\rightarrow$ « Il est ignoré des officiels » / « L'indifférence des officiels le frappe ».}
    \item \textbf{Fusion / Subordination :} Relier deux phrases simples en une complexe. \textit{Ex: « Il est vieux. Il est fier. » $\rightarrow$ « Malgré son âge, il conserve sa fierté. »}
\end{enumerate}
\end{propriete}

\begin{methode}[Réduire un paragraphe à une phrase dense (Entraînement)]
Prenez un paragraphe de 5-6 phrases. Appliquez l'algorithme :
\begin{enumerate}
    \item \textbf{Surlignez} le(s) verbe(s) d'action principale(s) (pas les verbes d'état « être, avoir, sembler » sauf s'ils portent le jugement).
    \item \textbf{Entourez} les sujets de ces verbes.
    \item \textbf{Barrez} au stylo rouge : tous les adjectifs qualificatifs non évaluatifs, tous les compléments de lieu/temps, toutes les relatives, toutes les énumérations d'exemples (garder la classe : « fruits » au lieu de « pommes, poires, bananes »).
    \item \textbf{Relisez} ce qui reste. C'est votre « squelette ».
    \item \textbf{Réécrivez} ce squelette en français fluide en utilisant la nominalisation et la subordination.
\end{enumerate}
\end{methode}

\begin{exoresolu}[Entraînement guidé : De l'extrait au résumé de paragraphe]
\textbf{Texte source (Extrait fictif style \emph{Vieux nègre}, 85 mots) :}
\begin{quote}
\textit{« Le gouverneur, vêtu de son uniforme blanc impeccable, monta sur l'estrade au son de l'hymne national joué par une fanfare militaire un peu essoufflée. Il prononça un long discours fleuri, plein de mots creux sur la gloire de la France et le devoir accompli. Il ne regarda jamais le vieux soldat dans les yeux. Quand il eut fini, il tendit une petite boîte en velours rouge. Le vieux nègre la prit, l'ouvrit, vit la médaille, et la referma sans un mot. »}
\end{quote}
\tcblower
\textbf{Travail sur brouillon (Application de la méthode) :}
\begin{enumerate}
    \item \textbf{Découpage :} 1 unité de sens (la remise de médaille).
    \item \textbf{Verbes noyaux :} \textbf{monta, prononça, ne regarda pas, tendit, prit,'ouvrit, vit, referma}.
    \item \textbf{Suppression expansions :} « vêtu de... », « au son de... », « plein de mots creux... », « quand il eut fini », « sans un mot ».
    \item \textbf{Squelette :} Gouverneur monte estrade $\rightarrow$ prononce discours creux $\rightarrow$ ignore vieux soldat $\rightarrow$ tend boîte $\rightarrow$ Vieux prend/ouvre/voit/referme.
    \item \textbf{Hiérarchisation :} L'essentiel = Contraste entre solennité vide du gouverneur et silence digne du vieux.
\end{enumerate}

\textbf{Reformulation candidate (28 mots) :}
\begin{quote}
\textit{Lors d'une cérémonie officielle creuse, le gouverneur remet distraitement sa médaille à l'ancien combattant qui l'accepte en silence, digne face à l'indifférence.}
\end{quote}

\textbf{Vérification :} Fidèle ? Oui. Bref ? 28/85 $\approx$ \num{33}\% (acceptable pour un paragraphe dense, le global visera 25\%). Reformulé ? Oui (nominalisation « cérémonie creuse », « indifférence », fusion phrases). Impersonnel ? Oui (pas de « je », pas de « nous »).
\end{exoresolu}

\courssub{Étape 3 : Assurer la cohérence du résumé global}

Avoir 5-6 phrases résumées en marge ne fait pas un résumé. Il faut les \textbf{souder} pour former un texte unique qui coule de source.

\begin{definition}[Cohérence textuelle : les trois piliers]
\begin{enumerate}
    \item \textbf{Connecteurs logiques (Relieurs) :} Ils explicitent la relation entre les idées (cause, conséquence, opposition, concession, énumération). \textit{Ex: \textbf{Ainsi, Par conséquent, Cependant, En outre, Enfin}.}
    \item \textbf{Reprises anaphoriques (Chaînes de référence) :} Éviter de répéter « Le vieux nègre », « Le gouverneur ». Utilisez pronoms (\textit{il, ce dernier, celui-ci}), synonymes (\textit{le vétéran, l'ancien combattant, le récipiendaire}), ou nominalisations (\textit{ce silence, cette remise}).
    \item \textbf{Progression thématique :} L'ordre des idées suit la logique de l'auteur (chronologique, argumentative, analytique). Ne pas mélanger la fin et le début.
\end{enumerate}
\end{definition}

\begin{attention}[Erreur fatale : Le « Résumé-Patchwork » (ou « Catalogue »)]
\begin{quote}
\textit{« Le gouverneur arrive. Il fait un discours. Il ne regarde pas le vieux. Il donne la médaille. Le vieux la prend. Il la regarde. Il ne dit rien. »}
\end{quote}
$\rightarrow$ \textbf{Zéro point en expression/cohérence.} C'est une liste d'actions, pas un résumé.
\textbf{Correction :} Utilisez la subordination et la nominalisation pour lier : \textit{« Après un discours vide où il ignore le vétéran, le gouverneur lui remet la distinction que ce dernier accepte silencieusement. »} (1 phrase au lieu de 7).
\end{attention}

\courssub{Les interdits absolus du résumé (Check-list de relecture)}

Avant de rendre votre copie, passez au crible ces \textbf{5 interdits}. Un seul peut faire chuter la note de manière drastique.

\begin{center}
\begin{tabular}{|c|p{12cm}|}\hline
\rowcolor{poporangeL}
\textbf{Interdit} & \textbf{Conséquence / Correction} \\ \hline
1. \textbf{Copier-coller} phrases ou segments $>$ 5 mots consécutifs. & \textbf{Note 0 en reformulation.} $\rightarrow$ Réécrire systématiquement. \\ \hline
2. \textbf{Ajouter} une idée, un exemple, un jugement personnel (« \textit{à mon avis, c'est triste} »). & \textbf{Hors sujet / Infidélité.} $\rightarrow$ Se mettre dans la peau de l'auteur. \\ \hline
3. \textbf{Commenter} le style de l'auteur (« \textit{l'auteur utilise une métaphore} »). & \textbf{Hors consigne.} $\rightarrow$ Résumer le \textbf{contenu}, pas la forme (sauf si sujet explicite). \\ \hline
4. \textbf{Modifier le sens} (contresens, inversion cause/effet). & \textbf{Pénalité lourde fidélité.} $\rightarrow$ Relire le texte source avant finalisation. \\ \hline
5. \textbf{Dépasser la longueur} (hors tolérance $\pm 10\%$). & \textbf{Pénalité concision.} $\rightarrow$ Compter les mots \textbf{avant} la copie propre. \\ \hline
\end{tabular}
\end{center}

\begin{savaistu}[Le résumé au Cameroun : spécificités du Probatoire / Bac Tech]
Au Probatoire et au Baccalauréat technique (séries F, G, H, ST, etc.), le texte à résumer fait souvent \textbf{350 à 500 mots}. Le résumé attendu tourne donc autour de \textbf{90 à 125 mots}.
Le correcteur attend une \textbf{présentation soignée} :
\begin{itemize}
    \item Titre : \textbf{RÉSUMÉ} (centré, majuscules, souligné).
    \item Un seul bloc (pas de paragraphes multiples sauf si texte très long et structure marquée).
    \item Nombre de mots indiqué entre parenthèses à la fin : \textit{(102 mots)}.
    \item Aucune citation, aucun tiret de dialogue, aucune référence de ligne (p. 12, l. 5).
\end{itemize}
\end{savaistu}

\begin{exercice}[Entraînement autonome : Contraction progressive]
\textbf{Consigne :} À partir du paragraphe ci-dessous (extrait type \emph{Vieux nègre}, $\approx$ 70 mots), effectuez le travail complet sur votre cahier de brouillon :
\begin{enumerate}
    \item Identifiez l'unité de sens et le verbe noyau.
    \item Barrez les expansions (exemples, circonstances, répétitions).
    \item Écrivez l'idée essentielle en 1 phrase (marge).
    \item Reformulez cette phrase en utilisant \textbf{au moins 2 techniques} (nominalisation + synonymie OU passif + fusion).
    \item Comptez les mots de votre phrase finale.
\end{enumerate}

\textbf{Texte :}
\textit{« Le vieil Africain, dont le visage ridé raconte toutes les souffrances d'une vie de labeur sous le joug colonial, reste planté devant l'officier français. Ce dernier, jeune et arrogant, lui tend un bout de métal brillant sans même daigner lui serrer la main. Le vieux regarde l'objet, puis l'officier, et crache par terre, geste unique de sa révolte muette. »}
\end{exercice}

\begin{corrige}[Exercice n : Contraction progressive]
\textbf{1. Analyse :} Unité de sens = Remise de médaille / Confrontation muette. Verbes noyaux : \textbf{reste, tend, regarde, crache}.
\textbf{2. Élagage :} Supprimer « dont le visage... colonial », « jeune et arrogant », « sans même daigner... », « geste unique de sa révolte muette » (interprétation à garder mais à reformuler).
\textbf{3. Idée essentielle (marge) :} Le vétéran reçoit la médaille de l'officier arrogant et manifeste son refus en crachant.
\textbf{4. Reformulations possibles (visant $\approx$ 20-25 mots) :}
\begin{itemize}
    \item \textit{Devant l'arrogance de l'officier qui lui tend la médaille sans égards, le vieux soldat exprime sa révolte par un crachat méprisant.} (23 mots) \checkmark \textit{(Synonymie « égards/daigner », fusion cause/effet)}.
    \item \textit{La remise dédaigneuse de la distinction par l'officier provoque le crachat silencieux du vétéran, unique réponse à l'humiliation.} (21 mots) \checkmark \textit{(Nominalisation « remise », « réponse », passif « provoque »)}.
\end{itemize}
\textbf{5. Vérification :} Fidèle ? Oui. Bref ? Oui. Reformulé ? Oui (zéro mot copié de +5 lettres consécutifs). Cohérent ? Oui (liaison cause/effet).
\end{corrige}

\begin{aretenir}[Synthèse de la section : La check-list du contracteur expert]
\begin{enumerate}
    \item \textbf{Lis} $\rightarrow$ \textbf{Découpe} (unités de sens) $\rightarrow$ \textbf{Sélectionne} (1 idée / unité).
    \item \textbf{Élagage} radical : exemples, adjectifs « décoratifs », circonstances, répétitions.
    \item \textbf{Reformule} : Nominalise, subordonne, synonymise, change la voix. \textbf{Jamais} copier-coller.
    \item \textbf{Assemble} : Connecteurs logiques + Reprises anaphoriques = Texte fluide.
    \item \textbf{Compte} : Mots totaux $\approx$ 1/4 original ($\pm 10\%$). Note le chiffre final.
    \item \textbf{Relit} : Sujet ? Verbe ? Accords ? Sens respecté ? \textbf{Titre « RÉSUMÉ » + (X mots)}.
\end{enumerate}
\end{aretenir}

La maîtrise de cette mécanique (\textbf{Analyse $\rightarrow$ Sélection $\rightarrow$ Reformulation $\rightarrow$ Assemblage}) est la clé de la réussite. Dans la prochaine section, nous verrons comment les notions de langue — \textbf{énoncés constatifs vs performatifs} — et de stylistique — \textbf{tonalités satirique et polémique} — éclairent la compréhension fine de \emph{Le Vieux nègre et la médaille} pour affiner encore votre sélection d'idées.

\coursec{Langue : énoncés constatifs et performatifs}

Dans la section précédente, nous avons appris à reformuler les idées essentielles d'un texte : supprimer les redondances, généraliser les exemples, remplacer les longues périphrases par des mots simples. Mais reformuler suppose d'abord de comprendre ce que fait le langage. Or, quand le commandant, dans \emph{Le Vieux nègre et la médaille} de Ferdinand Oyono, déclare la cérémonie ouverte, il ne se contente pas de décrire le monde : il le transforme. Dire, c'est parfois faire. Cette leçon de langue étudie deux grandes familles d'énoncés — les énoncés \cle{constatifs} et les énoncés \cle{performatifs} — et montre pourquoi cette distinction est précieuse pour le résumé : au discours indirect, on ne rapporte pas de la même manière une description et un acte de parole.

\courssub{L'énoncé constatif : décrire le monde}

\begin{definition}[Énoncé constatif]
Un énoncé \cle{constatif} est un énoncé qui décrit une réalité, qui constate un fait. Sa propriété essentielle est d'être \textbf{vrai ou faux} : on peut le confronter à la réalité et vérifier s'il lui correspond.
\end{definition}

Prenons une phrase simple : « Meka porte une médaille sur la poitrine. » Cette phrase décrit une situation. Pour savoir si elle est correcte, il suffit de regarder le vieux Meka : soit la médaille est là (l'énoncé est vrai), soit elle n'y est pas (l'énoncé est faux). C'est le critère fondamental du constatif : il possède une \cle{valeur de vérité}.

\begin{exemplebox}[Énoncés constatifs]
\begin{itemize}
\item « Le vieux Meka a reçu une médaille de l'administration coloniale. » (vrai dans le roman)
\item « La médaille a été perdue pendant la nuit de la tempête. » (vrai dans le roman)
\item « Meka est un jeune homme plein d'avenir. » (faux : Meka est un vieillard)
\item « Il pleut sur le village de Doum. » (vérifiable : vrai ou faux selon le moment)
\end{itemize}
\end{exemplebox}

Remarquons que le constatif ne dépend pas de la personne grammaticale ni du temps du verbe : « je marche », « tu es fatigué », « les Blancs sont arrivés hier » sont tous constatifs, car tous peuvent être démentis par un « c'est faux ! ».

\courssub{L'énoncé performatif : agir par les mots}

\begin{experience}[Observation : quand parler, c'est agir]
Imaginons trois situations de la vie courante.
\begin{enumerate}
\item Le chef de village dit, devant l'assemblée réunie : « Je déclare la réunion ouverte. »
\item Un père dit à sa fille : « Je te promets une bicyclette si tu réussis ton examen. »
\item Un élève dit à son camarade : « Il fait chaud dans cette salle. »
\end{enumerate}
\textbf{Observations.} Dans les situations 1 et 2, une fois les mots prononcés, quelque chose a changé dans le monde : la réunion est désormais ouverte, une promesse existe désormais, avec ses obligations. Aucune action physique supplémentaire n'a été nécessaire. Dans la situation 3, en revanche, rien n'a changé : la température de la salle est exactement la même avant et après la phrase.\par
\textbf{Interprétation.} Il existe donc des énoncés dont le simple fait d'être prononcés \emph{accomplit} une action : ce sont les énoncés performatifs. Les autres se contentent de décrire : ce sont les constatifs.
\end{experience}

\begin{definition}[Énoncé performatif]
Un énoncé \cle{performatif} est un énoncé dont l'émission, dans des circonstances appropriées, constitue l'accomplissement de l'action qu'il énonce. Dire, c'est faire : prononcer l'énoncé, c'est réaliser l'acte (promettre, ordonner, déclarer, baptiser, féliciter, condamner). Cette notion a été théorisée par le philosophe John L. \cle{Austin}.
\end{definition}

\begin{exemplebox}[Un performatif dans la cérémonie de la médaille]
Lors de la cérémonie de remise de la médaille à Meka, le commandant prononce : « Je déclare la cérémonie ouverte. » Au moment précis où ces mots sont dits par cette personne-là, la cérémonie est \emph{réellement} ouverte. Il n'y a rien d'autre à faire : les mots ont suffi. De même, quand l'officier dit à Meka « Je vous félicite au nom de l'administration », la félicitation est accomplie par la phrase elle-même.
\end{exemplebox}

\begin{savaistu}[John L. Austin : faire des choses avec des mots]
La théorie des actes de langage doit beaucoup au philosophe anglais John Langshaw Austin (1911-1960), dont les conférences ont été publiées en 1962 sous le titre \emph{How to Do Things with Words} (en français : \emph{Quand dire, c'est faire}). Austin y montre que le langage ne sert pas seulement à décrire le monde, mais aussi à agir sur lui : baptiser un navire, marier deux personnes, parier, ordonner. Cette découverte a révolutionné la linguistique et la philosophie du langage au XX\textsuperscript{e} siècle.
\end{savaistu}

\courssub{Les marques du performatif}

Comment reconnaître un énoncé performatif ? Il présente des marques formelles assez régulières.

\begin{aretenir}[Signes typiques de l'énoncé performatif]
\begin{itemize}
\item Un \textbf{verbe performatif} : un verbe qui nomme l'acte de parole accompli (\emph{déclarer, promettre, ordonner, féliciter, remercier, baptiser, condamner, pardonner, jurer, interdire, autoriser}).
\item La \textbf{première personne} du singulier ou du pluriel : \emph{je} (ou \emph{nous}).
\item Le \textbf{présent de l'indicatif}, qui coïncide avec le moment de la parole.
\item Éventuellement la formule solennelle « \textbf{par la présente} » : « Je déclare, par la présente, la session ouverte. »
\end{itemize}
Modèle : \emph{je} + \emph{verbe performatif} + \emph{présent} $\rightarrow$ acte accompli.
\end{aretenir}

\begin{attention}[Piège : toute phrase à la première personne n'est pas performative]
C'est l'erreur la plus fréquente. « Je marche », « je mange », « je pense » sont à la première personne du présent, pourtant ce sont des \textbf{constatifs} : ils décrivent une action que je fais avec mes jambes ou mon esprit, pas avec mes mots. On peut répondre « c'est faux, tu ne marches pas, tu es assis ». En revanche, on ne peut pas répondre « c'est faux » à « je te promets de venir » : la promesse existe dès qu'elle est dite. Seuls les verbes qui nomment un \emph{acte de parole} (dire quelque chose, c'est le faire) rendent l'énoncé performatif.
\end{attention}

\courssub{Les conditions de réussite d'un acte performatif}

Un constatif est vrai ou faux. Un performatif, lui, \textbf{réussit ou échoue}. Austin appelle « conditions de félicité » les conditions nécessaires pour que l'acte soit valide.

\begin{propriete}[Conditions de réussite d'un performatif]
Pour qu'un énoncé performatif accomplisse réellement son acte, il faut :
\begin{enumerate}
\item l'\textbf{autorité du locuteur} : celui qui parle doit être habilité à accomplir l'acte (le commandant peut déclarer la cérémonie ouverte ; un passant ne le peut pas) ;
\item des \textbf{circonstances appropriées} : le bon lieu, le bon moment, le bon rituel (on ne baptise pas un enfant pendant un match de football) ;
\item un \textbf{destinataire recevable} et une exécution correcte et complète de la formule.
\end{enumerate}
Si l'une de ces conditions manque, l'acte est \textbf{manqué} : il est nul, sans effet.
\end{propriete}

\begin{exemplebox}[Performatifs réussis et manqués]
\begin{itemize}
\item \textbf{Réussi} : le commandant, devant la foule rassemblée pour la fête, dit « Je déclare la cérémonie ouverte » $\rightarrow$ la cérémonie est ouverte.
\item \textbf{Manqué} : un écolier, seul dans sa cour, crie « Je déclare la cérémonie ouverte » $\rightarrow$ rien n'est ouvert : il n'a ni l'autorité ni les circonstances.
\item \textbf{Manqué} : « Je te lègue ma maison » dit en plaisantant à un inconnu dans un taxi $\rightarrow$ aucun legs n'existe sans testament en bonne forme.
\end{itemize}
\end{exemplebox}

\courssub{Deux outils pour distinguer constatif et performatif}

\begin{methode}[Distinguer un constatif d'un performatif en deux tests]
\begin{enumerate}
\item \textbf{Le test de vérité.} Pose la question : « Cet énoncé peut-il être vrai ou faux ? Peut-on répondre "c'est faux" ? » Si oui, c'est un \textbf{constatif}. (« Meka a reçu une médaille » $\rightarrow$ on peut répondre « c'est faux » : constatif.)
\item \textbf{Le test de "par la présente".} Insère mentalement « par la présente » dans la phrase. Si la phrase reste acceptable et solennelle, c'est un \textbf{performatif}. (« Je déclare, par la présente, la cérémonie ouverte » $\rightarrow$ correct : performatif. « Je marche, par la présente » $\rightarrow$ absurde : constatif.)
\end{enumerate}
En cas de doute, demande-toi enfin : « En disant cela, le locuteur décrit-il une action ou l'accomplit-il ? »
\end{methode}

\begin{popfigure}
\begin{center}
\begin{tikzpicture}[
  box/.style={draw=popblue, thick, rounded corners, fill=popblueL, text width=5.2cm, align=center, minimum height=1.1cm, font=\small},
  qbox/.style={draw=poporange, thick, rounded corners, fill=poporangeL, text width=4.6cm, align=center, minimum height=1.1cm, font=\small},
  res/.style={draw=popgreen, thick, rounded corners, fill=popgreenL, text width=3.4cm, align=center, minimum height=1cm, font=\small\bfseries},
  res2/.style={draw=poppurple, thick, rounded corners, fill=poppurpleL, text width=3.4cm, align=center, minimum height=1cm, font=\small\bfseries},
  arr/.style={-{Stealth[length=3mm]}, thick, popdark}
]
\node[qbox] (q1) at (0,0) {L'énoncé peut-il être jugé \textbf{vrai ou faux} ?};
\node[res, align=center] (c1) at (-4.6,-2.2){CONSTATIF\\ (il décrit le monde)};
\node[qbox, align=center] (q2) at (4.6,-2.2){Peut-on insérer « par la présente » ?\\ (je + verbe de parole + présent)};
\node[res2, align=center] (p1) at (4.6,-4.6){PERFORMATIF\\ (il accomplit l'acte)};
\node[res, align=center] (c2) at (0,-4.6){CONSTATIF\\ (première personne simple)};
\draw[arr] (q1) -- (c1) node[midway, above left, font=\small\bfseries, popgreen]{OUI};
\draw[arr] (q1) -- (q2) node[midway, above right, font=\small\bfseries, poporange]{NON};
\draw[arr] (q2) -- (p1) node[midway, right, font=\small\bfseries, poppurple]{OUI};
\draw[arr] (q2) -- (c2) node[midway, above left, font=\small\bfseries, popgreen]{NON};
\end{tikzpicture}
\end{center}
\legende{Arbre de décision : un énoncé est-il performatif ? Premier test : la vérité. Deuxième test : l'insertion de « par la présente ».}
\end{popfigure}

\courssub{Constatifs et performatifs dans \emph{Le Vieux nègre et la médaille}}

Le roman de Ferdinand Oyono met en scène une grande cérémonie officielle : la remise d'une médaille au vieux Meka pour services rendus à l'administration coloniale. Or toute cérémonie officielle est un festival d'énoncés performatifs : discours, déclarations, félicitations, remises solennelles. C'est une mine d'exemples pour notre leçon.

\begin{popfigure}
\begin{center}
\begin{tikzpicture}
\node[draw=popblue, thick, rounded corners, fill=popblueL, text width=6.6cm, align=left, font=\small, anchor=north west] (g) at (0,0) {
\textbf{\large Énoncés constatifs}\\[2pt]
\textbullet\ « Meka a donné ses terres à la mission. »\\[1pt]
\textbullet\ « Le vieil homme attend sous le soleil depuis le matin. »\\[1pt]
\textbullet\ « La médaille brille sur la poitrine de Meka. »\\[1pt]
\textbullet\ « Il a perdu sa médaille pendant la nuit. »\\[3pt]
\emph{Critère : vrai ou faux, vérifiable.}
};
\node[draw=poppurple, thick, rounded corners, fill=poppurpleL, text width=6.6cm, align=left, font=\small, anchor=north west] (d) at (7.6,0) {
\textbf{\large Énoncés performatifs}\\[2pt]
\textbullet\ « Je déclare la cérémonie ouverte. »\\[1pt]
\textbullet\ « Je vous félicite au nom de l'administration. »\\[1pt]
\textbullet\ « Nous vous remettons cette médaille. »\\[1pt]
\textbullet\ « Je vous remercie de vos services. »\\[3pt]
\emph{Critère : dire, c'est faire ; réussite ou échec.}
};
\end{tikzpicture}
\end{center}
\legende{Tableau à deux colonnes : énoncés constatifs et performatifs illustrés par des exemples inspirés de la cérémonie de remise de la médaille dans le roman d'Oyono.}
\end{popfigure}

Observons la cérémonie de plus près. Les discours officiels du commandant et des notables sont truffés de performatifs : on \emph{déclare}, on \emph{félicite}, on \emph{remercie}, on \emph{remet} la médaille. Chacun de ces verbes, prononcé par la bonne personne (le représentant de l'administration, qui a l'autorité) dans les bonnes circonstances (la fête officielle, devant la foule), accomplit réellement l'acte. Mais Oyono, avec l'ironie qu'on lui connaît, montre aussi la fragilité de ces beaux discours : les paroles officielles célèbrent la « fidélité » de Meka, alors que la réalité décrite par les constatifs du narrateur — la pauvreté du vieillard, ses terres perdues, sa médaille vite perdue elle aussi — dément cruellement les belles promesses coloniales. Le contraste entre les performatifs solennels des Blancs et les constatifs accablants du récit nourrit la tonalité satirique de l'œuvre, que nous étudierons dans la prochaine leçon de langue.

\courssub{Utilité pour le résumé : rapporter les énoncés au discours indirect}

Pourquoi cette distinction est-elle décisive pour la contraction de texte ? Parce qu'un résumé \textbf{rapporte} les paroles et les idées sans les reproduire : il les transpose au \cle{discours indirect}. Or la transposition diffère selon la nature de l'énoncé.

\begin{methode}[Rapporter un énoncé dans un résumé]
\begin{enumerate}
\item \textbf{Identifier} le type d'énoncé (constatif ou performatif) grâce aux deux tests.
\item \textbf{Constatif} : le rapporter avec un verbe déclaratif (\emph{dire que, affirmer que, constater que, rappeler que}) + subordonnée. Exemple : « La médaille est perdue » $\rightarrow$ \emph{Le narrateur rapporte que Meka a perdu sa médaille.}
\item \textbf{Performatif} : le rapporter avec un nom d'action ou un verbe qui nomme l'acte (\emph{ouvrir, féliciter, remercier, promettre}). Exemple : « Je déclare la cérémonie ouverte » $\rightarrow$ \emph{Le commandant ouvre la cérémonie} ou \emph{Le commandant déclare la cérémonie ouverte.}
\item \textbf{Ne jamais reproduire} les guillemets ni la première personne dans un résumé : on condense l'acte de parole en une proposition sobre.
\end{enumerate}
\end{methode}

\begin{attention}[Erreurs fréquentes dans le résumé]
\begin{itemize}
\item Recopier les discours officiels avec leurs guillemets : le résumé doit être entièrement au discours indirect, à la troisième personne.
\item Confondre l'acte accompli et sa description : « le commandant dit que la cérémonie est ouverte » (il la décrit) n'a pas le même sens que « le commandant ouvre la cérémonie » (il l'accomplit).
\item Oublier de neutraliser les marques de solennité (« par la présente », formules de politesse) qui alourdissent inutilement le résumé.
\end{itemize}
\end{attention}

\begin{exoresolu}[Constatif ou performatif ?]
Énoncé. Pour chacun des énoncés suivants, indique s'il est constatif ou performatif, en justifiant par le test de vérité ou le test de « par la présente » ; puis transpose-le au discours indirect comme dans un résumé.
\begin{enumerate}
\item « Je déclare la cérémonie ouverte », dit le commandant.
\item « Meka a donné ses terres à la mission catholique. »
\item « Je vous remercie de votre fidélité », dit l'officier à Meka.
\item « Je traverse le village à pied. »
\end{enumerate}
\tcblower
Solution détaillée.
\begin{enumerate}
\item \textbf{Performatif.} Le test de vérité échoue (on ne peut pas répondre « c'est faux ») ; le test de « par la présente » réussit (« je déclare, par la présente... »). Verbe de parole + première personne + présent, locuteur habilité. \emph{Résumé : Le commandant ouvre officiellement la cérémonie.}
\item \textbf{Constatif.} On peut vérifier et répondre « c'est faux » ; « par la présente » est impossible. \emph{Résumé : Le narrateur rappelle que Meka a donné ses terres à la mission.}
\item \textbf{Performatif.} Remercier se fait en parlant : la phrase accomplit le remerciement. \emph{Résumé : L'officier remercie Meka de sa fidélité.}
\item \textbf{Constatif.} Première personne du présent, mais « traverser » n'est pas un acte de parole : on peut dire « c'est faux, tu es assis ». \emph{Résumé : Le personnage dit qu'il traverse le village à pied.}
\end{enumerate}
\end{exoresolu}

\begin{exercice}[À toi de jouer]
Dans chaque cas, classe l'énoncé (constatif ou performatif), justifie ton choix, puis propose une formulation de résumé au discours indirect.
\begin{enumerate}
\item « Je te promets que je viendrai à la fête », dit Amalia à son mari.
\item « La pluie a détruit le pont pendant la nuit. »
\item « Je condamne l'accusé à deux ans de prison », dit le juge.
\item « Nous déclarons la récolte excellente cette année », écrivent les journalistes.
\end{enumerate}
\end{exercice}

\begin{corrige}[Exercice : À toi de jouer]
\begin{enumerate}
\item \textbf{Performatif} : promettre est un acte de parole accompli en le disant (« je te promets, par la présente » est acceptable). \emph{Résumé : Amalia promet à son mari de venir à la fête.}
\item \textbf{Constatif} : vérifiable, vrai ou faux ; « par la présente » impossible. \emph{Résumé : Le narrateur constate que la pluie a détruit le pont dans la nuit.}
\item \textbf{Performatif} : le juge, autorité compétente, accomplit la condamnation en la prononçant. \emph{Résumé : Le juge condamne l'accusé à deux ans de prison.}
\item \textbf{Constatif} (piège !) : malgré le verbe « déclarer », les journalistes n'ont pas le pouvoir de rendre la récolte excellente en parlant ; ils rapportent un fait vérifiable. \emph{Résumé : Les journalistes annoncent que la récolte est excellente cette année.}
\end{enumerate}
\end{corrige}

\begin{aretenir}[L'essentiel de la leçon]
\begin{itemize}
\item Un énoncé \cle{constatif} décrit le monde : il est \textbf{vrai ou faux}.
\item Un énoncé \cle{performatif} accomplit l'action qu'il énonce : il \textbf{réussit ou échoue}.
\item Marques du performatif : verbe d'acte de parole, première personne, présent, éventuellement « par la présente ».
\item Conditions de réussite : autorité du locuteur, circonstances appropriées, formule correctement exécutée.
\item Attention : « je marche » n'est pas performatif, car marcher n'est pas un acte de parole.
\item Dans un résumé, on rapporte tout au discours indirect : les constatifs avec « dire que », les performatifs avec le verbe ou le nom de l'acte (« ouvrir », « féliciter », « remercier »).
\end{itemize}
\end{aretenir}

Cette maîtrise des actes de langage nous servira doublement dans la suite de l'unité : d'abord pour repérer, dans les discours officiels du roman, le décalage ironique entre les paroles solennelles et la réalité vécue par Meka — c'est le ressort de la tonalité satirique que nous étudierons ensuite ; ensuite pour produire un résumé fidèle et élégant, où chaque parole du texte sera rapportée avec le verbe juste.

\coursec{Stylistique : les tonalités satirique et polémique}

Après avoir analysé la force illocutoire des énoncés — distinguer ce qui \emph{constate} (énoncé constatif) de ce qui \emph{agit} (énoncé performatif) —, nous devons maintenant scruter l'\textbf{attitude} globale de l'énonciateur. Un texte ne se contente pas d'informer ou d'agir ; il \emph{prend position}. C'est ce que l'on appelle la \cle{tonalité} (ou \cle{registre}). Dans \emph{Le Vieux nègre et la médaille}, Ferdinand Oyono ne se contente pas de raconter une cérémonie ; il juge, il raille, il accuse. Comprendre \emph{comment} il le fait est indispensable pour produire un résumé fidèle qui rend compte non seulement du \emph{quoi}, mais du \emph{comment}.

\courssub{La notion de tonalité : l'attitude de l'énonciateur perceptible dans le texte}

\begin{definition}[Tonalité (ou registre)]
La \textbf{tonalité} est l'attitude affective, intellectuelle et morale que l'énonciateur adopte à l'égard de son sujet et de son destinataire. Elle se manifeste par des choix lexicaux, syntaxiques, rhétoriques et ponctuels qui colorent l'ensemble du discours. On la perçoit comme une \emph{voix} qui résonne derrière les mots : ironique, lyrique, tragique, comique, polémique, satirique, etc.
\end{definition}

\begin{aretenir}[Identifier la tonalité]
Pour nommer la tonalité d'un texte, on procède par \textbf{convergence d'indices} :
\begin{enumerate}
    \item Le \textbf{champ lexical} (mots péjoratifs / mélioratifs, vocabulaire familier / soutenu).
    \item Les \textbf{figures de style} récurrentes (ironie, hyperbole, antiphrase...).
    \item La \textbf{syntaxe} (phrases courtes/longues, interrogations rhétoriques, exclamations).
    \item La \textbf{ponctuation} (points de suspension, points d'exclamation, tirets).
    \item L'\textbf{énonciation} (marques de subjectivité : \emph{je}, adverbes d'opinion, modalisateurs).
\end{enumerate}
C'est la \textbf{répétition} et la \textbf{cohérence} de ces indices qui fixent le registre dominant.
\end{aretenir}

Dans \emph{Le Vieux nègre et la médaille}, la tonalité générale est \textbf{satirique}, avec des pointes \textbf{polémiques} lors des confrontations directes (discours du commandant, sermons du prêtre, monologue final de Meka). Le rire n'est pas gratuit : il est une \cle{arme critique}.

\courssub{La tonalité satirique : railler les vices et les ridicules pour mieux les corriger}

\begin{definition}[Satire]
La \textbf{satire} est un genre (ou un registre) qui vise à \textbf{critiquer les mœurs, les institutions ou les travers humains} en les tournant en \cle{ridicule}. Elle utilise le rire comme instrument de \cle{correction} morale ou sociale. L'auteur satirique feint souvent d'adhérer à ce qu'il dénonce pour mieux en montrer l'absurdité (\cle{ironie}, \cle{antiphrase}).
\end{definition}

\begin{propriete}[Les quatre piliers de l'arsenal satirique chez Oyono]
Dans \emph{Le Vieux nègre et la médaille}, quatre procédés reviennent systématiquement pour construire l'effet satirique :
\begin{itemize}
    \item \textbf{L'ironie} : dire le contraire de ce que l'on pense, en comptant sur l'intelligence du lecteur pour décoder le vrai sens. \emph{Ex :} « Il était heureux comme un roi » (alors que Meka est humilié).
    \item \textbf{L'antiphrase} : figure d'ironie où un mot est employé dans un sens opposé à sa signification propre. \emph{Ex :} La « médaille » (symbole d'honneur) devient marque de \emph{servilité}.
    \item \textbf{La caricature} : grossissement outrancier des traits physiques ou moraux pour rendre le personnage ridicule. \emph{Ex :} Le commandant bedonnant, le prêtre gourmand.
    \item \textbf{L'hyperbole} : exagération volontaire pour souligner le décalage entre la réalité et la représentation officielle.
\end{itemize}
\end{propriete}

\begin{popfigure}
\begin{center}
\begin{tabularx}{\linewidth}{|l|X|X|X|}
\hline
\rowcolor{popblueL}
\textbf{Procédé satirique} & \textbf{Définition / Mécanisme} & \textbf{Exemple dans \emph{Le Vieux nègre et la médaille}} & \textbf{Effet produit} \\
\hline
\textbf{Ironie} & Dire le contraire de sa pensée ; double énonciation. & « Meka \emph{rayonnait} de bonheur » (p. 15) alors qu'il est terrassé par la fatigue et la honte. & Distance critique ; complicité lecteur/auteur. \\
\hline
\textbf{Antiphrase} & Mot/Expression utilisé à contre-sens (sous-type d'ironie). & La « \emph{médaille d'honneur} » remise à un vieillard exploité ; le « \emph{chef} » Meka sans autorité réelle. & Dénonciation de l'imposture coloniale. \\
\hline
\textbf{Caricature} & Grossissement outrancier de traits physiques/moraux. & Le Commandant : « \emph{ventre bedonnant}, \emph{face rubiconde}, \emph{voix de stentor} » (p. 22). Le Père Supérieur : « \emph{bedaine}, \emph{goinfre} ». & Désacralisation de l'autorité ; ridicule du pouvoir. \\
\hline
\textbf{Hyperbole} & Exagération volontaire pour frapper les esprits. & « \emph{Des milliers d'années} de civilisation » apportées par le colon ; la cérémonie décrite comme un « \emph{événement cosmique} ». & Mise en évidence du mensonge idéologique. \\
\hline
\end{tabularx}
\end{center}
\legende{Tableau récapitulatif des procédés satiriques dans \emph{Le Vieux nègre et la médaille} (F. Oyono).}
\end{popfigure}

\begin{exemplebox}[La médaille : antiphrase centrale du roman]
Le titre même de l'œuvre condense la satire. La \textbf{médaille} est, par définition, une récompense pour un \emph{mérite}, un \emph{courage}, un \emph{service exceptionnel}.
\begin{itemize}
    \item \textbf{Sens propre (discours officiel) :} La France récompense la « fidélité » de Meka.
    \item \textbf{Sens réel (discours satirique) :} La France « récompense » 50 ans d'exploitation, de travail forcé, de vol de terres, d'humiliations et la perte de ses deux fils.
\end{itemize}
L'antiphrase transforme l'objet symbolique en \textbf{preuve matérielle du crime colonial}. Meka ne reçoit pas une médaille ; il reçoit le \emph{reçu} de sa propre aliénation. Le vieil homme « \emph{porte} la médaille comme une croix » (double sens : décorations / supplice).
\end{exemplebox}

\courssub{La tonalité polémique : l'attaque frontale contre l'adversaire}

\begin{definition}[Polémique]
La \textbf{polémique} (du grec \emph{polemos}, « guerre ») est une \cle{argumentation agressive} qui vise à \textbf{réfuter violemment une thèse adverse} ou à \textbf{détruire la crédibilité d'un adversaire}. Contrairement à la satire qui rit \emph{de} son objet en le déformant, la polémique s'adresse \emph{à} un contradicteur réel ou supposé pour le terrasser. Le ton est grave, incisif, souvent injurieux (\cle{invective}).
\end{definition}

\begin{propriete}[Procédés polémiques majeurs]
\begin{itemize}
    \item \textbf{L'accumulation (énumération) :} Entasser les arguments, les griefs, les preuves pour accabler l'adversaire par le nombre et la masse. « Vous nous avez pris nos terres, nos femmes, nos fils, notre dignité... »
    \item \textbf{L'invective :} Insulte, injure directe, vocabulaire violemment péjoratif (\emph{voleurs, assassins, hypocrites, chiens}).
    \item \textbf{L'interrogation rhétorique (argumentative) :} Question qui n'attend pas de réponse car elle contient sa propre réponse accusatrice. « \emph{Où sont nos fils ?} » « \emph{Qu'avez-vous fait de notre or ?} »
    \item \textbf{L'apostrophe :} S'adresser directement à l'absent ou à l'abstrait pour l'accuser. « \emph{Ô France ! Ô Civilisation !} »
\end{itemize}
\end{propriete}

\begin{attention}[Piège fréquent : Confondre Satirique et Polémique]
\emph{Erreur :} « Le texte est polémique car l'auteur critique la colonisation. »
\emph{Correction :} La \textbf{satire} critique \textbf{par le rire et le détour} (ironie, caricature). La \textbf{polémique} critique \textbf{par l'attaque frontale et l'argumentation directe} (invective, accumulation).
\begin{itemize}
    \item \textbf{Passage satirique :} La description de la cérémonie (chap. 3-4) : le ridicule des officiels, l'absurdité du protocole. $\rightarrow$ \emph{Registre satirique.}
    \item \textbf{Passage polémique :} Le monologue intérieur de Meka à la fin (chap. 6) ou les répliques de Kelara : « \emph{Vous êtes des voleurs, des mangeurs d'hommes !} » $\rightarrow$ \emph{Registre polémique.}
\end{itemize}
Un même texte (comme le roman d'Oyono) peut \textbf{alterner} les deux : la satire prépare le terrain, la polémique porte le coup de grâce.
\end{attention}

\courssub{Analyse de passages : la satire de l'administration coloniale et du clergé}

Le roman opère une \textbf{double satire} : le pouvoir civil (administration) et le pouvoir spirituel (mission) sont complices dans l'oppression.

\begin{experience}[Analyse guidée : La scène de la remise de médaille (Chap. 4)]
\textbf{Extrait ciblé (reconstitution fidèle) :} « Le Commandant, dans son uniforme neuf, brillait de tous ses galons. Il posa sur la poitrine flétrie du vieux nègre la petite croix de métal... \emph{— C'est la France qui te remercie, Meka,} dit-il d'une voix émue. »
\begin{enumerate}
    \item \textbf{Caricature physique :} « \emph{uniforme neuf}, \emph{brillait de tous ses galons} » vs « \emph{poitrine flétrie du vieux nègre} ». Contraste saisissant : la puissance clinquante vs la vie usée.
    \item \textbf{Antiphrase performative :} L'énoncé « \emph{C'est la France qui te remercie} » est un \textbf{performatif} (voir section 3) \emph{mensonger}. Il \emph{devrait} créer une dette de gratitude ; il \emph{révèle} une dette de sang.
    \item \textbf{Ironie narrative :} Le narrateur rapporte l'émotion du Commandant (« \emph{d'une voix émue} ») que le lecteur sait feinte ou hypocrite.
    \item \textbf{Effet :} Le lecteur rit \emph{jaune}. Le ridicule tue l'autorité.
\end{enumerate}
\end{experience}

\begin{savaistu}[Contexte historique : La « médaille de la fidélité » au Cameroun]
Sous administration française (mandat SDN puis tutelle ONU), l'administration coloniale décernait effectivement des \textbf{médailles d'honneur} ou de la \textbf{médaille de la fidélité} aux chefs traditionnels et « évolués » loyaux. C'était un outil de \cle{gouvernance indirecte} : acheter la collaboration des élites locales pour asseoir l'ordre colonial. Oyono, fils de notable et ancien fonctionnaire colonial, connaît parfaitement ce mécanisme. La satire du roman est donc une \textbf{contre-histoire} documentée.
\end{savaistu}

\begin{experience}[Analyse guidée : La satire du clergé (Chap. 2 - Le repas du Père Supérieur)]
\textbf{Extrait ciblé :} Le Père Supérieur mange « \emph{avec une voracité de loup} », « \emph{sa bedaine rebondissait} », il boit « \emph{du vin de messe} » en quantité. Il bénit la médaille « \emph{au nom du Père, du Fils et du Saint-Esprit} » puis réclame sa part de poulet.
\begin{itemize}
    \item \textbf{Caricature morale :} Le prêtre incarne la \cle{gloutonnerie} (péché capital) au lieu de la charité.
    \item \textbf{Parodie liturgique :} La bénédiction de la médaille (objet profane, symbole d'oppression) parodie la consécration eucharistique. Le sacré est instrumentalisé pour légitimer le profane (le pouvoir colonial).
    \item \textbf{Complicité :} L'Église mange à la table du Commandant. Satire de l'alliance \emph{Trône et Autel}.
\end{itemize}
\end{experience}

\begin{aretenir}[Effets produits sur le lecteur : le rire comme arme politique]
\begin{enumerate}
    \item \textbf{Désacralisation :} Le rire ôte au pouvoir colonial son aura de légitimité, de grandeur, de mission civilisatrice. Le Commandant devient un bouffon bedonnant ; le Prêtre, un goinfre.
    \item \textbf{Prise de distance critique :} L'ironie force le lecteur (camerounais des années 50, ou élève d'aujourd'hui) à ne pas subir le discours officiel mais à l'\emph{analyser}.
    \item \textbf{Identification à la victime :} En riant \emph{avec} Meka (ou de ses bourreaux), le lecteur partage sa lucidité. La satire crée une \cle{communauté de lecture} contre l'oppresseur.
    \item \textbf{Transition vers le tragique :} Le rire satirique s'effondre à la fin (mort de Meka, prise de conscience de Kelara). La satire laisse place à la \textbf{polémique} puis au \textbf{tragique}. Le rire était une \emph{résistance} ; le silence final est une \emph{révolte}.
\end{enumerate}
\end{aretenir}

\courssub{De la tonalité au résumé : signaler sans reproduire}

C'est un point \textbf{crucial pour l'examen} : dans un résumé (contraction de texte), on \textbf{ne doit jamais imiter la tonalité} de l'auteur (faire de l'ironie soi-même), mais on doit \textbf{la signaler et la qualifier} avec un vocabulaire précis.

\begin{methode}[Identifier et rendre compte de la tonalité dans un résumé]
\begin{enumerate}
    \item \textbf{Repérer les indices} (champ lexical, figures, syntaxe, ponctuation) $\rightarrow$ \emph{Quels sont les procédés dominants ?}
    \item \textbf{Nommer le registre} $\rightarrow$ \emph{Satirique ? Polémique ? Ironique ? Tragique ? Lyrique ?}
    \item \textbf{Préciser la cible} $\rightarrow$ \emph{De quoi rit-on ? Qui est attaqué ? (L'administration, le clergé, l'idéologie coloniale).}
    \item \textbf{Formuler au style indirect / rapporté} en utilisant des verbes de \cle{parole intellectuelle} ou \cle{d'attitude} :
    \begin{itemize}
        \item \emph{L'auteur \textbf{raille} / \textbf{tourne en dérision} / \textbf{dénonce par l'ironie} / \textbf{caricature}...}
        \item \emph{L'auteur \textbf{attaque violemment} / \textbf{accable} / \textbf{réfute} / \textbf{invective}...}
        \item \emph{Le texte adopte un \textbf{ton satirique} / \textbf{une tonalité polémique}...}
    \end{itemize}
    \item \textbf{Intégrer cette mention} dans la phrase qui résume le passage concerné.
\end{enumerate}
\end{methode}

\begin{exoresolu}[Du texte au résumé : rendre compte de la satire]
\textbf{Extrait (imaginé, style Oyono) :}
« Le Grand Chef Blanc, ventre à l'air, décorations au vent, descendit de sa voiture climatisée. Il tendit la breloque clinquante au vieillard courbé sous le poids des ans et des impôts. \emph{— Reçois, mon fils, le gage de notre amour éternel,} psalmodia-t-il. Le vieux ne bronchait pas, les yeux secs. »

\textbf{Question :} Rédiger la phrase de résumé correspondante en signalant la tonalité.

\textbf{Solution détaillée :}
\begin{enumerate}
    \item \textbf{Indices :} « ventre à l'air », « breloque clinquante », « psalmodia », « gage de notre amour éternel » (antiphrase), contraste physique.
    \item \textbf{Registre :} \textbf{Satirique} (caricature + ironie + antiphrase).
    \item \textbf{Cible :} L'hypocrisie de l'autorité coloniale (le « Grand Chef Blanc »).
    \item \textbf{Rédaction du résumé :}
    > \emph{« Dans une scène à la \textbf{tonalité satirique}, l'auteur \textbf{caricature} l'autorité coloniale (le "Grand Chef Blanc") et \textbf{dénonce par l'ironie} l'hypocrisie du discours officiel ("gage d'amour éternel") qui masque l'exploitation réelle (le vieux "courbé sous le poids des ans et des impôts"). »}
\end{enumerate}
\textbf{Erreur à ne pas commettre :} « Le Grand Chef Blanc, tout bedonnant, donne une breloque au vieux en lui mentant qu'il l'aime. » $\rightarrow$ \emph{Ici, l'élève \textbf{reproduit} la satire (vocabulaire familier "breloque", "bedonnant", jugement "mentant") au lieu de la \textbf{rapporter}.}
\end{exoresolu}

\courssub{Synthèse visuelle : de la dérision légère à l'attaque polémique violente}

Les deux registres ne sont pas étanches ; ils forment un \textbf{continuum de l'agressivité verbale}. Le schéma ci-dessous positionne les procédés selon leur intensité conflictuelle.

\begin{popfigure}
\begin{center}
\begin{tikzpicture}[
    axis/.style={very thick, ->, >=Stealth, popdark},
    tick/.style={thick, popdark},
    label/.style={font=\small\bfseries, text=popdark, anchor=north},
    box/.style={rounded corners=4pt, inner sep=6pt, font=\small, align=center},
    satirique/.style={box, fill=popgreenL, draw=popgreen, text=popgreen!80!black},
    transition/.style={box, fill=poporangeL, draw=poporange, text=poporange!80!black},
    polemique/.style={box, fill=poppinkL, draw=poppink, text=poppink!80!black},
    proc/.style={font=\footnotesize, text=popmuted, align=center, anchor=north}
]
% Axe principal
\draw[axis] (0,0) -- (14,0) node[right, label] {Intensité de l'agressivité verbale $\uparrow$};

% Graduations
\foreach \x/\txt in {1/Légère, 3.5/Modérée, 6/Forte, 8.5/Très forte, 11/Extrême, 13/Paroxysme}
    \draw[tick] (\x,0.15) -- (\x,-0.15) node[below, font=\tiny, text=popmuted] {\txt};

% Zones colorées (fond)
\fill[popgreenL!30] (0,-1.2) rectangle (5,1.2);
\fill[poporangeL!30] (5,-1.2) rectangle (9,1.2);
\fill[poppinkL!30] (9,-1.2) rectangle (14,1.2);

% Étiquettes de zones
\node[satirique, align=center] at (2.5, 1.5){ZONE \textbf{SATIRIQUE}\\(Rire critique, distance)};
\node[transition, align=center] at (7, 1.5){ZONE \textbf{MIXTE}\\(Ironie aigre, sarcasme)};
\node[polemique, align=center] at (11.5, 1.5){ZONE \textbf{POLÉMIQUE}\\(Attaque frontale, rupture)};

% Exemples de procédés placés sur l'axe
% Satirique
\node[proc] at (1, -0.6) {Sous-entendu};
\node[proc] at (2, -0.6) {Antiphrase};
\node[proc] at (3, -0.6) {Ironie douce};
\node[proc] at (4, -0.6) {Caricature};
\node[proc] at (5, -0.6) {Parodie};

% Transition / Sarcasme
\node[proc, text=poporange!80!black] at (6, -0.6) {Sarcasme};
\node[proc, text=poporange!80!black] at (7, -0.6) {Ironie cinglante};
\node[proc, text=poporange!80!black] at (8, -0.6) {Réduction à l'absurde};

% Polémique
\node[proc, text=poppink!80!black] at (9, -0.6) {Accumulation};
\node[proc, text=poppink!80!black] at (10, -0.6) {Interrog. rhétorique};
\node[proc, text=poppink!80!black] at (11, -0.6) {Invective};
\node[proc, text=poppink!80!black] at (12, -0.6) {Anathème / Excommunication};
\node[proc, text=poppink!80!black] at (13, -0.6) {Imprécation};

% Flèches indicatives
\draw[popgreen, thick, <->] (0.5, -1) -- (4.5, -1) node[midway, below, font=\tiny, popgreen] {Détour / Second degré};
\draw[poporange, thick, <->] (5.5, -1) -- (8.5, -1) node[midway, below, font=\tiny, poporange] {Ambiguïté calculée};
\draw[poppink, thick, <->] (9.5, -1) -- (13.5, -1) node[midway, below, font=\tiny, poppink] {Frontal / Premier degré};

% Positionnement "Vieux Nègre"
\node[draw=popblue, thick, fill=popblueL, rounded corners, inner sep=4pt, font=\small\bfseries, anchor=south] at (7, -2.2) {\emph{Le Vieux nègre et la médaille} : Principalement 2 $\rightarrow$ 5 (Satire), pointes à 9-10 (Polémique finale)};
\end{tikzpicture}
\end{center}
\legende{Continuum des tonalités critiques : de la satire (détour par le rire) à la polémique (affrontement direct). Le roman d'Oyono navigue principalement en zone satirique pour basculer en zone polémique au dénouement.}
\end{popfigure}

\begin{aretenir}[Bilan pour le résumé (Intégration)]
\begin{itemize}
    \item \textbf{Satirique} = \emph{L'auteur raille / tourne en dérision / caricature / use d'ironie / d'antiphrase pour dénoncer [Cible].}
    \item \textbf{Polémique} = \emph{L'auteur attaque / accable / invective / réfute violemment [Adversaire/Thèse] par l'accumulation / l'interrogation rhétorique.}
    \item \textbf{Mixte} = \emph{Le texte use d'un ton satirique qui vire au polémique lors de [Passage clé].}
\end{itemize}
Maîtriser ce vocabulaire métalangagier (\cle{railler}, \cle{dénoncer ironiquement}, \cle{caricaturer}, \cle{accabler}, \cle{invectiver}) est la marque d'un résumé de \textbf{niveau Probatoire/Baccalauréat Technique}.
\end{aretenir}

\coursec{Production et amélioration du résumé}

Après avoir analysé les tonalités satirique et polémique qui traversent \emph{Le Vieux nègre et la médaille}, nous disposons désormais des clés de lecture pour saisir la charge critique de l'œuvre. Cette compréhension fine est le socle indispensable : on ne peut résumer fidèlement un texte dont on n'a pas perçu l'intention profonde, ses ironies, ses déplacements de sens. La contraction de texte n'est pas un exercice de raccourcissement mécanique ; c'est un acte de \cle{reformulation intelligente} qui suppose d'avoir identifié la thèse, l'architecture argumentative et les procédés stylistiques qui la portent. Cette section vous guide pas à pas, de la première lecture à la version finale prête pour l'épreuve du Probatoire, en intégrant les exigences de justification qui font la différence entre un résumé correct et un résumé excellent.

\courssub{Les cinq étapes de la production du résumé}

La rédaction d'un résumé suit un processus itératif en cinq temps. Chaque étape a son objectif propre et ses critères de réussite. Le schéma ci-dessous synthétise cette chaîne opératoire ; gardez-le en mémoire comme une \cle{check-list mentale} le jour de l'examen.

\begin{popfigure}
\begin{tikzpicture}[
  node distance=1.2cm and 1.5cm,
  stepbox/.style={rectangle, rounded corners, draw=popblue, thick, fill=popblueL, minimum width=2.8cm, minimum height=1.1cm, align=center, font=\small},
  arrow/.style={-Stealth, thick, popdark},
  labelnode/.style={font=\tiny, text=popmuted, align=center}
]
\node[stepbox, align=center] (lire){\textbf{1. Lire \& comprendre}\\(2--3 min)\\\textit{Repérer thèse, plan, ton}};
\node[stepbox, right=of lire, align=center] (sel){\textbf{2. Sélectionner}\\(3--4 min)\\\textit{Idées forces + articulations}};
\node[stepbox, right=of sel, align=center] (ref){\textbf{3. Reformuler}\\(5--7 min)\\\textit{Propre vocabulaire, DI}};
\node[stepbox, right=of ref, align=center] (red){\textbf{4. Rédiger}\\(8--10 min)\\\textit{Présent, impersonnel, fluide}};
\node[stepbox, right=of red, align=center] (cor){\textbf{5. Corriger}\\(3--5 min)\\\textit{Nombre de mots, langue, fidélité}};
\draw[arrow] (lire) -- (sel);
\draw[arrow] (sel) -- (ref);
\draw[arrow] (ref) -- (red);
\draw[arrow] (red) -- (cor);
\node[labelnode, below=0.3cm of lire] {Question clé :\newline« Quel est le message central ? »};
\node[labelnode, below=0.3cm of sel] {Question clé :\newline« Qu'est-ce qui est \emph{indispensable} ? »};
\node[labelnode, below=0.3cm of ref] {Question clé :\newline« Comment le dire \emph{autrement} sans trahir ? »};
\node[labelnode, below=0.3cm of red] {Question clé :\newline« Le texte tient-il seul, sans l'original ? »};
\node[labelnode, below=0.3cm of cor] {Question clé :\newline« Respecte-t-on la contrainte de longueur ? »};
\end{tikzpicture}
\legende{Organigramme des 5 étapes de production du résumé (durées indicatives pour un texte de 400--500 mots au Bac).}
\end{popfigure}

\begin{methode}[Rédiger le résumé en 5 étapes chronométrées — méthode Bac]
\textbf{Étape 1 — Lecture active (2--3 min).} Lisez le texte une première fois sans écrire. Identifiez : le thème, la thèse (quelle position l'auteur défend-il ?), le plan (combien de parties ? quelles transitions ?), le ton (ici : satirique, polémique). Soulignez \emph{uniquement} les mots-clés qui portent la thèse.
\textbf{Étape 2 — Sélection hiérarchisée (3--4 min).} Relisez en notant dans la marge : \textcircled{1} idées forces (thèses principales), \textcircled{2} arguments clés, \textcircled{3} connecteurs logiques (donc, cependant, en somme). Éliminez : exemples longs, citations, détails anecdotiques, répétitions.
\textbf{Étape 3 — Reformulation au brouillon (5--7 min).} Réécrivez chaque idée force \emph{avec vos mots}, au \cle{discours indirect} (il affirme que..., l'auteur montre que...), au \cle{présent de vérité générale}. Fusionnez les phrases courtes, supprimez les relativités inutiles.
\textbf{Étape 4 — Rédaction propre (8--10 min).} Recopiez en un paragraphe unique (sauf si le texte source a des parties très marquées justifiant deux paragraphes). Vérifiez l'enchaînement : connecteurs, anaphores, progression logique. Pas de « je », pas de « nous », pas de jugements personnels.
\textbf{Étape 5 — Contrôle final (3--5 min).} \textbf{Comptez les mots} (voir méthode ci-dessous). Vérifiez : fidélité au sens, cohérence interne, correction orthographique/grammaticale, respect de la contrainte de longueur (±10\%).
\end{methode}

\courssub{Rédaction : règles d'or et conventions}

Le résumé obéit à des conventions strictes qui le distinguent de toute autre production écrite. Les maîtriser, c'est s'assurer la note de \cle{forme} (souvent 4 à 6 points sur 20 au Probatoire).

\begin{propriete}[Conventions de rédaction du résumé]
Le résumé se rédige :
\begin{itemize}
\item au \textbf{présent de vérité générale} (l'auteur \emph{dénonce}, il \emph{montre}, le texte \emph{met en lumière}) ;
\item au \textbf{discours indirect} (pas de guillemets, pas de citations longues ; on rapporte : « l'auteur affirme que... », « il souligne que... ») ;
\item à la \textbf{troisième personne} (impersonnalité totale : pas de « je pense que », « à mon avis ») ;
\item en \textbf{phrases complètes et syntaxiquement autonomes} (le résumé doit être compréhensible sans avoir lu le texte source) ;
\item en \textbf{un paragraphe unique} (sauf instruction contraire ou texte source très long à structure marquée) ;
\item \textbf{sans commentaire, sans développement personnel, sans exemple ajouté}.
\end{itemize}
\end{propriete}

\begin{attention}[Piège fréquent — L'intrusion du sujet résumant]
\emph{« Je trouve que l'auteur a raison de critiquer la colonisation. »} $\rightarrow$ \textbf{Zéro point} sur la partie « fidélité / impersonnalité ».
\emph{« L'auteur critique la colonisation. »} $\rightarrow$ \textbf{Correct}.
Le résumé rend compte de \textbf{ce que le texte dit}, pas de \textbf{ce que vous en pensez}. Même si le texte est polémique (comme \emph{Le Vieux nègre et la médaille}), votre résumé reste \emph{constatif} : il rapporte la polémique, il ne la prolonge pas.
\end{attention}

\begin{exemplebox}[Application sur un extrait du \emph{Vieux nègre et la médaille}]
\textbf{Texte source (simplifié) :} « Le vieux nègre, décoré de la médaille coloniale, exhibe sa distinction avec une fierté naïve qui force le rire et la pitié. La médaille, symbole de sa soumission, brille sur sa poitrine comme un miroir aux alouettes. »
\\
\textbf{Respect des conventions :}
\begin{itemize}
\item \textit{Présent / DI / 3e personne :} « L'auteur dépeint un vieux décoré qui arbore fièrement sa médaille coloniale, symbole ironique de son aliénation. »
\item \textit{Pas de citation :} on ne garde pas « miroir aux alouettes » entre guillemets ; on reformule : « symbole illusoire » ou « leurre ».
\item \textit{Impersonnalité :} aucune trace du lecteur.
\end{itemize}
\end{exemplebox}

\courssub{Maîtriser la longueur : compter, ajuster, calibrer}

Au Probatoire, la consigne précise généralement : « Résumez le texte en \textbf{X mots} (±10\%) ». Une erreur de comptage ou un non-respect de la fourchette coûte cher (souvent 2 à 4 points). Voici la méthode fiable.

\begin{methode}[Compter et ajuster le nombre de mots]
\textbf{1. Compter les mots du texte source.} Règle officielle : un mot = toute suite de caractères séparée par un espace. Les contractions (c'est, n'est, qu'il) comptent pour \textbf{un mot}. Les nombres écrits en chiffres (1945) comptent pour \textbf{un mot}. La ponctuation ne compte pas.
\textbf{2. Calculer la cible.} Exemple : consigne « 80 mots ±10\% » $\rightarrow$ fourchette \textbf{72 à 88 mots}. Visez le milieu (80).
\textbf{3. Compter votre brouillon.} Numérotez chaque mot au stylo (1, 2, 3...) ou utilisez la méthode des lignes : comptez les mots d'une ligne type, multipliez par le nombre de lignes.
\textbf{4. Ajuster.}
\begin{itemize}
\item \textbf{Trop long ?} Compression : supprimez les adjectifs non essentiels, fusionnez deux phrases par une subordination, remplacez une périphrase par un verbe précis (ex. « il fait en sorte de » $\rightarrow$ « il veille à »).
\item \textbf{Trop court ?} Expansion contrôlée : rétablissez un connecteur logique supprimé, précisez un agent implicite (ex. « on dit » $\rightarrow$ « l'auteur affirme »), développez une conséquence logique implicite.
\end{itemize}
\textbf{5. Recompter la version finale.} Recomptez \emph{après} la recopie propre. Une erreur de 1--2 mots est tolérée dans la fourchette ; au-delà, la pénalité s'applique.
\end{methode}

\begin{exemplebox}[Compression / Expansion — Exemple chiffré]
\textbf{Cible : 50 mots ±10\% (45--55).}
\\
\textbf{Brouillon (62 mots) :} « L'auteur décrit un vieux nègre qui porte fièrement sa médaille coloniale. Cette médaille représente en réalité sa soumission au colonisateur. Le ton est satirique car l'auteur se moque de la naïveté du vieux. Il montre aussi que la colonisation a détruit la dignité des Africains. La médaille brille comme un leurre. »
\\
\textbf{Compression (50 mots) :} « L'auteur dépeint un vieux Africain arborant fièrement sa médaille coloniale, symbole ironique de son aliénation. Par une satire cinglante, il dénonce la naïveté du décoré et révèle comment la colonisation a spolié la dignité des peuples, transformant la distinction en leurre brillant. » \textit{(50 mots exacts)}
\end{exemplebox}

\courssub{Relecture critique et grille d'auto-évaluation}

La relecture n'est pas une relecture « impressionniste ». Elle s'appuie sur une \cle{grille critériée} que vous devez intérioriser. Au Bac, l'évaluateur utilise une grille voisine ; l'auto-évaluation et l'évaluation croisée entre pairs (en classe) sont les meilleurs entraînements.

\begin{aretenir}[Grille d'auto-évaluation / évaluation croisée (4 critères, 5 pts chacun = 20 pts)]
\begin{center}
\begin{tabularx}{\linewidth}{|l|X|X|X|X|}\hline
\textbf{Critère} & \textbf{5 pts — Excellent} & \textbf{4 pts — Bon} & \textbf{2--3 pts — Moyen} & \textbf{0--1 pt — Insuffisant} \\ \hline
\textbf{Fidélité} & Thèse + arguments clés + articulations respectés ; sens exact & Thèse + idées principales ; 1 omission mineure & Thèse approximative ; idées secondaires confondues avec principales & Contre-sens ; idées inventées ; avis personnel \\ \hline
\textbf{Reformulation} & Vocabulaire propre, DI maîtrisé, présent VG, phrases fluides & Reformulation correcte, 1--2 maladresses (citation résiduelle, temps) & Reformulation pauvre, calque syntaxique, passages en DI/DS mêlés & Copier-coller ; citations longues ; français approximatif \\ \hline
\textbf{Cohérence} & Enchaînement logique, connecteurs variés, anaphores claires, paragraphe uni & Liens logiques présents, 1 rupture mineure & Enchaînement heurté, répétitions, anaphores floues & Incohérent ; liste d'idées sans liens ; illisible \\ \hline
\textbf{Langue / Longueur} & 0 faute ; longueur dans la fourchette & 1--2 fautes mineures ; longueur dans la fourchette & 3--5 fautes ; longueur hors fourchette (±10\% non respecté) & $>$5 fautes ; longueur très éloignée ; texte incompréhensible \\ \hline
\end{tabularx}
\end{center}
\end{aretenir}

\begin{methode}[Auto-évaluation en 3 minutes (à faire sur votre brouillon avant recopie)]
\begin{enumerate}
\item \textbf{Fidélité :} Cochez dans le texte source chaque idée reprise. Y a-t-il une idée force manquante ? Une idée secondaire promue au rang principal ?
\item \textbf{Reformulation :} Surlignez les verbes introducteurs (affirmer, montrer, souligner, démontrer, ironiser...). Sont-ils variés ? Le discours indirect est-il constant ?
\item \textbf{Cohérence :} Lisez à voix haute (mentalement). Les connecteurs (ainsi, toutefois, en outre, par conséquent) font-ils avancer la démonstration ?
\item \textbf{Langue :} Chasse aux accords (sujet/verbe, participe passé), aux homophones (à/a, et/est, ou/où, son/sont), à la ponctuation.
\item \textbf{Longueur :} Recomptez. Si hors fourchette $\rightarrow$ retour compression/expansion (étape 4).
\end{enumerate}
\end{methode}

\courssub{Amélioration par la réécriture : du résumé faible au résumé réussi}

Rien ne vaut l'analyse contrastée d'un résumé insuffisant et de sa version corrigée. L'exercice ci-dessous met en lumière les erreurs types et les leviers de progression.

\begin{exoresolu}[Du résumé faible au résumé amélioré — Analyse annotée]
\textbf{Texte source (extrait fictif, style \emph{Vieux nègre et la médaille}) :}
« Le vieux Tchicaya, médaillé de l'Ordre du Mérite colonial, défile au 14 juillet. Sa poitrine bombée, son sourire béat disent sa fierté. Mais la médaille pèse. Elle pèse le poids des terres confisquées, des impôts forcés, des fils partis au front pour une patrie qui les méprise. Le colonisateur tape dans le dos du "bon indigène" pendant que l'administrateur signe l'arrêté d'expropriation. La médaille n'honore pas ; elle marque le bétail. »

\textbf{Consigne :} Résumer en 50 mots ±10\% (45--55 mots).

\medskip
\textbf{Résumé faible (élève A — 38 mots, hors fourchette, note estimée 7/20) :}
« Le vieux Tchicaya est fier de sa médaille coloniale. Mais l'auteur dit que c'est ironique car la médaille cache la spoliation des terres et l'exploitation. Le colonisateur flatte l'indigène mais lui vole ses terres. La médaille marque les soumis. Je trouve ça révoltant. »

\medskip
\textbf{Diagnostic (grille) :}
\begin{itemize}
\item \textbf{Fidélité (2/5) :} Thèse présente mais arguments réduits ; « l'auteur dit que » (DS) ; avis personnel final (« Je trouve ça révoltant »).
\item \textbf{Reformulation (1/5) :} Discours direct/indirect mêlés ; vocabulaire proche du texte (« cache », « flatte », « vole ») ; phrase « Je trouve... » interdite.
\item \textbf{Cohérence (2/5) :} Enchaînement basique (Mais, car) ; pas de connecteur de conséquence ; « La médaille marque les soumis » isolé.
\item \textbf{Langue/Longueur (2/5) :} 38 mots (hors fourchette basse) ; « l'auteur dit que » familier ; pas de faute grave mais style faible.
\end{itemize}

\tcblower

\textbf{Résumé amélioré (version corrigée — 51 mots, dans la fourchette, note estimée 18/20) :}
« L'auteur dépeint le vieux Tchicaya, fier exhibant sa médaille coloniale lors du défilé du 14 juillet. Par une ironie cinglante, il révèle que cette distinction masque la spoliation des terres, l'exploitation fiscale et le sacrifice des jeunes pour une métropole méprisante. La médaille, loin d'honorer, stigmatise la soumission du "bon indigène" complice de son propre asservissement. »

\medskip
\textbf{Analyse des corrections (leviers de progression) :}
\begin{tabularx}{\linewidth}{|l|X|X|}\hline
\textbf{Problème initial} & \textbf{Correction apportée} & \textbf{Principe mobilisé} \\ \hline
Avis personnel (« Je trouve... ») & Supprimé ; remplacé par « L'auteur révèle / dépeint » & Impersonnalité / DI \\ \hline
Discours direct (« l'auteur dit que ») & « L'auteur révèle que / dépeint... » & Verbes introducteurs variés + DI \\ \hline
Vocabulaire calqué (« cache », « flatte », « vole ») & « Masque », « distinction », « spoliation », « sacrifice », « stigmatise » & Reformulation lexicale précise \\ \hline
Connecteurs faibles (Mais, car) & « Par une ironie cinglante », « Loin d'honorer » & Connecteurs à valeur argumentative \\ \hline
Longueur 38 mots (trop court) & Expansion contrôlée : précision des mécanismes (fiscalité, sacrifice, complicité) & Expansion informative, non ornementale \\ \hline
Absence du ton satirique & « Ironie cinglante », « fierté béate » (impliqué par « fier exhibant ») & Fidélité au ton de l'œuvre étudiée \\ \hline
\end{tabularx}
\end{exoresolu}

\begin{popfigure}
\begin{tikzpicture}[
  node distance=0.8cm and 1cm,
  box/.style={rectangle, draw=popdark, thick, rounded corners, fill=popcream, minimum width=7cm, text width=6.5cm, align=left, font=\small},
  arrow/.style={-Stealth, popgreen, thick}
]
\node[box, align=center] (faible){\textbf{Résumé FAIBLE (38 mots)}\\$\times$ Avis personnel (« Je trouve »)\\$\times$ Discours direct (« l'auteur dit que »)\\$\times$ Vocabulaire calqué\\$\times$ Connecteurs basiques\\$\times$ Longueur hors fourchette\\$\times$ Ton satirique absent};
\node[box, right=2.5cm of faible, align=center] (fort){\textbf{Résumé AMÉLIORÉ (51 mots)}\\$\checkmark$ Impersonnalité totale\\$\checkmark$ Discours indirect varié\\$\checkmark$ Reformulation lexicale précise\\$\checkmark$ Connecteurs argumentatifs\\$\checkmark$ Longueur dans la fourchette\\$\checkmark$ Ton satirique restitué};
\draw[arrow, bend left=15] (faible.east) to node[midway, above, font=\small, text=popgreen] {Réécriture guidée par la grille} (fort.west);
\end{tikzpicture}
\legende{Comparaison annotée : résumé faible vs résumé amélioré — chaque étiquette renvoie à un critère de la grille d'évaluation.}
\end{popfigure}

\courssub{Justifier sa démarche : l'exigence du savoir-faire}

Le programme de Terminale technique insiste sur le savoir-faire : « Rédiger et justifier une démarche, une réponse ou une conclusion ». Au Probatoire, il peut vous être demandé, en fin d'épreuve ou en oral, d'\textbf{expliquer vos choix de sélection et de reformulation}. Cette justification n'est pas un exercice accessoire : elle prouve que votre résumé est le fruit d'un processus conscient, non du hasard.

\begin{exemplebox}[Modèle de justification orale/écrite (3--4 minutes)]
\textbf{Structure type :}
\begin{enumerate}
\item \textbf{Identification de la thèse source :} « Le texte de Mongo Beti dénonce l'aliénation du vieux décoré par la médaille coloniale, symbole de sa complicité involontaire. »
\item \textbf{Choix de sélection :} « J'ai retenu trois idées forces : (1) la fierté naïve du vieux, (2) l'ironie de la médaille qui masque la spoliation, (3) la complicité du "bon indigène". J'ai écarté le détail du défilé du 14 juillet (circonstance) et la description physique (secondaire). »
\item \textbf{Choix de reformulation :} « J'ai remplacé "sourire béat" par "fierté naïve" pour garder la nuance ironique. J'ai transformé "la médaille pèse le poids des terres confisquées" en "la médaille masque la spoliation" pour condenser l'image métaphorique en concept. J'ai utilisé les verbes "dépeint", "révèle", "stigmatise" pour varier l'introduction des idées. »
\item \textbf{Gestion de la longueur :} « Mon brouillon faisait 62 mots. J'ai compressé en fusionnant les deux premières phrases et en remplaçant "le colonisateur tape dans le dos..." par "pour une métropole méprisante" (gain de 12 mots). »
\item \textbf{Vérification finale :} « J'ai recompté : 51 mots. J'ai relu pour les accords et la ponctuation. Le résumé est autonome : on comprend la thèse sans avoir lu le texte. »
\end{enumerate}
\end{exemplebox}

\begin{aretenir}[Grille d'évaluation au Probatoire / Baccalauréat technique]
La grille officielle du Probatoire / Baccalauréat technique répartit généralement les 20 points ainsi (variable selon académies, mais structure stable) :
\begin{center}
\begin{tabular}{|l|c|c|}\hline
\textbf{Critère} & \textbf{Points} & \textbf{Indicateurs clés} \\ \hline
Fidélité au fond (thèse, arguments, articulations) & 5--6 & Idées forces toutes présentes, pas de contre-sens, pas d'ajout \\ \hline
Qualité de la reformulation (vocabulaire, DI, temps) & 5--6 & Vocabulaire propre, DI constant, présent VG, pas de citation \\ \hline
Cohérence et organisation (connecteurs, progression) & 4--5 & Texte uni, logique, anaphores claires, paragraphage pertinent \\ \hline
Correction linguistique (orthographe, grammaire, syntaxe) & 3--4 & 0--1 faute = max ; $>$3 fautes = pénalité forte \\ \hline
Respect de la contrainte de longueur (±10\%) & 1--2 (souvent bonus/malus) & Dans la fourchette = points pleins ; hors fourchette = 0 ou malus \\ \hline
\textbf{Total} & \textbf{20} &  \\ \hline
\end{tabular}
\end{center}
\textbf{Conseil stratégique :} Les 10--12 points de « fidélité + reformulation » se gagnent \emph{avant} la rédaction finale (étapes 1--3). La cohérence et la langue se soignent à l'étape 4--5. Ne sacrifiez jamais l'étape de sélection/réformulation au brouillon pour « gagner du temps » : c'est là que se joue la note.
\end{aretenir}

\courssub{Intégration et transfert : une compétence pour la vie professionnelle}

La contraction de texte n'est pas un exercice scolaire enfermé dans l'épreuve de français. C'est une \cle{compétence transverse} à haute valeur ajoutée dans le monde professionnel et administratif camerounais.

\begin{savaistu}[Savoir résumer : la compétence n°1 des concours administratifs camerounais (ENAM, IRIC, ENSET, concours professionnels)]
\begin{itemize}
\item \textbf{ENAM (École Nationale d'Administration et de Magistrature) :} L'épreuve de « Note de synthèse » exige de condenser un dossier de 20--30 pages en 3--4 pages structurées. C'est une contraction de texte \emph{à grande échelle} : même méthode (sélection $\rightarrow$ reformulation $\rightarrow$ rédaction), même exigence de fidélité, d'impersonnalité et de longueur contrainte.
\item \textbf{IRIC (Institut des Relations Internationales du Cameroun) :} Les épreuves de « Résumé de documents » et « Note pour le ministre » testent la capacité à extraire l'essentiel de rapports diplomatiques, de discours, de traités.
\item \textbf{Fonction publique (tous ministères) :} Rédiger un \emph{compte-rendu de réunion}, une \emph{note de service}, un \emph{rapport d'activité} = résumer des échanges ou des données pour un décideur pressé.
\item \textbf{Secteur privé (banques, assurances, ONG, entreprises) :} « Executive summary » de business plans, synthèses d'appels d'offres, veille concurrentielle.
\end{itemize}
\textbf{Transférabilité directe :} Les 5 étapes apprises ici (lire, sélectionner, reformuler, rédiger, corriger) sont \emph{identiques} à celles du métier de chargé d'études, de secrétaire général, de diplomate, de journaliste, de chef de projet. Maîtriser le résumé au Bac, c'est s'offrir un atout durable pour l'insertion professionnelle.
\end{savaistu}

\begin{exercice}[Entraînement intégré — Résumé + Justification]
\textbf{Texte support (extrait adapté, $\sim$180 mots) :}
« La satire de Mongo Beti ne vise pas seulement le colonisateur ; elle frappe aussi, et peut-être plus cruellement, le colonisé qui accepte sa condition. Le vieux nègre, figure du "bon indigène", incarne cette trahison consentie. Sa médaille n'est pas une récompense ; c'est un collier. Le rire de l'auteur n'est pas moqueur, il est lucide : il dénonce le mécanisme par lequel le pouvoir colonial transforme la victime en gardien de sa propre prison. La cérémonie du 14 juillet, avec ses drapeaux, ses discours, sa distribution de médailles, est une mise en scène où le dominé applaudit son dominateur. La lucidité de l'écrivain consiste à refuser le pathos pour préférer l'ironie froide, qui seul permet de nommer l'innommable : la complicité objective de l'opprimé à son oppression. »

\textbf{Consignes :}
\begin{enumerate}
\item Rédigez un résumé de ce texte en \textbf{45 mots ±10\% (40--50 mots)}. Appliquez la méthode en 5 étapes. Chronométrez-vous (25 min max).
\item Sur une feuille annexe, rédigez votre \textbf{justification de démarche} (modèle ci-dessus) en 10 lignes max.
\item Auto-évaluez-vous sur la grille (4 critères). Notez votre score /20.
\item Échangez avec un pair : évaluez son résumé sur la grille ; comparez vos justifications.
\end{enumerate}
\end{exercice}

\begin{corrige}[Exercice n -- Proposition de résumé et justification]
\textbf{Résumé (48 mots) :}
« Mongo Beti satirise le "bon indigène" qui, décoré le 14 juillet, célèbre sa propre aliénation. La médaille, loin d'honorer, figure le collier de la soumission consentie. Par une ironie glacée, l'auteur dénonce le mécanisme colonial qui fait du dominé le complice objectif de son oppression, refusant tout pathos pour nommer l'innommable. »

\textbf{Justification (schéma) :}
\begin{itemize}
\item \textbf{Thèse :} Complicité de l'opprimé via la médaille coloniale (ironie, pas pathos).
\item \textbf{Sélection :} 3 idées forces : (1) figure du "bon indigène", (2) médaille = collier/soumission, (3) ironie froide vs pathos. Écarté : détails cérémonie (drapeaux, discours), qualificatifs « lucide », « innommable » (reformulés).
\item \textbf{Reformulation :} « figure du "bon indigène" » gardé (terme auteur) ; « célébrer sa propre aliénation » pour « applaudit son dominateur » ; « figure le collier » pour « n'est pas une récompense ; c'est un collier » ; « complice objectif » repris (concept clé).
\item \textbf{Longueur :} Brouillon 56 mots $\rightarrow$ compression : fusion phrases 1--2, « refusant tout pathos pour nommer l'innommable » condensé en finale.
\end{itemize}
\textbf{Auto-évaluation visée :} Fidélité 5/5, Reformulation 5/5, Cohérence 5/5, Langue 5/5, Longueur OK $\rightarrow$ 20/20.
\end{corrige}

Cette section achève le cœur méthodologique de la leçon. La suivante — \emph{Contexte de l'œuvre, bilan et intégration} — replacera \emph{Le Vieux nègre et la médaille} dans l'histoire littéraire et idéologique de Mongo Beti, et vérifiera l'acquisition de l'ensemble des savoirs de l'unité par des activités de synthèse.

\coursec{Contexte de l'œuvre, bilan et intégration}

Dans la section précédente, vous avez appris à produire et à améliorer un résumé : sélectionner les idées essentielles, les reformuler, les enchaîner avec fluidité. Mais un résumé — comme toute lecture — gagne à être éclairé par le \cle{contexte} de l'œuvre. Comprendre \emph{quand}, \emph{où} et \emph{pourquoi} Ferdinand Oyono a écrit \emph{Le Vieux nègre et la médaille} permet de saisir la portée exacte du récit : Meka n'est pas seulement un vieil homme ridiculisé, il est le symbole d'une Afrique trompée par la colonisation. Cette dernière section fait le bilan de l'étude de l'œuvre et prépare l'épreuve de contraction au Baccalauréat technique.

\courssub{Le contexte historique : le Cameroun sous domination française}

\begin{definition}[Contexte historique d'une œuvre]
Le \cle{contexte historique} d'une œuvre est l'ensemble des événements, des institutions et des mentalités de l'époque où elle a été écrite et où se déroule son action. Le connaître aide à comprendre les personnages, les conflits et le message de l'auteur.
\end{definition}

Pour comprendre \emph{Le Vieux nègre et la médaille}, il faut remonter à l'histoire du Cameroun. Après la défaite de l'Allemagne en 1916 pendant la Première Guerre mondiale, le Cameroun, ancienne colonie allemande, est partagé entre la France et la Grande-Bretagne. La plus grande partie passe sous \cle{mandat français}, confié par la Société des Nations, puis sous \emph{tutelle} de l'ONU après 1945. Le pays accède à l'indépendance le 1\ier{} janvier 1960.

C'est dans ce Cameroun administré par la France que se situe le roman, publié en \textbf{1956}, quatre ans avant l'indépendance. Deux idéologies coloniales y sont dénoncées :

\begin{itemize}
\item L'\cle{assimilation} : politique qui prétend faire des colonisés des « Français à part entière », à condition qu'ils adoptent la langue, la religion, les habits et les manières des colonisateurs. Meka en est le produit : il a donné ses terres aux missionnaires, envoyé ses deux fils mourir à la guerre « pour la France », et il croit sincèrement à l'amitié des Blancs.
\item La « \cle{mission civilisatrice} » : discours qui justifie la colonisation en prétendant apporter la « civilisation », le christianisme et le progrès aux peuples africains jugés « arriérés ». Oyono montre que cette prétendue générosité cache en réalité le mépris et l'exploitation.
\end{itemize}

La médaille offerte à Meka résume cette hypocrisie : on récompense sa soumission, non sa dignité. Le lendemain de la fête, humilié et battu par les gendarmes, Meka découvre que la médaille ne fait de lui ni un ami ni un égal des Blancs.

\begin{popfigure}
\begin{center}
\begin{tikzpicture}[x=1.55cm, font=\small]
\draw[thick, popline, -{Stealth}] (0,0) -- (7.4,0);
\foreach \x/\d in {0/1916, 1.5/1939, 3/1945, 4.5/1956, 6/1960}{
  \draw[thick, popdark] (\x,-0.12) -- (\x,0.12);
  \node[below, popdark] at (\x,-0.12) {\d};
}
\node[above, align=center, text width=2.6cm, popblue] at (0,0.35) {Début du mandat\\français au Cameroun};
\draw[popblue, thin] (0,0.12) -- (0,0.3);
\node[above, align=center, text width=2.4cm, poppurple] at (1.5,0.35) {Seconde Guerre\\mondiale : soldats\\camerounais engagés};
\draw[poppurple, thin] (1.5,0.12) -- (1.5,0.3);
\node[above, align=center, text width=2.2cm, popgreen] at (3,0.35) {Tutelle de l'ONU\\(après-guerre)};
\draw[popgreen, thin] (3,0.12) -- (3,0.3);
\fill[poporange] (4.5,0) circle (0.09);
\node[above, align=center, text width=2.6cm, poporange] at (4.5,0.35) {\textbf{1956 : publication du roman}\\\emph{Le Vieux nègre et la médaille}};
\draw[poporange, thin] (4.5,0.12) -- (4.5,0.3);
\node[above, align=center, text width=2.6cm, poppink] at (6,0.35) {1\ier{} janvier 1960 :\\indépendance du Cameroun};
\draw[poppink, thin] (6,0.12) -- (6,0.3);
\end{tikzpicture}
\end{center}
\legende{Frise chronologique : du mandat français (1916) à l'indépendance du Cameroun (1960). Le roman paraît en 1956, en pleine période de contestation coloniale.}
\end{popfigure}

\begin{popfigure}
\begin{center}
\begin{tikzpicture}[font=\small, scale=1.05]
% Contour très simplifié du Cameroun
\draw[thick, popdark, fill=popcream] (0,0) -- (1.6,-0.2) -- (2.6,0.4) -- (3.1,1.6) -- (2.7,2.8) -- (2.2,3.9) -- (1.4,4.6) -- (0.7,3.9) -- (-0.2,3.1) -- (-0.5,1.9) -- (-0.3,0.8) -- cycle;
% Zone forestière du Sud (cadre du roman)
\fill[popgreenL, opacity=0.7] (-0.3,0.8) -- (0,0) -- (1.6,-0.2) -- (2.2,0.7) -- (1.2,1.3) -- (0,1.3) -- cycle;
\draw[thick, popgreen] (-0.3,0.8) -- (0,0) -- (1.6,-0.2) -- (2.2,0.7) -- (1.2,1.3) -- (0,1.3) -- cycle;
% Villes
\fill[popink] (0.5,0.35) circle (0.07); \node[right, popink] at (0.55,0.35) {Douala};
\fill[popink] (1.1,0.85) circle (0.07); \node[above right, popink] at (1.05,0.9) {Yaoundé};
\fill[poporange] (1.7,0.55) circle (0.09); \node[right, poporange] at (1.75,0.55) {\textbf{Doumé} (région de l'action)};
% Nord
\node[popmuted, align=center] at (1.2,3.4) {Nord\\(régions sèches)};
% Légende
\node[below, align=left, text width=7.5cm] at (1.3,-0.5) {\textcolor{popgreen}{\rule{0.4cm}{0.15cm}} Zone forestière du Sud : villages, missions et postes administratifs du roman};
% Nord arrow
\draw[-{Stealth}, thick, popdark] (3.6,3.6) -- (3.6,4.4); \node[right, popdark] at (3.6,4.0) {N};
\end{tikzpicture}
\end{center}
\legende{Croquis schématique du Cameroun sous tutelle française. L'action du roman se déroule dans le Sud forestier (région de Doumé, village de Meka), entre villages africains, mission catholique et poste de commandement. Échelle indicative, contours simplifiés.}
\end{popfigure}

\courssub{Le contexte littéraire : la littérature africaine francophone des années 1950}

\emph{Le Vieux nègre et la médaille} s'inscrit dans un moment capital : la naissance de la littérature africaine francophone moderne. Dans les années 1950, des écrivains africains prennent la plume pour raconter la colonisation \emph{du point de vue des colonisés}, et non plus des colons.

\begin{definition}[La Négritude]
La \cle{Négritude} est un mouvement littéraire et culturel né dans les années 1930-1950, fondé par Aimé Césaire (Martinique), Léopold Sédar Senghor (Sénégal) et Léon-Gontran Damas (Guyane). Il revendique la fierté de l'identité noire, la valeur des cultures africaines et refuse l'assimilation qui nie ces cultures.
\end{definition}

À côté de la Négritude (surtout poétique), se développe le \cle{roman de la contestation coloniale} : des romanciers comme Mongo Beti (\emph{Le Pauvre Christ de Bomba}, 1956), Camara Laye (\emph{L'Enfant noir}, 1953) et Ferdinand Oyono dénoncent, par le récit, les mensonges de la « mission civilisatrice ». Oyono se distingue par son arme favorite : l'\cle{ironie} et la satire, qui ridiculisent l'administration coloniale et les missionnaires sans jamais avoir besoin de les accuser directement.

\begin{savaistu}[Un étudiant devenu classique]
Ferdinand Oyono écrit \emph{Une vie de boy} et \emph{Le Vieux nègre et la médaille} alors qu'il est encore étudiant en France, dans les années 1950. Ces deux romans deviennent immédiatement des classiques de la littérature africaine. Plus tard, Oyono sera diplomate et ministre au Cameroun : l'écrivain contestataire est devenu un serviteur de l'État indépendant qu'il appelait de ses vœux.
\end{savaistu}

\textbf{Une vieille habitude : comparer avec \emph{Une vie de boy}.} Au Baccalauréat, il est utile de rapprocher les deux romans d'Oyono. \emph{Une vie de boy} (1956) raconte Toundi, un jeune domestique chez les Blancs, qui croit d'abord à leur supériorité avant de découvrir leur cruauté et d'en mourir. Comparaison :

\begin{center}
\begin{tabularx}{\linewidth}{|l|X|X|}
\hline
\rowcolor{popblueL} \textbf{Critère} & \textbf{\emph{Une vie de boy}} & \textbf{\emph{Le Vieux nègre et la médaille}} \\
\hline
Héros & Toundi, jeune domestique & Meka, vieux paysan \\
\hline
Illusion initiale & Admiration pour les Blancs & Fierté de la médaille et de l'« amitié » des Blancs \\
\hline
Désillusion & Découvre la cruauté, meurt & Humilié après la fête, perd ses illusions \\
\hline
Arme de l'auteur & Ironie, regard naïf & Satire, grotesque, tonalité polémique \\
\hline
Message & La colonisation détruit celui qui s'y soumet & La colonisation trompe puis brise l'Africain fidèle \\
\hline
\end{tabularx}
\end{center}

Dans les deux cas, le schéma est identique : \textbf{illusion $\rightarrow$ épreuve $\rightarrow$ éveil}. C'est une « vieille habitude » d'Oyono : montrer un Africain naïf que la réalité coloniale finit par réveiller.

\courssub{Bilan de l'étude de l'œuvre : Meka, l'Africain trompé puis éveillé}

Dressons le portrait de \cle{Meka}. Au début du roman, c'est un vieillard soumis et crédule : il a tout donné à la colonisation (ses terres, ses fils, sa foi) et il attend en retour une reconnaissance. La médaille est pour lui le sommet de sa vie. Mais la fête tourne au ridicule : ivre, il s'égare, est arrêté et frappé par les gendarmes — les mêmes hommes qui l'applaudissaient la veille. Le symbole est limpide :

\begin{exemplebox}[La pluie, signe de l'échec de l'illusion]
La \cle{pluie} qui interrompt brutalement la fête de la médaille symbolise l'échec de l'illusion coloniale : la cérémonie grandiose préparée par les Blancs se dégonfle comme une baudruche, exactement comme l'amitié qu'ils prétendaient offrir à Meka. La nature elle-même semble dénoncer la mascarade. Le lendemain, l'humiliation de Meka confirme ce que la pluie annonçait : la médaille n'était qu'un décor.
\end{exemplebox}

Meka est donc le \textbf{symbole de l'Africain trompé puis éveillé} : trompé par les promesses de l'assimilation, éveillé par l'expérience du mépris. À la fin du roman, il a perdu ses illusions mais gagné sa lucidité — et la solidarité de son village, qui rit avec lui des Blancs. C'est toute la tonalité \emph{satirique} et \emph{polémique} étudiée dans cette leçon : Oyono fait rire pour mieux faire réfléchir, et le rire devient une arme de contestation.

\begin{attention}[Erreur fréquente au Bac]
Ne plaquez \textbf{jamais} des connaissances historiques générales sur le texte sans les relier à lui. Écrire « la colonisation était une domination violente » ne vaut rien si vous ne montrez pas \emph{où} et \emph{comment} le roman le dit (la médaille, la fête, la pluie, l'arrestation de Meka). Règle d'or : \textbf{une idée générale = une preuve précise tirée du texte}.
\end{attention}

\courssub{Synthèse des acquis : de la lecture à la production du résumé}

Cette leçon vous a mené sur un chemin complet, celui-là même que demande l'épreuve du Baccalauréat technique :

\begin{enumerate}
\item \textbf{Lire et comprendre} : analyser un extrait du \emph{Vieux nègre et la médaille}, dégager le thème, la thèse et la tonalité (satirique, polémique).
\item \textbf{Identifier les idées essentielles} : distinguer l'essentiel de l'accessoire, supprimer exemples, répétitions et détails.
\item \textbf{Reformuler} : dire les idées de l'auteur avec vos propres mots, sans copier-coller ni ajouter votre avis.
\item \textbf{Rédiger et améliorer} : rédiger un texte bref, lié, correct, puis le relire pour corriger la langue (accords, connecteurs, énoncés constatifs).
\item \textbf{Situer l'œuvre} : relier le texte à son contexte historique et littéraire pour en mesurer la portée.
\end{enumerate}

\begin{methode}[Gérer son temps à l'épreuve de contraction (50 minutes)]
\begin{enumerate}
\item \textbf{10 min — Lecture active} : lisez le texte deux fois, soulignez les idées principales, repérez la thèse de l'auteur.
\item \textbf{15 min — Sélection et plan} : classez les idées essentielles, éliminez exemples et répétitions, ordonnez en 3 ou 4 étapes logiques.
\item \textbf{20 min — Rédaction} : reformulez avec vos mots, enchaînez avec des connecteurs (d'abord, ensuite, enfin ; cependant, c'est pourquoi), respectez le nombre de mots imposé.
\item \textbf{5 min — Relecture} : corrigez orthographe et accords, vérifiez qu'aucune phrase n'est recopiée du texte, comptez les mots.
\end{enumerate}
\end{methode}

\begin{exoresolu}[Appliquer le contexte dans un résumé]
\textbf{Énoncé.} Voici une phrase de résumé : « Meka reçoit une médaille. » Améliorez-la en y intégrant le sens de l'œuvre, sans dépasser 25 mots.
\tcblower
\textbf{Solution.} « En récompensant Meka d'une médaille dérisoire, les colonisateurs masquent leur mépris ; humilié dès le lendemain de la fête, le vieil Africain perd ses illusions. » (27 mots — à resserrer : « Récompensé d'une médaille dérisoire puis humilié dès le lendemain, Meka découvre le mépris caché derrière l'amitié coloniale. » = 21 mots.) La phrase relie l'événement (la médaille) à sa signification (l'illusion coloniale), sans hors-texte gratuit.
\end{exoresolu}

\courssub{Compte-rendu et remédiation}

La dernière séance est consacrée à la \textbf{correction collective} des résumés produits et à la \cle{remédiation}. Démarche en classe :

\begin{enumerate}
\item \textbf{Lecture de deux ou trois productions} d'élèves (anonymisées) : la classe repère les qualités (fidélité, brièveté, enchaînement) et les défauts (recopie du texte, avis personnel, hors-sujet historique).
\item \textbf{Grille d'autoévaluation} : chaque élève coche — ai-je respecté le nombre de mots ? reformulé sans recopier ? gardé toutes les idées essentielles ? évité mon opinion personnelle ?
\item \textbf{Réécriture ciblée} : chacun reprend uniquement le point qui a échoué (reformulation, connecteurs, ou respect de la limite de mots).
\end{enumerate}

\begin{center}
\begin{tabularx}{\linewidth}{|l|X|X|}
\hline
\rowcolor{popblueL} \textbf{Difficulté observée} & \textbf{Cause fréquente} & \textbf{Remédiation} \\
\hline
Phrases recopiées du texte & Peur de reformuler & Remplacer chaque verbe et chaque connecteur par un synonyme \\
\hline
Avis personnel inséré & Confusion résumé / commentaire & Bannir « je pense que » ; ne garder que la pensée de l'auteur \\
\hline
Idées essentielles oubliées & Lecture trop rapide & Souligner une idée par paragraphe avant de rédiger \\
\hline
Hors-texte historique & Connaissances plaquées & Relier chaque fait historique à un passage précis du texte \\
\hline
Dépassement du nombre de mots & Détails et exemples conservés & Supprimer exemples, répétitions, énumérations \\
\hline
\end{tabularx}
\end{center}

\begin{exercice}[Pour vérifier vos acquis]
\begin{enumerate}
\item Rédigez en 5 lignes un résumé du chapitre de la fête de la médaille, en utilisant au moins trois connecteurs logiques.
\item Expliquez en quelques phrases ce que la pluie symbolise dans le roman, en citant un passage précis.
\item Comparez Meka et Toundi (\emph{Une vie de boy}) : quelle « vieille habitude » narrative Oyono utilise-t-il dans les deux romans ?
\end{enumerate}
\end{exercice}

\begin{corrige}[Exercice : éléments de réponse]
\begin{enumerate}
\item Exemple attendu : « D'abord, Meka, fier de sa médaille, se rend à la fête organisée par les Blancs. Ensuite, enivré, il s'égare dans la nuit. Cependant, la pluie interrompt la cérémonie. Enfin, arrêté et frappé par les gendarmes, il comprend que l'amitié des colonisateurs n'était qu'une illusion. »
\item La pluie symbolise l'échec de l'illusion coloniale : elle gâche la fête comme le réel gâche les belles promesses faites à Meka (preuve : la cérémonie interrompue, les invités dispersés).
\item Dans les deux romans, Oyono suit le schéma illusion $\rightarrow$ épreuve $\rightarrow$ éveil : un Africain naïf admire les Blancs, subit leur mépris, et perd ses illusions.
\end{enumerate}
\end{corrige}

\begin{aretenir}[L'essentiel de la leçon]
\begin{itemize}
\item \emph{Le Vieux nègre et la médaille} (1956) dénonce l'assimilation et la « mission civilisatrice » dans le Cameroun sous tutelle française.
\item Meka incarne l'Africain trompé puis éveillé ; la pluie symbolise l'effondrement de l'illusion coloniale.
\item L'œuvre appartient au roman de la contestation coloniale, contemporain de la Négritude (Césaire, Senghor, Damas).
\item Le résumé exige : lecture active, sélection des idées essentielles, reformulation, rédaction brève et relique, sans avis personnel ni recopie.
\item Au Bac : 10 min de lecture, 15 min de sélection, 20 min de rédaction, 5 min de relecture.
\end{itemize}
\end{aretenir}

\newpage

\coursec{Activité d'intégration}

\coursec{Leçon 22 : Rédiger un résumé de texte}

\courssub{Objectifs de l'activité}

Cette activité d'intégration clôt l'unité consacrée à la contraction de texte. Elle place l'élève en situation d'examen, proche de l'épreuve de Français du Probatoire et du Baccalauréat technique, et vise les objectifs suivants :

\begin{itemize}
\item \cle{Lire activement} un extrait du \emph{Vieux nègre et la médaille} de Ferdinand Oyono et repérer ses idées essentielles ;
\item \cle{Hiérarchiser} les idées dans un tableau (idées principales, idées secondaires, exemples à éliminer) ;
\item \cle{Reformuler} fidèlement les idées de l'auteur avec ses propres mots, sans copier les phrases du texte ;
\item \cle{Rédiger} un résumé de 90 à 100 mots, cohérent, correct sur le plan de la langue, et signalant la tonalité satirique du passage ;
\item \cle{Évaluer} la production d'un camarade à l'aide d'une grille critériée et \cle{améliorer} sa propre copie en justifiant ses choix de reformulation.
\end{itemize}

\begin{aretenir}[Savoir-faire visés]
Appliquer les notions de l'unité à une situation concrète ; rédiger et justifier une démarche, une réponse ou une conclusion. Ces deux savoir-faire officiels sont ici pleinement mobilisés : la démarche de contraction est appliquée à un texte littéraire camerounais, puis chaque élève justifie par écrit deux de ses choix de reformulation.
\end{aretenir}

\courssub{Situation}

\textbf{Contexte.} La municipalité de Sangmélima prépare la cérémonie du 20 mai. Le maire doit prononcer un discours sur le thème : « Honorer ceux qui ont servi ». Le journal du lycée technique de la ville, \emph{Le Compas}, consacre son prochain numéro à la littérature camerounaise et à la dénonciation des honneurs trompeurs. La rédaction du journal, dont tu fais partie, a choisi de publier un \textbf{résumé} de la scène de la remise de la médaille dans \emph{Le Vieux nègre et la médaille} de Ferdinand Oyono, afin de montrer aux lecteurs comment l'écrivain camerounais dénonce, avec humour, l'hypocrisie coloniale. Mais le format du journal impose une contrainte stricte : le résumé ne doit pas dépasser \textbf{100 mots}.

\textbf{Extrait à résumer (environ 370 mots).} L'enseignant distribue l'extrait suivant, adapté de la scène où Meka reçoit sa médaille :

\begin{quote}
Le grand jour arriva. Meka, raide dans son costume neuf qui le gênait au cou, attendait au milieu de la foule des notables. Le commandant, casque blanc sur la tête, s'avança vers lui avec un sourire large comme la place du marché. « Mon brave Meka, dit-il d'une voix forte pour que tout le monde entende, la France n'oublie pas ses fidèles serviteurs. Tu as donné tes deux fils pour la mère patrie, tu as donné tes terres pour la mission. Reçois cette médaille, signe de notre reconnaissance éternelle. »

Meka sentit le métal froid sur sa poitrine. La foule applaudit. Les Blancs hochaient la tête avec satisfaction, comme s'ils venaient d'accomplir une grande œuvre. Pourtant, au fond de lui, le vieil homme ne ressentait aucune joie. Il pensait à ses fils morts loin d'ici, dans une guerre qui n'était pas la leur. Il pensait à ses terres qu'on lui avait prises pour y construire la mission, et sur lesquelles il ne pouvait plus même pas marcher librement.

« Dis quelque chose, vieux ! » murmura le commandant en lui serrant le bras.

Meka ouvre la bouche, mais aucun mot ne sortit. Que pouvait-il dire ? Remercier pour cette médaille qui ne rendrait ni ses fils ni ses terres ? Le commandant, gêné, déclara : « L'émotion l'étouffe, voyez-vous ! Ces braves gens ne savent pas exprimer leur gratitude, mais leur cœur parle pour eux ! » Et la foule d'applaudir plus fort.

Le soir, assis devant sa case, Meka regardait la médaille qui brillait dans la pénombre. Ses voisins venaient l'admirer. « Tu es un homme important maintenant », disait l'un. Meka hochait la tête sans répondre. Une médaille contre deux fils et des terres : voilà le compte que faisait la France. Il se demanda si cette rondelle de métal pourrait nourrir sa vieillesse, si elle pourrait consoler sa femme qui pleurait encore ses enfants chaque nuit. Le métal était lourd, mais le vide dans son cœur l'était davantage.
\end{quote}

\textbf{Tâche.} Tu rédigeras, pour \emph{Le Compas}, un résumé de cet extrait en \textbf{90 à 100 mots}, fidèle aux idées de l'auteur, entièrement reformulé, cohérent et correct, en signalant la \cle{tonalité satirique} du passage. Ta production sera ensuite évaluée par un camarade selon la grille officielle, puis améliorée.

\courssub{Consignes}

Le travail se déroule en cinq étapes, d'abord en groupes de quatre, puis individuellement.

\begin{enumerate}
\item \textbf{Lecture active (groupe).} Lis l'extrait deux fois. À la deuxième lecture, surligne les idées essentielles et barre au crayon les détails secondaires (descriptions, répétitions, exemples illustratifs).
\item \textbf{Tableau des idées (groupe).} Construis collectivement un tableau à trois colonnes : \emph{idées principales} (à conserver), \emph{idées secondaires} (à condenser), \emph{détails à éliminer}. Vérifie que les idées principales couvrent tout le texte, du début à la fin.
\item \textbf{Rédaction individuelle.} Rédige seul ton résumé de 90 à 100 mots. Reformule chaque idée avec tes propres mots (synonymes, changements de constructions, regroupements), utilise des connecteurs logiques (\emph{d'abord, cependant, en effet, finalement}) et signale la tonalité satirique par une formule comme : « l'auteur raille », « Oyono dénonce avec ironie ». Compte tes mots et indique le total à la fin.
\item \textbf{Évaluation croisée (binôme).} Échange ta copie avec celle d'un camarade d'un autre groupe. Évalue sa production à l'aide de la grille ci-dessous, note chaque critère et écris un conseil d'amélioration.
\begin{center}
\begin{tabularx}{\linewidth}{|l|X|c|}
\hline
\rowcolor{popblueL} \textbf{Critère} & \textbf{Attendu} & \textbf{Note} \\
\hline
Fidélité & Toutes les idées essentielles, aucune idée étrangère au texte & /4 \\
\hline
Reformulation & Aucune phrase copiée, vocabulaire et constructions personnels & /4 \\
\hline
Cohérence & Enchaînement logique, connecteurs, un seul paragraphe fluide & /4 \\
\hline
Langue & Orthographe, grammaire, conjugaison, ponctuation & /4 \\
\hline
Longueur & Entre 90 et 100 mots, total indiqué & /4 \\
\hline
\end{tabularx}
\end{center}
\item \textbf{Réécriture justifiée.} À partir des remarques reçues, réécris une version améliorée de ton résumé et justifie par écrit \textbf{deux} de tes choix de reformulation (par exemple : pourquoi avoir remplacé telle expression, regroupé telles idées).
\end{enumerate}

\courssub{Langue : énoncés constatifs et performatifs}

\begin{definition}[Énoncé constatif]
Un énoncé \cle{constatif} décrit un état de fait, affirme quelque chose qui peut être vrai ou faux. Il porte sur la réalité et relève de la constatation.
\textbf{Exemple dans le texte :} « Meka sentit le métal froid sur sa poitrine. » (Description d'une sensation physique, vérifiable dans la fiction).
\end{definition}

\begin{definition}[Énoncé performatif]
Un énoncé \cle{performatif} ne décrit pas, il \textbf{réalise l'action} qu'il énonce par le seul fait d'être prononcé par une personne habilitée, dans un contexte approprié (felicity conditions). Il ne peut être ni vrai ni faux, mais \emph{heureux} (réussi) ou \emph{malheureux} (raté).
\textbf{Exemples dans le texte :}
\begin{itemize}
\item « Reçois cette médaille, signe de notre reconnaissance éternelle. » (Le commandant \emph{effectue} la remise de la médaille en le disant : acte déclaratif).
\item « L'émotion l'étouffe, voyez-vous ! » (Constatif en surface, mais performatif socialement : le commandant \emph{impose} l'interprétation du silence de Meka, il \emph{fabrique} la gratitude).
\end{itemize}
\end{definition}

\begin{attention}[Distinction cruciale pour l'examen]
Ne pas confondre \emph{phrase performative canonique} (verbe performatif à la 1ère personne du présent indicatif : « Je te décerne... », « Je promets... ») et \cle{énoncé performatif} au sens large de la théorie des actes de langage (Austin, Searle) qui inclut les injonctions (impératif), les déclarations institutionnelles, les promesses, les menaces. Au Probatoire/Bac, on attend la capacité à identifier l'acte de langage accompli : \emph{dire, c'est faire}.
\end{attention}

\begin{exemplebox}[Application guidée]
Repère dans l'extrait trois énoncés performatifs prononcés par le commandant et nomme l'acte de langage accompli pour chacun.
\begin{enumerate}
\item « La France n'oublie pas ses fidèles serviteurs. » $\rightarrow$ Acte d'\textbf{engagement} / de promesse implicite (engage l'institution).
\item « Reçois cette médaille... » $\rightarrow$ Acte de \textbf{déclaration} institutionnelle (change le statut de Meka).
\item « L'émotion l'étouffe... » $\rightarrow$ Acte d'\textbf{imposition de sens} / de définition de la réalité (pouvoir discursif du colonisateur).
\end{enumerate}
\end{exemplebox}

\courssub{Stylistique : les tonalités satirique et polémique}

\begin{definition}[Tonalité satirique]
La \cle{tonalité satirique} vise à \textbf{rire de} quelque chose ou quelqu'un pour le corriger ou le dénoncer. Elle use de l'ironie (dire le contraire de ce qu'on pense), de l'exagération, de la caricature, du décalage entre le discours et la réalité. L'auteur prend de la distance, souvent avec une once d'amertume.
\textbf{Dans le texte :} Le discours du commandant (« reconnaissance éternelle », « mère patrie ») est violemment décalé par la réalité intérieure de Meka (douleur, perte, vide). La foule applaudit « comme s'ils venaient d'accomplir une grande œuvre » : l'ironie narrative souligne l'absurdité et l'hypocrisie de la cérémonie.
\end{definition}

\begin{definition}[Tonalité polémique]
La \cle{tonalité polémique} est plus directe et agressive : elle \textbf{attaque} frontalement une idée, une institution, un adversaire. Elle use de l'argumentation virulente, de l'invective, de la réfutation systématique, sans nécessairement passer par le rire.
\textbf{Lien avec l'œuvre :} \emph{Le Vieux nègre et la médaille} est globalement une œuvre polémique (contre le système colonial), mais cette scène relève principalement du \emph{satirique} : la dénonciation passe par le rire amer, le grotesque de la situation, le contraste saisissant entre le faste apparent et la misère réelle.
\end{definition}

\begin{aretenir}[Repérer et signaler la tonalité au résumé]
Signaler la tonalité dans un résumé, c'est écrire explicitement : « Oyono \emph{raille} / \emph{dénonce avec ironie} / \emph{tourne en dérision} les honneurs trompeurs de la colonisation. » Ce n'est pas donner son avis personnel, c'est rendre compte fidèlement de l'intention de l'auteur perçue dans le texte.
\end{aretenir}

\courssub{Ressources à mobiliser}

\begin{aretenir}[La méthode du résumé en cinq étapes]
\begin{enumerate}
\item \textbf{Lire} deux fois le texte : lecture globale puis lecture analytique avec surlignage.
\item \textbf{Dégager} le thème et la thèse de l'auteur ; repérer la tonalité (ici satirique) et les actes de langage performatifs.
\item \textbf{Hiérarchiser} : distinguer idées principales, idées secondaires et détails à supprimer (exemples, répétitions, descriptions pittoresques).
\item \textbf{Reformuler} : synonymes, changement de construction (nominalisation, passage actif/passif, subordination), regroupement d'idées voisines ; \textbf{jamais} de phrase copiée.
\item \textbf{Rédiger et vérifier} : un paragraphe unique, connecteurs logiques, langue correcte, respect du nombre de mots imposé (ici 90 à 100, soit environ le quart du texte source).
\end{enumerate}
\end{aretenir}

\begin{attention}[Pièges fréquents à éviter]
\begin{itemize}
\item \textbf{Copier-coller} des phrases du texte : c'est la faute la plus sanctionnée. Reformule toujours.
\item Ajouter son \textbf{avis personnel} ou des idées absentes du texte : le résumé est fidèle, pas critique.
\item Oublier de \textbf{compter les mots} : hors des limites, le critère « longueur » est perdu.
\item Négliger la \textbf{tonalité} : ici, il faut montrer qu'Oyono se moque des honneurs coloniaux.
\item Confondre résumé et \textbf{commentaire} : pas d'introduction ni de conclusion développées, un seul paragraphe.
\end{itemize}
\end{attention}

\begin{savaistu}[Ferdinand Oyono, écrivain et homme d'État camerounais]
Ferdinand Oyono (1929-2010), né à Ngoulémakong près de Sangmélima (région du Sud), est l'un des grands romanciers camerounais de langue française. \emph{Le Vieux nègre et la médaille} (1956) dénonce, avec une ironie mordante, l'injustice de la colonisation : le héros Meka a tout sacrifié (ses fils, ses terres) et ne reçoit en échange qu'une médaille sans valeur. Oyono fut aussi diplomate (ambassadeur) et ministre de la République du Cameroun (Fonction publique, Affaires étrangères).
\end{savaistu}

\courssub{Démarche et corrigé commenté}

\begin{exoresolu}[Résumer la scène de la remise de la médaille (90 à 100 mots)]
Énoncé : à partir de l'extrait distribué, produire le tableau des idées hiérarchisées, puis un résumé de 90 à 100 mots signalant la tonalité satirique, et justifier deux choix de reformulation.
\tcblower
\textbf{Étape 1 — Lecture active.} À la deuxième lecture, on surligne : le discours du commandant et la remise de la médaille (acte performatif) ; le contraste entre les apparences (foule qui applaudit, satisfaction des Blancs) et le ressenti de Meka (aucune joie, douleur) ; les sacrifices (deux fils morts, terres confisquées) ; le silence de Meka récupéré et détourné par le commandant ; la méditation du soir (la médaille ne nourrit ni ne console). On barre : « casque blanc », « raide dans son costume neuf », « sourire large comme la place du marché », propos des voisins le soir (détails pittoresques/illustratifs).

\textbf{Étape 2 — Tableau des idées hiérarchisées.}
\begin{center}
\begin{tabularx}{\linewidth}{|X|X|X|}
\hline
\rowcolor{popblueL} \textbf{Idées principales (à garder)} & \textbf{Idées secondaires (à condenser)} & \textbf{Détails à éliminer} \\
\hline
Meka reçoit une médaille pour ses sacrifices (fils, terres) & Le discours emphatique du commandant (performatif) & Costume neuf, casque blanc \\
\hline
Il n'éprouve aucune joie : contraste apparence/réalité & Réactions de la foule et des Blancs & « sourire large comme la place du marché » \\
\hline
Son silence est récupéré/détourné par le commandant & Le commandant gêné, interprétation fausse & Propos des voisins le soir \\
\hline
Le soir, il médite : la médaille ne compense rien & La douleur de sa femme (évoquée) & « rondelle de métal », pénombre \\
\hline
\end{tabularx}
\end{center}

\textbf{Étape 3 — Résumé rédigé (modèle, 93 mots).}

\emph{Lors d'une cérémonie officielle, l'administrateur colonial remet une médaille au vieux Meka pour récompenser le sacrifice de ses deux fils et de ses terres. Pourtant, derrière les discours pompeux et les applaudissements, le héros ne ressent que douleur et amertume : cet honneur ne saurait combler ses pertes. Muet face à l'injonction de remercier, il voit son silence détourné par l'officier qui le proclame reconnaissant. Le soir, Meka mesure le vide de cette récompense, incapable de nourrir ni de consoler. Oyono raille ainsi, avec une ironie amère, les honneurs trompeurs de la colonisation.} (93 mots)

\textbf{Étape 4 — Justification de deux choix de reformulation.}
\begin{enumerate}
\item « \emph{remet une médaille au vieux Meka pour récompenser le sacrifice de ses deux fils et de ses terres} » remplace tout le paragraphe du discours du commandant (« Mon brave Meka... Reçois cette médaille... ») : j'ai condensé l'acte performatif de déclaration et son contenu en une seule proposition infinitive, en supprimant le style direct et les formules rhétoriques (« mère patrie », « reconnaissance éternelle ») qui relèvent de l'ironie déjà signalée in fine.
\item « \emph{honneurs trompeurs} » reformule « une médaille qui ne rendrait ni ses fils ni ses terres / le vide de cette récompense » : la nominalisation (adjectif + nom) résume en deux mots une idée développée sur plusieurs phrases, ce qui économise le lexique tout en nommant explicitement la visée satirique de l'auteur.
\end{enumerate}

\textbf{Vérification finale} : fidélité (les quatre idées principales présentes), reformulation (aucune phrase copiée, structures nouvelles), cohérence (connecteurs \emph{pourtant, ainsi}, progression chronologique), langue (correcte), longueur (93 mots, dans la fourchette 90-100).
\end{exoresolu}

\courssub{Bilan de l'activité}

Cette activité a permis de mettre en évidence plusieurs acquis et difficultés.

\begin{itemize}
\item \textbf{Sur la méthode} : le résumé n'est pas un raccourci mécanique mais un travail en étapes — lire, hiérarchiser, reformuler, rédiger, vérifier. Le tableau des idées est l'outil décisif : il garantit la fidélité et évite les oublis ou les ajouts.
\item \textbf{Sur la reformulation} : les élèves ont constaté que condenser suppose de changer de point de vue grammatical (passage du style direct à la narration, nominalisations, regroupements syntaxiques), et non de remplacer mot à mot par des synonymes.
\item \textbf{Sur la langue (énoncés performatifs)} : l'analyse du discours du commandant a montré que la cérémonie repose sur des actes de langage (déclaration, engagement, imposition de sens) que le résumé doit rendre par des verbes d'action (\emph{remet, proclame, impose}) plutôt que par la citation.
\item \textbf{Sur la tonalité} : signaler la satire (« Oyono raille », « ironie amère ») prouve que le résumé rend compte non seulement du \emph{quoi} (les idées), mais aussi du \emph{comment} (l'intention critique de l'auteur).
\item \textbf{Sur l'évaluation} : la grille critériée (fidélité, reformulation, cohérence, langue, longueur, chacune sur 4) a rendu les élèves capables de s'auto-évaluer comme à l'examen. Le débriefing collectif a dégagé les erreurs fréquentes — phrases copiées, avis personnels, mots non comptés, tonalité oubliée — qui feront l'objet de la séance de \cle{remédiation}.
\end{itemize}

\begin{exercice}[Pour s'entraîner à la maison]
Choisis un passage de 350 à 400 mots d'un autre chapitre du \emph{Vieux nègre et la médaille} (par exemple la visite du commandant au village au début du roman). Rédige un résumé de 90 à 100 mots en suivant les cinq étapes de la méthode, en identifiant au passage un énoncé performatif et en signalant la tonalité. Auto-évalue ta production avec la grille officielle. Apporte ta production à la séance de remédiation.
\end{exercice}

\newpage

\coursec{Méthodes et exercices résolus}

\coursec{Rédiger un résumé de texte}

\courssub{Lecture active et compréhension de l'œuvre (Exposé N°1)}

\begin{methode}[Lecture annotée et repérage des informations clés]
    \textbf{Objectif :} Maîtriser le contenu du texte pour en extraire la substantifique moelle.
    \begin{enumerate}
        \item \textbf{Lecture globale (5 min) :} Lire le texte d'une traite sans s'arrêter au vocabulaire difficile. Identifier : l'auteur, le titre, le genre (roman), le thème général (colonialisme, ironie, condition du « vieux nègre »).
        \item \textbf{Lecture détaillée (15 min) :} Relire en annotant dans la marge :
              \begin{itemize}
                  \item \cle{Personnages} : Qui parle ? Qui agit ? (Meka, le prêtre, l'administrateur, Kelara, etc.)
                  \item \cle{Espaces/Temps} : Le village (Doum), la ville, la cérémonie. Chronologie : avant/après la médaille.
                  \item \cle{Actions principales} : Quête de la médaille $\rightarrow$ Cérémonie $\rightarrow$ Désillusion/Vol.
                  \item \cle{Mots-clés / Champs lexicaux} : Honneur, médaille, blanc, noir, argent, fête, humiliation.
              \end{itemize}
        \item \textbf{Structuration :} Découper le texte en \cle{séquences logiques} (ex: 1. L'annonce, 2. Les préparatifs, 3. La cérémonie, 4. Le vol, 5. Le constat). Donner un titre à chaque séquence.
        \item \textbf{Synthèse orale :} Raconter l'histoire en 3 phrases maximum (Début, Milieu, Fin).
    \end{enumerate}
    \textbf{Réflexe Bac :} Ne jamais commencer un résumé sans avoir fait ce travail de « squelette » du texte.
\end{methode}

\begin{exoresolu}[Analyse d'un extrait type (simulant l'Exposé 1)]
\textbf{Énoncé :} Lisez l'extrait suivant (inspiré du roman) et répondez aux questions.
\begin{quote}
\emph{« Meka ne dormit pas de la nuit. La médaille brillait au fond de sa malle, comme un petit soleil qu'il aurait capturé. Demain, le Commandant lui épinglerait sur la poitrine ce morceau de métal qui valait plus que toutes les récoltes de cacao de Doum. Il revoyait déjà la scène : les tambours, la foule, les blancs qui salueraient \textbf{lui}, Meka, l'ancien porteur, le fidèle serviteur. Kelara, sa femme, avait mis les bouchées doubles : poulet, vin de palme, tout le village mangerait à sa gloire. Mais une ombre traversait son rêve : et si on la lui volait, cette médaille, avant même qu'il ne l'ait eue ? »}
\end{quote}
\textbf{Questions :}
\begin{enumerate}
    \item Identifiez le personnage principal, son statut social et son objectif.
    \item Relevez deux indices montrant l'importance symbolique de la médaille pour Meka.
    \item Quel sentiment domine le dernier paragraphe ? Justifiez par un mot du texte.
    \item Proposez un titre à cette séquence narrative.
\end{enumerate}
\tcblower
\textbf{Solution détaillée :}
\begin{enumerate}
    \item \textbf{Personnage :} Meka. \textbf{Statut :} « ancien porteur », « fidèle serviteur » (colonisé, subalterne). \textbf{Objectif :} Recevoir la médaille du Commandant lors de la cérémonie du lendemain.
    \item \textbf{Indices symboliques :}
          \begin{itemize}
              \item « brillait [...] comme un petit soleil qu'il aurait capturé » (comparaison $\rightarrow$ valeur inestimable, lumière, pouvoir).
              \item « valait plus que toutes les récoltes de cacao de Doum » (hyperbole $\rightarrow$ valeur d'échange supérieure à la richesse réelle du village).
          \end{itemize}
    \item \textbf{Sentiment :} L'\cle{angoisse} / la \cle{crainte} / le \cle{doute}. \textbf{Justification :} « Mais une \textbf{ombre} traversait son rêve », « et si on la lui \textbf{volait} ». Le champ lexical de l'inquiétude (ombre, vol, peur implicite) contredit l'euphorie initiale.
    \item \textbf{Titre proposé :} « La veille de la cérémonie : entre rêve de gloire et peur du destin » ou « L'obsession de la médaille ».
\end{enumerate}
\textbf{Conclusion :} Cet extrait pose l'enjeu central : la médaille comme symbole de reconnaissance illusoire par le colonisateur, source d'aliénation et de fragilité pour le colonisé.
\end{exoresolu}

\courssub{La contraction de texte : reformuler les idées essentielles}

\begin{methode}[De la paraphrase à la reformulation synthétique]
    \textbf{Objectif :} Passer du texte original (source) à une version condensée (cible) sans trahir le sens, en changeant la forme (vocabulaire, syntaxe).
    \begin{enumerate}
        \item \textbf{Segmenter :} Découper le paragraphe en \cle{unités de sens} (une idée = un trait de soulignement ou un numéro).
        \item \textbf{Hiérarchiser :} Classer les unités : \cle{Idée principale (IP)} (indispensable) vs \cle{Détails/Exemples (D)} (supprimables ou fusionnables).
        \item \textbf{Dépouiller (Supprimer) :} Rayez : adjectifs qualificatifs non essentiels, répétitions, exemples, détails descriptifs, connecteurs de remplissage.
        \item \textbf{Reformuler (Réécrire) :}
              \begin{itemize}
                  \item \cle{Généraliser} : « poulet, vin de palme, ignames » $\rightarrow$ « un festin » / « des vivres ».
                  \item \cle{Substituer} : Remplacer les termes précis par des hyperonymes (Meka $\rightarrow$ le protagoniste / le vieux ; médaille $\rightarrow$ la décoration / l'insigne).
                  \item \cle{Changer la structure} : Phrase affirmative $\rightarrow$ nominale ; Discours direct $\rightarrow$ indirect ; Actif $\rightarrow$ Passif (si agent non important).
              \end{itemize}
        \item \textbf{Lier :} Utiliser des \cle{connecteurs logiques} (donc, car, mais, ensuite, finalement) pour assurer la cohésion.
        \item \textbf{Vérifier :} Le sens est-il conservé ? Le ton est-il neutre/objectif ? (Pas d'interprétation personnelle).
    \end{enumerate}
    \textbf{Formule magique :} \textbf{IP + IP + IP = Résumé} (Les détails meurent, les idées survivent).
\end{methode}

\begin{exoresolu}[Exercice de reformulation progressive]
\textbf{Énoncé :} Contractez le paragraphe suivant en \textbf{une seule phrase de 30 mots maximum} en gardant l'essentiel.
\begin{quote}
\emph{« Le jour J, le soleil cognait dur sur la place publique. Meka, vêtu de son pagne neuf, attendait, le cœur battant la chamade. Le Commandant arriva dans sa grosse voiture noire, entouré de soldats. Il épingla la médaille sur la poitrine de Meka, lui serra la main brièvement, puis repartit aussi vite qu'il était venu, laissant Meka seul avec sa breloque sous les acclamations forcées de la foule. »}
\end{quote}
\tcblower
\textbf{Solution détaillée (étapes visibles) :}
\begin{enumerate}
    \item \textbf{Segmentation / Hiérarchisation :}
          \begin{itemize}
              \item [D] Le soleil cognait / place publique (Cadre $\rightarrow$ accessoire).
              \item [D] Meka vêtu de neuf / cœur battant (Détails psychologiques/physiques $\rightarrow$ accessoires).
              \item [IP] Le Commandant arrive / épingle la médaille / serre la main / repart vite (\textbf{Action centrale}).
              \item [IP] Meka reste seul avec la médaille (\textbf{Conséquence/État final}).
              \item [D] Acclamations forcées de la foule (Détail ironique $\rightarrow$ peut être gardé si place, sinon coupé).
          \end{itemize}
    \item \textbf{Reformulation / Fusion :}
          \begin{itemize}
              \item « Le Commandant arrive, épingle la médaille, serre la main, repart » $\rightarrow$ « Le Commandant remet brièvement la médaille à Meka puis s'en va. »
              \item « Meka reste seul avec sa breloque » $\rightarrow$ « Meka demeure seul avec sa décoration. »
          \end{itemize}
    \item \textbf{Liaison :} « Le Commandant remet brièvement la médaille à Meka puis s'en va, \textbf{laissant} ce dernier seul avec sa décoration. »
    \item \textbf{Comptage :} 19 mots. \checkmark Contrainte respectée.
\end{enumerate}
\textbf{Réponse finale :} \emph{« Le Commandant remet brièvement la médaille à Meka puis s'en va, laissant ce dernier seul avec sa décoration. »}
\end{exoresolu}

\courssub{Langue : Énoncés constatifs et performatifs}

\begin{methode}[Identifier et distinguer Constatif vs Performatif]
    \textbf{Objectif :} Maîtriser la distinction austinienne pour l'analyse de texte et l'expression précise.
    \begin{enumerate}
        \item \textbf{Analyser le verbe (critère principal) :}
              \begin{itemize}
                  \item \cle{Performatif} : Verbe d'action langagière à la \textbf{1ère personne du singulier}, \textbf{présent de l'indicatif} (souvent), \textbf{actif} (ex: \emph{je promets, je jure, je déclare, je nomme, je parie, j'ordonne}). \textbf{Test :} On peut ajouter « \emph{par la présente} » ou « \emph{ici} » $\rightarrow$ \emph{« Je \textbf{vous déclare} mari et femme (par la présente) »}. L'énoncé \textbf{réalise} l'action qu'il énonce.
                  \item \cle{Constatif} : Verbe d'état, d'opinion, de description (tous les autres temps, toutes les autres personnes). \textbf{Test :} On peut demander « \emph{Est-ce vrai ?} » / « \emph{Est-ce faux ?} » $\rightarrow$ \emph{« Il pleut » (Vrai/Faux)}. L'énoncé \textbf{constate} une réalité.
              \end{itemize}
        \item \textbf{Analyser la force illocutoire (contexte) :} Un même verbe peut être performatif ou non selon le contexte. \emph{« Je fume »} (constatif, description) vs \emph{« Je vous promets d'arrêter de fumer »} (performatif, engagement).
        \item \textbf{Repérer les performatifs explicites vs implicites :}
              \begin{itemize}
                  \item Explicite : Verbe performatif présent (\emph{Je vous ordonne de partir}).
                  \item Implicite (indirect) : Autre structure mais même force (\emph{Partez !} = impératif / ordre ; \emph{La porte est ouverte} = ordre implicite de fermer).
              \end{itemize}
    \end{enumerate}
    \textbf{Au Bac :} On vous demande souvent de \textbf{classer} des énoncés ou de \textbf{transformer} un constatif en performatif (ex: « Il est nommé chef » $\rightarrow$ « Je le nomme chef »).
\end{methode}

\begin{exoresolu}[Classification et transformation d'énoncés]
\textbf{Énoncé :}
\begin{enumerate}
    \item Classez les énoncés suivants en \textbf{Constatif (C)} ou \textbf{Performatif (P)}. Justifiez brièvement.
    \begin{itemize}
        \item[a)] « Je vous déclare époux et épouse. »
        \item[b)] « Le Vieux Nègre a reçu une médaille. »
        \item[c)] « Je parie que Meka va la perdre. »
        \item[d)] « Oyono dénonce le colonialisme. »
        \item[e)] « Je vous ordonne de rendre la médaille. »
    \end{itemize}
    \item Transformez l'énoncé constatif (b) en énoncé performatif explicite (imaginez le contexte : le Commandant parle).
\end{enumerate}
\tcblower
\textbf{Solution détaillée :}
\textbf{1. Classification :}
\begin{itemize}
    \item[a)] \textbf{P (Performatif).} Verbe \emph{déclare} (1ère pers. sg, présent, actif). Test : « Par la présente, je vous déclare... » $\rightarrow$ L'acte de mariage est accompli \textbf{par} la parole.
    \item[b)] \textbf{C (Constatif).} Verbe \emph{a reçu} (3ème pers. sg, passé composé). Test : « Est-ce vrai ? » $\rightarrow$ Oui/Non. Cela décrit un fait passé.
    \item[c)] \textbf{P (Performatif).} Verbe \emph{parie} (1ère pers. sg, présent, actif). Test : « Par la présente, je parie... » $\rightarrow$ L'engagement du pari est pris \textbf{maintenant}.
    \item[d)] \textbf{C (Constatif).} Verbe \emph{dénonce} (3ème pers. sg, présent). Description d'une intention/littérature. Pas d'acte juridique immédiat.
    \item[e)] \textbf{P (Performatif).} Verbe \emph{ordonne} (1ère pers. sg, présent, actif). Test : « Par la présente, j'ordonne... » $\rightarrow$ L'ordre est donné \textbf{par} l'énonciation.
\end{itemize}
\textbf{2. Transformation (b) :}
\begin{itemize}
    \item Contexte : Le Commandant remet la médaille.
    \item Constatif : « Le Vieux Nègre a reçu une médaille. »
    \item \textbf{Performatif :} \emph{« Je vous remets cette médaille. »} ou \emph{« Je vous décore de cette médaille. »}
    \item \textbf{Justification :} Passage à la 1ère personne, présent de l'indicatif, verbe d'action langagière (\emph{remets}, \emph{décore}) réalisant l'acte de décoration.
\end{itemize}
\end{exoresolu}

\courssub{Stylistique : Les tonalités satirique et polémique}

\begin{methode}[Repérer et analyser la satire et la polémique]
    \textbf{Objectif :} Démontrer comment l'auteur critique en riant (satire) ou en attaquant (polémique).
    \begin{enumerate}
        \item \textbf{Définir la cible (La \cle{visée}) :} Qui/Qu'est-ce qui est visé ? (Le colonisateur, l'administration, l'aliénation du colonisé, la médaille vide de sens).
        \item \textbf{Identifier les procédés de la \cle{Tonalité Satirique} (Rire critique) :}
              \begin{itemize}
                  \item \cle{Ironie / Antiphrase} : Dire le contraire de ce qu'on pense (ex: « \emph{ce morceau de métal qui valait plus que toutes les récoltes} »).
                  \item \cle{Hyperbole / Exagération} : Grossir le trait (ex: « \emph{petit soleil} », fête démesurée pour une breloque).
                  \item \cle{Dérision / Caricature} : Rendre ridicule (Meka qui ne dort pas, le Commandant qui « arrive et repart » comme une étoile filante).
                  \item \cle{Champ lexical dépréciatif} : Mots péjoratifs (breloque, ferraille, singe, comédiens).
              \end{itemize}
        \item \textbf{Identifier les procédés de la \cle{Tonalité Polémique} (Attaque directe/argumentative) :}
              \begin{itemize}
                  \item \cle{Accusation directe / Verbes forts} : \emph{Exploiter, mentir, voler, assassiner, bafouer}.
                  \item \cle{Argumentation} : Raisons logiques, exemples concrets (le vol de la médaille, l'hypocrisie du prêtre).
                  \item \cle{Apostrophe / Interpellation} : S'adresser à l'ennemi (ex: « \emph{Ô vous, les blancs...} »).
                  \item \cle{Champ lexical de la révolte / justice} : Liberté, dignité, droits, honte.
              \end{itemize}
        \item \textbf{Analyser l'articulation :} Chez Oyono, la \textbf{satire masque souvent la polémique}. Le rire (Meka naïf) fait mal (critique du système).
    \end{enumerate}
    \textbf{Tableau récapitulatif :}
    \begin{center}
    \begin{tabularx}{\linewidth}{|l|X|X|}\hline
    \rowcolor{popblueL} \textbf{Critère} & \textbf{Satirique} & \textbf{Polémique} \\ \hline
    \textbf{But} & Corriger par le rire, divertir en instruisant & Convaincre, combattre, détruire une thèse/ennemi \\ \hline
    \textbf{Ton} & Moqueur, léger (en surface), indirect & Grave, violent, direct, engagé \\ \hline
    \textbf{Armes} & Ironie, caricature, parodie, sous-entendu & Argumentation, invective, réquisitoire, preuves \\ \hline
    \textbf{Effet visé} & Rire jaune, prise de conscience douce & Indignation, adhésion, action, rupture \\ \hline
    \end{tabularx}
    \end{center}
\end{methode}

\begin{exoresolu}[Analyse de tonalités sur extraits courts]
\textbf{Énoncé :} Identifiez la tonalité dominante (Satirique ou Polémique) et citez \textbf{deux procédés} précis pour chaque extrait.
\begin{quote}
\textbf{Extrait A :} \emph{« Le prêtre fit un sermon magnifique. Il parla de l'humilité, de la pauvreté volontaire, du royaume des cieux. Meka, ému, pensa à sa médaille. Le prêtre ne parla pas de la médaille. Il parla du ciel. C'était plus distingué. »}
\textbf{Extrait B :} \emph{« Ils nous ont pris nos terres. Ils nous ont pris nos bras. Ils nous donnent une médaille en toc pour nous faire oublier qu'ils nous ont volé notre dignité. Rendez-nous nos terres, gardez votre ferraille ! »}
\end{quote}
\tcblower
\textbf{Solution détaillée :}
\textbf{Extrait A :} \textbf{Tonalité Satirique.}
\begin{enumerate}
    \item \textbf{Ironie / Antiphrase :} « \emph{C'était plus distingué} » $\rightarrow$ Sous-entendu : le prêtre évite le sujet qui fâche (l'argent, le matériel) par hypocrisie / distinction de classe. Le narrateur se moque de cette « distinction » qui cache la cupidité.
    \item \textbf{Contraste / Parallélisme ironique :} Opposition « \emph{royaume des cieux} » (spirituel) vs « \emph{médaille} » (matériel, terre-à-terre). Meka « \emph{ému} » pense à l'objet, pas à l'âme. Le narrateur souligne le décalage entre le discours religieux et la réalité triviale.
\end{enumerate}
\textbf{Extrait B :} \textbf{Tonalité Polémique.}
\begin{enumerate}
    \item \textbf{Accusation directe / Anaphore :} « \emph{Ils nous ont pris... Ils nous ont pris... Ils nous donnent...} » (Répétition du sujet « Ils » = colonisateurs). Énumération des vols (terres, bras, dignité) $\rightarrow$ Réquisitoire implacable.
    \item \textbf{Apostrophe / Impératif injonctif :} « \emph{Rendez-nous nos terres, gardez votre ferraille !} » $\rightarrow$ S'adresse directement à l'oppresseur. Verbes d'action (\emph{rendez, gardez}). Rejet violent du symbole (médaille = « \emph{ferraille} » = mépris).
\end{enumerate}
\end{exoresolu}

\courssub{Production et amélioration du résumé (Épreuve type Bac)}

\begin{methode}[Rédaction du résumé : La méthode en 4 temps (Épreuve officielle)]
    \textbf{Objectif :} Produire un résumé fidèle, cohérent, dans la limite de mots imposée (souvent 1/4 ou 1/3 du texte, ou nombre fixe ex: 80-100 mots).
    \begin{enumerate}
        \item \textbf{Temps 1 : Analyse du sujet \& Contraintes (5 min).} Lire le texte support. Noter : \textbf{Nombre de mots imposé} (ex: 90 mots $\pm$ 10\%). Identifier le \textbf{type de texte} (narratif, argumentatif) et le \textbf{sujet central}.
        \item \textbf{Temps 2 : Repérage \& Sélection des Idées Forces (15 min).}
              \begin{itemize}
                  \item Numéroter les paragraphes.
                  \item Souligner \textbf{UNE seule idée force par paragraphe} (sujet + prédicat essentiel).
                  \item \textbf{Éliminer} : exemples, descriptions, dialogues, répétitions, commentaires de l'auteur, détails anecdotiques.
              \end{itemize}
        \item \textbf{Temps 3 : Rédaction du Brouillon structuré (20 min).}
              \begin{itemize}
                  \item \textbf{Présenter la source :} « Dans \emph{Le Vieux Nègre et la Médaille} de Ferdinand Oyono, ... » (ou « Le texte relate... »).
                  \item \textbf{Enchaîner les idées} par des connecteurs logiques (\emph{Dans un premier temps, Ensuite, Cependant, Finalement}).
                  \item \textbf{Reformuler} (changer vocabulaire/syntaxe, discours indirect, présent de narration).
                  \item \textbf{Compter les mots} (1 trait = 1 mot). Ajuster : développer si trop court (rajouter idée oubliée), couper si trop long (supprimer adjectifs, fusionner phrases).
              \end{itemize}
        \item \textbf{Temps 4 : Copie propre \& Relecture (10 min).}
              \begin{itemize}
                  \item Recopier en soignant l'orthographe/grammaire.
                  \item \textbf{Vérification finale :} Sens fidèle ? Limite mots respectée ? Français correct ? Pas de « Je » (sauf citation) ? Pas de commentaire personnel ?
              \end{itemize}
    \end{enumerate}
    \textbf{Astuce comptage :} Comptez les mots d'une ligne type $\times$ nombre de lignes $\approx$ total. Marquez le total en bas de page.
\end{methode}

\begin{exoresolu}[Sujet type Bac : Résumé d'un texte inédit (simulation Exposé N°2)]
\textbf{Énoncé :} \emph{Vous résumerez le texte ci-après en \textbf{80 mots maximum} ($\pm$ 10\%), en en dégageant les idées essentielles.}
\begin{quote}
\emph{« La cérémonie achevée, Meka rentra chez lui, le cœur gonflé d'une fierté nouvelle. Il posa la médaille sur la table, à côté de la lampe à pétrole, et l'admira longuement. Kelara prépara le festin. Les voisins vinrent, touchèrent la médaille, félicitèrent l'heureux élu. La nuit tomba. Meka s'endormit, la main sur son trésor. Au petit matin, la table était vide. La médaille avait disparu. Kelara pleurait. Les voisins avaient fui. Seule restait la lampe, témoin muet de la supercherie. Meka comprit alors que la médaille n'était pas un honneur, mais un piège. »}
\end{quote}
\tcblower
\textbf{Solution détaillée :}
\textbf{Étape 1 : Contraintes.} 80 mots $\pm$ 10\% $\rightarrow$ \textbf{Entre 72 et 88 mots}.
\textbf{Étape 2 : Idées forces (1 par séquence) :}
\begin{enumerate}
    \item Meka fier, admire sa médaille chez lui (Début).
    \item Fête populaire, admiration des voisins (Apogée).
    \item Nuit / Sommeil (Transition).
    \item Vol de la médaille au matin / Fuite des voisins (Chute).
    \item Prise de conscience : médaille = piège / illusion (Résolution/Sens).
\end{enumerate}
\textbf{Étape 3 : Rédaction brouillon (reformulation + liaison) :}
\emph{« Après la cérémonie, Meka rentre fier et expose sa médaille. Kelara organise une fête où les villageois admirent l'insigne. La nuit venue, Meka s'endort dessus. Mais au matin, la médaille a disparu ; les invités ont fui. Seule la lampe reste. Meka réalise alors que cette décoration n'était qu'un leurre, un piège tendu par l'administration coloniale. »}
\textbf{Étape 4 : Comptage \& Ajustement.}
Comptons : \emph{Après(1) la2 cérémonie,3 Meka4 rentre5 fier6 et7 expose8 sa9 médaille.10 Kelara11 organise12 une13 fête14 où15 les16 villageois17 admirent18 l'insigne.19 La20 nuit21 venue,22 Meka23 s'endort24 dessus.25 Mais26 au27 matin,28 la29 médaille30 a31 disparu ;32 les33 invités34 ont35 fui.36 Seule37 la38 lampe39 reste.40 Meka41 réalise42 alors43 que44 cette45 décoration46 n'était47 qu'un48 leurre,49 un50 piège51 tendu52 par53 l'administration54 coloniale55.} $\rightarrow$ \textbf{55 mots}.
\textbf{Problème :} Trop court (min 72). Il faut développer légèrement (ajouter détails significatifs : « fierté nouvelle », « admiration longue », « supercherie »).
\textbf{Version finale (visant $\approx$ 80 mots) :}
\emph{« Après la cérémonie, Meka rentre chez lui, le cœur gonflé d'une fierté nouvelle. Il pose la médaille sur la table et l'admire longuement. Kelara prépare un festin ; les voisins viennent toucher l'insigne et le féliciter. La nuit tombe, Meka s'endort la main sur son trésor. Mais au petit matin, la table est vide : la médaille a disparu, les voisins ont fui. Seule la lampe témoigne de la supercherie. Meka comprend alors que cette décoration n'était pas un honneur, mais un piège cruel. »}
\textbf{Comptage final :} 84 mots. \checkmark \textbf{Dans la fourchette [72-88].}
\textbf{Conclusion :} Le résumé respecte la chronologie, les 5 idées forces, la reformulation (discours indirect, vocabulaire varié : insigne, leurre) et la contrainte quantitative.
\end{exoresolu}

\courssub{Contexte de l'œuvre, bilan et intégration}

\begin{methode}[Situer l'œuvre et synthétiser pour l'oral/écrit d'intégration]
    \textbf{Objectif :} Maîtriser les clés de lecture globales pour l'épreuve orale ou la dissertation.
    \begin{enumerate}
        \item \textbf{Contexte historique \& biographique (Fiche d'identité) :}
              \begin{itemize}
                  \item \textbf{Auteur :} Ferdinand \textbf{Oyono} (Cameroun, 1929-2010). Diplomate, haut fonctionnaire. Témoin de la colonisation puis de l'indépendance.
                  \item \textbf{Publication :} 1956 (Éditions René Julliard). \textbf{Période :} Fin de la colonisation française au Cameroun (guerre de libération UPC, répression). Pré-indépendance.
                  \item \textbf{Genre :} Roman \textbf{satirique} / \textbf{anti-colonial} / \textbf{réaliste}.
              \end{itemize}
        \item \textbf{Thèmes majeurs (Pour dissertation) :}
              \begin{itemize}
                  \item \cle{Aliénation / Mimétisme :} Meka veut ressembler au Blanc (médaille, pagne neuf, fête) $\rightarrow$ Perte d'identité.
                  \item \cle{Illusion de la reconnaissance :} La médaille = « breloque », « ferraille », « piège ». Le colonisateur ne donne que du vide.
                  \item \cle{Complicité / Trahison :} Kelara (femme), voisins (intérêt matériel), prêtre (hypocrisie), administration (manipulation).
                  \item \cle{La parole / L'oralité :} Style « conte », proverbes, dialogues en français « petit nègre » volontaire (caricature), humour.
              \end{itemize}
        \item \textbf{Structure narrative (Bilan Exposés 1 \& 2) :}
              \begin{itemize}
                  \item \textbf{Exposé 1 (Début) :} Quête de la médaille (Espoir, naïveté, préparation faste).
                  \item \textbf{Exposé 2 (Fin) :} Cérémonie $\rightarrow$ Vol $\rightarrow$ Désillusion lucide (Meka « comprend »).
              \end{itemize}
        \item \textbf{Intégration (Sujet type) :} \emph{« En quoi le rire est-il une arme chez Oyono ? »} $\rightarrow$ Plan : 1. Satire (rire de Meka/colonisés) 2. Polémique (rire jaune du colonisateur) 3. Prise de conscience (rire libérateur).
    \end{enumerate}
\end{methode}

\begin{exoresolu}[Sujet d'intégration : Dissertation courte / Question de synthèse]
\textbf{Énoncé :} \emph{« Le Vieux Nègre et la Médaille est une fable amère où le rire cache la colère. »} \textbf{Dans quelle mesure cette affirmation résume-t-elle l'esthétique et la visée de l'œuvre ?} (Réponse structurée en arguments principaux).
\tcblower
\textbf{Solution détaillée (Plan détaillé + Introduction/Conclusion types) :}
\textbf{Introduction :}
\begin{itemize}
    \item \textbf{Amorce :} Ferdinand Oyono, figure majeure de la littérature camerounaise d'expression française, publie en 1956 une charge virulente contre le système colonial.
    \item \textbf{Sujet :} Le roman relate la quête vaine de Meka, vieux porteur, pour obtenir une médaille d'honneur.
    \item \textbf{Problématique :} Comment l'humour et la satire permettent-ils de dévoiler la violence symbolique du colonialisme et d'éveiller les consciences ?
    \item \textbf{Annonce du plan :} Nous verrons que le rire est d'abord un \textbf{miroir déformant} (satire des personnages), puis une \textbf{arme politique} (polémique anti-coloniale), pour aboutir à une \textbf{prise de conscience tragique} (l'amertume finale).
\end{itemize}
\textbf{Développement :}
\textbf{I. Le rire satirique : démasquer la naïveté et l'hypocrisie (Le miroir déformant).}
\begin{itemize}
    \item \textbf{Meka, antihéros comique :} Naïveté touchante (ne dort pas, « petit soleil »), mimétisme ridicule (pagne neuf, fête pour une « breloque »). Le lecteur rit \textbf{de} lui (supériorité).
    \item \textbf{La galerie de portraits caricaturaux :} Le Commandant (brusquerie, indifférence), le Prêtre (hypocrisie, « plus distingué »), Kelara (intérêt matériel), les villageois (opportunistes). \textbf{Procédés :} Ironie, hyperbole, antiphrase, champ lexical dépréciatif.
    \item \textbf{Fonction :} Rendre visible l'aliénation. Le colonisé devient le jouet du colonisateur par sa propre volonté de reconnaissance.
\end{itemize}
\textbf{II. La polémique sous-jacente : accuser le système colonial (L'arme politique).}
\begin{itemize}
    \item \textbf{La médaille comme symbole vide :} « Ferraille », « toc », « piège ». L'administration « donne » pour mieux « prendre » (terres, bras, dignité). Le don est un leurre.
    \item \textbf{La violence symbolique :} La cérémonie est une mascarade (arrivée/départ rapide du Commandant). Le colonisateur ne daigne pas rester.
    \item \textbf{Le basculement vers le tragique :} Le vol final n'est pas un simple vol, c'est la \textbf{révélation} du sens réel : le colonisé ne possède rien, pas même l'honneur qu'on lui « donne ».
\end{itemize}
\textbf{III. L'amertume finale : le rire qui se fige en lucidité (La fable amère).}
\begin{itemize}
    \item \textbf{L'éveil de Meka :} « Meka comprit alors... » (Dernière phrase). Passage du statut d'objet (vieux nègre) à sujet lucide.
    \item \textbf{Fonction cathartique :} Le lecteur rit d'abord, puis ressent la honte/colère. Le rire n'est pas une fin en soi, c'est un \textbf{pédagogie de la révolte}.
    \item \textbf{Portée universelle :} Critique de toute domination qui infantilise l'autre pour l'exploiter.
\end{itemize}
\textbf{Conclusion :}
\begin{itemize}
    \item \textbf{Bilan :} L'affirmation est juste. L'esthétique d'Oyono (satire, oralité, ironie) sert une visée polémique radicale.
    \item \textbf{Ouverture :} Cette « fable amère » préfigure la littérature de la décolonisation (Beti, Mongo Beti, Sembène Ousmane) où l'humour devient résistance.
\end{itemize}
\end{exoresolu}

\begin{attention}[Les 5 erreurs qui coûtent des points au Bac (Résumé \& Analyse)]
    \begin{enumerate}
        \item \textbf{Hors-sujet / Non-respect de la contrainte de longueur :} Résumé trop long ($>$ 110\% de la limite) $\rightarrow$ \textbf{Note 0 ou pénalité lourde} (souvent note sur 20 divisée par 2 au-delà de 10\%). \textbf{Réflexe :} Compter \textbf{avant} la copie propre. Viser le milieu de la fourchette.
        \item \textbf{Copier-coller (Plagiat / Absence de reformulation) :} Reproduire des phrases entières du texte. \textbf{Consigne :} « Reformuler », « Dégager les idées ». \textbf{Sanction :} Note très faible (compétence « Production » non validée). \textbf{Parade :} Changer le vocabulaire (synonymes, hyperonymes), passer au discours indirect, modifier la syntaxe.
        \item \textbf{Ajout d'idées personnelles / Commentaire / Interprétation :} Écrire « \emph{Je pense que Meka est naïf} », « \emph{C'est triste} », « \emph{L'auteur veut dire que...} ». \textbf{Règle :} Le résumé est \textbf{objectif} et \textbf{fidèle} au texte source. Pas de « Je », pas de jugements de valeur.
        \item \textbf{Confusion Constatif / Performatif (Question de langue) :} Dire « \emph{Je déclare} est un constatif car c'est une phrase affirmative ». \textbf{Erreur fatale :} Oublier le critère \textbf{1ère personne + Présent + Verbe d'action + « Par la présente »}. \textbf{Parade :} Appliquer le test systématiquement.
        \item \textbf{Confondre Satire et Polémique / Citer sans analyser (Question de style) :} Dire « \emph{C'est satirique car il y a de l'ironie} » sans citer le passage précis ni expliquer \textbf{le mécanisme} (ex: « L'antiphrase "c'était plus distingué" souligne l'hypocrisie du prêtre »). \textbf{Exigence Bac :} \textbf{Citation} + \textbf{Nom du procédé} + \textbf{Explication de l'effet de sens / visée}.
    \end{enumerate}
\end{attention}

\newpage

\coursec{Exercices}

\courssub{Je vérifie mes connaissances}

\begin{exercice}[QCM : la contraction de texte]
Pour chaque question, une seule réponse est exacte. Entoure la bonne lettre.
\begin{enumerate}
\item Le résumé d'un texte de 400 mots, demandé en un quart du texte, doit compter environ :
\begin{enumerate}
\item[a)] 50 mots \quad b)] 100 mots \quad c)] 200 mots \quad d)] 400 mots
\end{enumerate}
\item Dans un résumé, il est interdit de :
\begin{enumerate}
\item[a)] reformuler les idées de l'auteur \quad b)] donner son avis personnel \quad c)] utiliser des connecteurs logiques \quad d)] employer le présent de narration
\end{enumerate}
\item L'énoncé « Je déclare la séance ouverte » est :
\begin{enumerate}
\item[a)] constatif \quad b)] descriptif \quad c)] performatif \quad d)] interrogatif
\end{enumerate}
\item La tonalité qui consiste à railler les défauts d'une personne ou d'une société pour les dénoncer est la tonalité :
\begin{enumerate}
\item[a)] lyrique \quad b)] satirique \quad c)] tragique \quad d)] épique
\end{enumerate}
\end{enumerate}
\end{exercice}

\begin{exercice}[Vrai ou faux, avec justification]
Dis si chaque affirmation est vraie ou fausse, puis justifie ta réponse en une ou deux phrases.
\begin{enumerate}
\item Un résumé peut contenir des phrases recopiées mot pour mot dans le texte.
\item « La pluie tombe sur le village de Doum » est un énoncé constatif.
\item La tonalité polémique vise à convaincre par le débat et l'attaque argumentée d'une thèse adverse.
\item Dans \emph{Le Vieux Nègre et la Médaille}, la médaille reçue par Meka symbolise la réussite sociale du vieillard.
\end{enumerate}
\end{exercice}

\begin{exercice}[Définitions]
Définis en deux ou trois phrases chacun des termes suivants, en donnant pour chacun un exemple tiré de \emph{Le Vieux Nègre et la Médaille} ou de ta propre invention :
\begin{enumerate}
\item la contraction de texte ;
\item l'énoncé performatif ;
\item la reformulation ;
\item l'ironie (procédé satirique).
\end{enumerate}
\end{exercice}

\begin{exercice}[Calculs de contraction]
Un texte à résumer compte 480 mots.
\begin{enumerate}
\item Calcule le nombre de mots du résumé si la consigne impose le quart du texte.
\item Calcule le nombre de mots si la consigne impose le tiers du texte.
\item Un élève a produit un résumé de 200 mots à partir de ce texte de 480 mots. Exprime en pourcentage la proportion du texte initial que représente ce résumé. Ce résumé respecte-t-il une consigne du quart ? Justifie.
\end{enumerate}
\end{exercice}

\courssub{Je m'entraîne}

\begin{exercice}[Classer les énoncés]
Voici dix énoncés inspirés de l'univers du roman. Recopie et complète le tableau en indiquant pour chacun s'il est \cle{constatif} ou \cle{performatif}, puis justifie ton classement pour les énoncés 3, 6 et 9.
\begin{enumerate}
\item Meka reçoit une médaille des mains du commandant.
\item Je te promets de venir à la cérémonie.
\item Le vent secoue les toits du village.
\item Je déclare la fête terminée.
\item Les anciens du village admirent la médaille.
\item Je vous remercie de votre accueil.
\item La pluie a détruit une partie de la case.
\item Je t'ordonne de sortir immédiatement.
\item Le commandant prononce un long discours.
\item Je félicite le vieux Meka pour sa bravoure.
\end{enumerate}
\begin{center}
\begin{tabularx}{\linewidth}{|c|X|X|}\hline
\rowcolor{popblueL} \textbf{N°} & \textbf{Constatif} & \textbf{Performatif} \\ \hline
1 & & \\ \hline
2 & & \\ \hline
\end{tabularx}
\end{center}
(Prolonge le tableau jusqu'à l'énoncé 10.)
\end{exercice}

\begin{exercice}[Reformulation ciblée]
Reformule chacune des phrases suivantes sans en changer le sens, en appliquant le procédé indiqué entre parenthèses.
\begin{enumerate}
\item Le commandant remet la médaille au vieux Meka. (changement de voix : passif)
\item La foule acclame le vieillard décoré. (nominalisation du verbe « acclamer »)
\item Meka est fier, car il a reçu une distinction. (remplacement du connecteur de cause par un autre)
\item Les villageois admirent la médaille brillante. (synonymie : « admirer », « brillante »)
\item La cérémonie terminée, Meka rentre au village. (transformation de la proposition subordonnée en proposition principale)
\end{enumerate}
\end{exercice}

\begin{exercice}[Analyse stylistique d'un passage]
Lis attentivement le passage suivant, inspiré de la scène de la remise de la médaille : « Le commandant, le ventre tendu dans son uniforme trop étroit, déclama un discours pompeux sur la "grande amitié" entre la France et ses "enfants" d'Afrique, tandis que Meka, écrasé par la chaleur, se demandait à quoi pouvait bien servir ce petit bout de métal accroché à sa poitrine. »
\begin{enumerate}
\item Relève trois procédés satiriques utilisés dans ce passage (ironie, exagération, guillemets moqueurs, antithèse, etc.).
\item Pour chaque procédé relevé, cite le passage exact et explique son effet dénonciateur.
\item Dis en trois lignes si ce passage relève aussi de la tonalité polémique, et pourquoi.
\end{enumerate}
\end{exercice}

\begin{exercice}[Correction d'un résumé fautif]
Voici le résumé produit par un élève à partir d'un extrait de 400 mots du roman (consigne : résumer en 100 mots) : « Le vieux Meka reçoit une médaille. "Le commandant, le ventre tendu dans son uniforme trop étroit, déclama un discours pompeux." Moi je trouve que ce roman est très intéressant et que Ferdinand Oyono écrit très bien. Ensuite il y a une fête au village avec de la nourriture et des boissons et tout le monde danse et chante et Meka est content et ses amis aussi sont contents et ils mangent beaucoup. La médaille est en or massif et vaut une fortune. En conclusion, ce texte parle de la colonisation. »
\begin{enumerate}
\item Relève dans ce résumé six erreurs parmi les suivantes : phrase copiée sur le texte, avis personnel, information inventée (hors-texte), non-respect de la longueur, incohérence de formulation, faute de langue.
\item Pour chaque erreur, cite le passage fautif et nomme le type d'erreur.
\item Propose une version améliorée du résumé, en 100 mots environ, qui respecte les règles de la contraction de texte.
\end{enumerate}
\end{exercice}

\courssub{Je me prépare au Bac}

\begin{exercice}[Résumé guidé type Baccalauréat]
\textbf{Texte support (extrait adapté, 400 mots environ).} « Meka avait attendu ce jour toute sa vie. Le commandant du cercle, suivi de tout son état-major, descendit de sa voiture devant la case du vieillard. Les tambours du village résonnaient, les femmes chantaient, les enfants couraient pieds nus dans la poussière. Le commandant épingla la médaille sur la poitrine osseuse de Meka et prononça un discours interminable où il était question de fidélité, de sacrifice et de la grande famille coloniale. Meka écoutait sans comprendre la moitié des mots, mais il souriait, car on lui avait dit de sourire. Le soir venu, la fête battit son plein : on tua des chèvres, on vida les calebasses de vin de palme. Mais au milieu de la liesse, Meka sentait confusément que cette médaille ne changerait rien à sa vie : sa case resterait en torchis, ses champs resteraient maigres, et demain il reprendrait sa houe comme chaque matin depuis soixante ans. La médaille brillait sur le mur de sa case, inutile et froide, tandis que dehors la pluie commençait à tomber sur le village endormi. »
\begin{enumerate}
\item Lis le texte deux fois, puis complète le tableau d'idées ci-dessous en rédigeant chaque idée essentielle en une phrase de ton cru (sans recopier le texte).
\begin{center}
\begin{tabularx}{\linewidth}{|l|X|}\hline
\rowcolor{popblueL} \textbf{Étapes du récit} & \textbf{Idée essentielle reformulée} \\ \hline
La cérémonie & \\ \hline
Le discours du commandant & \\ \hline
La fête du village & \\ \hline
Les réflexions de Meka & \\ \hline
La conclusion symbolique & \\ \hline
\end{tabularx}
\end{center}
\item Rédige ensuite ton résumé en 100 mots (tolérance : 90 à 110 mots), en reliant les idées par des connecteurs logiques variés.
\item Indique à la fin le nombre exact de mots de ton résumé.
\item Explique en trois lignes quelles suppressions (exemples, répétitions, descriptions secondaires) tu as opérées et pourquoi.
\end{enumerate}
\end{exercice}

\begin{exercice}[Langue et stylistique : sujet type Probatoire]
\textbf{Texte support.} « Je vous décore au nom de la France, déclara le commandant. Meka s'inclina. Les villageois applaudirent, et l'interprète traduisit : "Le grand chef blanc dit que tu es un bon et fidèle serviteur." Le vieillard sourit de toutes ses dents déchaussées, pendant que l'officier, rouge de chaleur, s'essuyait le front et songeait déjà à son déjeuner. »
\begin{enumerate}
\item Relève dans ce passage deux énoncés performatifs et deux énoncés constatifs. Justifie ton classement en t'appuyant sur les verbes employés.
\item Transforme l'énoncé performatif « Je vous décore au nom de la France » en énoncé constatif rapporté au discours indirect.
\item Relève deux procédés satiriques dans la description du commandant et explique leur effet.
\item Montre, en cinq lignes, comment ce passage illustre la distance entre le discours officiel colonial et la réalité vécue par les personnages.
\item Reformule la phrase « Le vieillard sourit de toutes ses dents déchaussées » de deux manières différentes sans en changer le sens.
\end{enumerate}
\end{exercice}

\begin{exercice}[Question de synthèse et résumé intégré type Bac]
\textbf{Situation.} Dans le cadre d'un débat littéraire organisé par le club de lecture de ton lycée technique sur le thème « Littérature et dénonciation de la colonisation », tu dois présenter une intervention écrite.
\begin{enumerate}
\item Situe \emph{Le Vieux Nègre et la Médaille} dans son contexte : auteur, date de publication (1956), pays et période historique évoqués (le Cameroun sous administration coloniale française). (4 lignes maximum)
\item En une quinzaine de lignes, montre comment Ferdinand Oyono dénonce l'illusion coloniale à travers le personnage de Meka : la médaille comme récompense vide, le contraste entre le discours officiel et la misère réelle, l'ironie du récit. Appuie ta démonstration sur deux exemples précis du roman.
\item Résume ta propre démonstration de la question 2 en 50 mots environ, en appliquant les règles de la contraction de texte (pas de citation, pas d'avis nouveau, connecteurs logiques).
\item Conclus ton intervention par un énoncé performatif adressé au public du club de lecture (par exemple une invitation à lire l'œuvre), puis explique pourquoi cet énoncé est performatif.
\end{enumerate}
\end{exercice}

\newpage

\coursec{Corrigés des exercices}

\begin{exemplebox}[Énoncés pour l'exercice 5 — Classification]
Les 10 énoncés à classer :\\
1. Le vent secoue les toits du village.\\
2. Je vous déclare mari et femme.\\
3. La pluie tombe sur le village de Doum.\\
4. Je promets de venir à la cérémonie.\\
5. Le commandant prononce un long discours.\\
6. Je vous remercie de votre accueil.\\
7. Meka reçoit une médaille.\\
8. J'ordonne le silence.\\
9. Les villageois applaudissent.\\
10. Je déclare la fête terminée.
\end{exemplebox}

\begin{corrige}[Exercice 1 — QCM : la contraction de texte]
\begin{enumerate}
\item \textbf{Réponse b) 100 mots.} Le quart de 400 mots se calcule ainsi : $400 \div 4 = 100$ mots. Le résumé doit donc compter environ 100 mots (avec une tolérance habituelle de $\pm 10\,\%$).
\item \textbf{Réponse b) donner son avis personnel.} Le résumé est un texte \cle{objectif} : il restitue fidèlement les idées de l'auteur, sans jugement, sans commentaire, sans opinion du résumiste. Reformuler, utiliser des connecteurs et employer le présent de narration sont au contraire recommandés.
\item \textbf{Réponse c) performatif.} L'énoncé « Je déclare la séance ouverte » \emph{accomplit} l'action qu'il énonce : le simple fait de prononcer cette phrase ouvre réellement la séance. Le verbe « déclarer » est un verbe performatif à la première personne du présent.
\item \textbf{Réponse b) satirique.} La tonalité \cle{satirique} consiste à railler, par l'ironie, l'exagération ou la caricature, les défauts d'un individu ou d'une société afin de les dénoncer et de faire réfléchir le lecteur.
\end{enumerate}
\end{corrige}

\begin{corrige}[Exercice 2 — Vrai ou faux, avec justification]
\begin{enumerate}
\item \textbf{Faux.} Un résumé ne doit contenir \emph{aucune} phrase recopiée mot pour mot : il repose sur la \cle{reformulation} des idées avec ses propres mots. Recopier le texte relève du plagiat et montre que l'élève n'a pas opéré la contraction attendue.
\item \textbf{Vrai.} « La pluie tombe sur le village de Doum » décrit un fait que l'on peut vérifier : l'énoncé peut être déclaré vrai ou faux. C'est la définition même de l'énoncé \cle{constatif}.
\item \textbf{Vrai.} La tonalité \cle{polémique} engage un débat argumenté : elle attaque une thèse adverse et défend une position pour convaincre l'auditoire, souvent avec vigueur.
\item \textbf{Faux.} La médaille symbolise au contraire l'\cle{illusion coloniale} : récompense vide et inutile, elle ne change rien à la misère de Meka. Elle dénonce l'hypocrisie du discours colonial qui honore en paroles sans améliorer la condition réelle des colonisés.
\end{enumerate}
\end{corrige}

\begin{corrige}[Exercice 3 — Définitions]
\begin{enumerate}
\item \textbf{La contraction de texte} est une technique de production écrite qui consiste à réduire un texte à une fraction imposée de sa longueur (quart, tiers\ldots) en n'en conservant que les idées essentielles, reformulées avec ses propres mots. \emph{Exemple :} résumer en 100 mots un extrait de 400 mots du \emph{Vieux Nègre et la Médaille}.
\item \textbf{L'énoncé performatif} est un énoncé qui accomplit l'action qu'il énonce au moment où il est prononcé, généralement grâce à un verbe d'action à la première personne du présent. \emph{Exemple :} « Je te promets de venir à la cérémonie » : dire cette phrase, c'est promettre réellement.
\item \textbf{La reformulation} consiste à exprimer une idée avec d'autres mots que ceux du texte d'origine, sans en changer le sens, grâce à la synonymie, à la nominalisation, au changement de voix ou à la transformation de phrases. \emph{Exemple :} « La foule acclame le vieillard » devient « L'acclamation du vieillard par la foule ».
\item \textbf{L'ironie} est un procédé satirique qui consiste à dire le contraire de ce que l'on pense, ou à présenter comme admirable ce qui est ridicule, afin de railler et de dénoncer. \emph{Exemple :} présenter la médaille de Meka comme une magnifique récompense alors qu'elle ne changera rien à sa pauvreté.
\end{enumerate}
\end{corrige}

\begin{corrige}[Exercice 4 — Calculs de contraction]
\begin{enumerate}
\item \textbf{Le quart du texte :}
\[ 480 \div 4 = 120 \text{ mots.} \]
Le résumé doit compter environ \textbf{120 mots}.
\item \textbf{Le tiers du texte :}
\[ 480 \div 3 = 160 \text{ mots.} \]
Le résumé doit compter environ \textbf{160 mots}.
\item \textbf{Proportion du résumé de 200 mots :}
\[ \frac{200}{480} \times 100 = \frac{20\,000}{480} \approx 41,7\,\%. \]
Le résumé représente environ \textbf{41,7\,\%} du texte initial. Or le quart correspond à $25\,\%$. Comme $41,7\,\% > 25\,\%$, le résumé est \textbf{trop long} : il ne respecte pas la consigne du quart. Il aurait fallu environ 120 mots, soit un écart de 80 mots en trop.
\end{enumerate}
\end{corrige}

\begin{corrige}[Exercice 5 — Classer les énoncés]
\begin{center}
\begin{tabularx}{\linewidth}{|c|X|X|}\hline
\rowcolor{popblueL} \textbf{N°} & \textbf{Constatif} & \textbf{Performatif} \\ \hline
1 & X & \\ \hline
2 & & X \\ \hline
3 & X & \\ \hline
4 & & X \\ \hline
5 & X & \\ \hline
6 & & X \\ \hline
7 & X & \\ \hline
8 & & X \\ \hline
9 & X & \\ \hline
10 & & X \\ \hline
\end{tabularx}
\end{center}
\textbf{Justifications demandées :}
\begin{itemize}
\item \textbf{Énoncé 1} (« Le vent secoue les toits du village ») : constatif, description d'un fait observable, vérifiable (vrai ou faux).
\item \textbf{Énoncé 2} (« Je vous déclare mari et femme ») : performatif, verbe « déclarer » à la 1\iere{} personne du présent, accomplit l'union.
\item \textbf{Énoncé 3} (« La pluie tombe sur le village de Doum ») : constatif, fait météorologique vérifiable.
\item \textbf{Énoncé 4} (« Je promets de venir à la cérémonie ») : performatif, verbe « promettre » à la 1\iere{} personne du présent, engage le locuteur.
\item \textbf{Énoncé 5} (« Le commandant prononce un long discours ») : constatif, rapporte une action à la 3\ieme{} personne, passé simple.
\item \textbf{Énoncé 6} (« Je vous remercie de votre accueil ») : performatif, verbe « remercier » à la 1\iere{} personne du présent, accomplit le remerciement.
\item \textbf{Énoncé 7} (« Meka reçoit une médaille ») : constatif, rapporte un fait (3\ieme{} personne, présent de narration).
\item \textbf{Énoncé 8} (« J'ordonne le silence ») : performatif, verbe « ordonner » à la 1\iere{} personne du présent, l'ordre est donné par l'énonciation.
\item \textbf{Énoncé 9} (« Les villageois applaudissent ») : constatif, fait observable (3\ieme{} personne).
\item \textbf{Énoncé 10} (« Je déclare la fête terminée ») : performatif, verbe « déclarer » à la 1\iere{} personne du présent, clôture l'événement.
\end{itemize}
\end{corrige}

\begin{attention}[Point de vigilance — Performatif vs Constatif]
Un même verbe peut être performatif ou non selon la personne et le temps : « Je promets » (1\iere{} personne, présent) est performatif, mais « Il promet » ou « J'ai promis » sont des constatifs qui rapportent un fait. Seuls les verbes dits « performatifs » (déclarer, promettre, ordonner, nommer, baptiser, jurer, parier\ldots) à la 1\iere{} personne du présent de l'indicatif (ou du subjonctif dans certains cas) produisent l'action.
\end{attention}

\begin{exemplebox}[Passage analysé — Exercice 7 (Extrait du roman, chap. 6)]
« Le commandant, rouge de chaleur, s'essuyait le front du revers de la main et songeait déjà à son déjeuner. Son ventre tendu dans son uniforme trop étroit faisait songer à une outre prête à crever. Il prononça son discours pompeux sur la "grande amitié" franco-africaine, sur la "grande famille" dont Meka et ses "enfants" faisaient partie. [...] Meka, écrasé par la chaleur, ne comprenait pas grand-chose aux paroles de l'officier. Il ne voyait que ce petit bout de métal qui brillait au soleil. »
\end{exemplebox}

\begin{corrige}[Exercice 7 — Analyse stylistique d'un passage]
\begin{enumerate}
\item \textbf{Trois procédés satiriques :}
\begin{itemize}
\item \textbf{La caricature physique (exagération) :} « le ventre tendu dans son uniforme trop étroit », « rouge de chaleur », « faisait songer à une outre prête à crever ». Le commandant est ridiculisé : son corps déborde de l'uniforme symbole de l'autorité coloniale. L'effet dénonciateur est de détruire la majesté du représentant de la France.
\item \textbf{Les guillemets moqueurs (ironie) :} « la "grande amitié" », « la "grande famille" », ses "enfants" d'Afrique. Les guillemets signalent que le narrateur ne croit pas à ces mots : l'amitié et la filiation sont des mensonges du discours colonial, qui masquent la domination.
\item \textbf{L'antithèse / le contraste :} le discours « pompeux » sur la grandeur coloniale s'oppose à la réalité vécue par Meka, « écrasé par la chaleur », pour qui la médaille n'est qu'un « petit bout de métal ». Ce contraste dénonce le fossé entre les belles paroles officielles et l'inutilité concrète de la récompense.
\end{itemize}
(On pouvait aussi relever l'ironie de « petit bout de métal » qui réduit la décoration prestigieuse à un objet insignifiant, ou le détail « songeait déjà à son déjeuner » qui révèle l'indifférence du commandant.)
\item \textbf{Tonalité polémique ?} Ce passage relève partiellement de la polémique : il ne se contente pas de railler, il attaque le discours colonial (la « grande famille » France-Afrique) en opposant frontalement la propagande officielle à la réalité. Cependant, la \cle{satire} domine, car la dénonciation passe surtout par le rire, la caricature et l'ironie plutôt que par une argumentation directe et structurée.
\end{enumerate}
\end{corrige}

\begin{corrige}[Exercice 8 — Correction d'un résumé fautif]
\begin{enumerate}
\item \textbf{Six erreurs relevées :}
\begin{center}
\begin{tabularx}{\linewidth}{|l|X|}\hline
\rowcolor{popblueL} \textbf{Passage fautif} & \textbf{Type d'erreur} \\ \hline
« Le commandant, le ventre tendu\ldots discours pompeux. » & Phrase copiée mot pour mot sur le texte (interdite, même entre guillemets) \\ \hline
« Moi je trouve que ce roman est très intéressant\ldots » & Avis personnel du résumiste (le résumé doit rester objectif) \\ \hline
« La médaille est en or massif et vaut une fortune. » & Information inventée (hors-texte) : rien dans le roman ne l'affirme \\ \hline
Résumé d'environ 120 mots au lieu de 100 & Non-respect de la longueur imposée (consigne : le quart) \\ \hline
« il y a une fête\ldots et tout le monde danse et chante\ldots et ils mangent beaucoup » & Accumulation de détails secondaires et de coordinations « et » (incohérence, absence de hiérarchisation) \\ \hline
« En conclusion, ce texte parle de la colonisation. » & Formule de dissertation hors de propos : un résumé n'a ni introduction ni conclusion personnelle \\ \hline
\end{tabularx}
\end{center}
\item \textbf{Version améliorée (environ 100 mots, consigne : le quart de 400 mots) :}

« Le vieux Meka reçoit une médaille des mains du commandant, qui prononce un discours officiel célébrant l'amitié coloniale. Une fête est ensuite organisée au village : les habitants dansent, chantent et partagent un repas en l'honneur du décoré. Pourtant, derrière les réjouissances, Meka comprend confusément que cette distinction ne changera rien à sa condition : sa pauvreté demeure intacte. La médaille apparaît ainsi comme une récompense vide, symbole de l'illusion entretenue par le discours colonial. »

(\emph{Comptage : 92 mots. Fourchette 90--110 mots respectée. L'essentiel est la hiérarchisation des idées, l'absence de copie, d'avis personnel et d'invention.})
\end{enumerate}
\end{corrige}

\begin{attention}[Points de vigilance — Bac / Probatoire]
Au Bac, toute phrase copiée, tout avis personnel (« je trouve que\ldots », « à mon avis ») ou toute information absente du texte est sanctionné. Respecte scrupuleusement la fourchette de mots annoncée (généralement $\pm 10\,\%$) et indique le nombre de mots à la fin de ta copie (ex. : « Nombre de mots : 98 »).
\end{attention}

\begin{exemplebox}[Extrait support — Exercice 9 et 10 (Cérémonie de remise de la médaille)]
« Le commandant arriva dans sa voiture noire. Il descendit, rouge de chaleur, s'essuya le front. Meka s'inclina. Le commandant lui épingla la médaille sur la poitrine. "Je vous décore au nom de la France", dit-il. L'interprète traduisit : "Le grand chef blanc dit que tu es un bon et fidèle serviteur." Les villageois applaudirent. Le commandant remonta en voiture, songeant déjà à son déjeuner. »
\end{exemplebox}

\begin{corrige}[Exercice 9 — Résumé guidé type Baccalauréat]
\begin{enumerate}
\item \textbf{Tableau d'idées (corrigé possible) :}
\begin{center}
\begin{tabularx}{\linewidth}{|l|X|}\hline
\rowcolor{popblueL} \textbf{Étapes du récit} & \textbf{Idée essentielle reformulée} \\ \hline
La cérémonie & Le commandant vient décorer Meka devant tout le village en fête. \\ \hline
Le discours du commandant & Le discours officiel célèbre la fidélité et la famille coloniale, sans que Meka le comprenne vraiment. \\ \hline
La fête du village & Le village célèbre l'événement par un grand festin. \\ \hline
Les réflexions de Meka & Meka devine que la médaille ne changera ni sa pauvreté ni son labeur quotidien. \\ \hline
La conclusion symbolique & La médaille, inutile, symbolise l'illusion de la récompense coloniale. \\ \hline
\end{tabularx}
\end{center}
\item \textbf{Résumé possible (100 mots $\pm 10\,\%$) :}

« Meka reçoit une médaille des mains du commandant, devant le village rassemblé en fête. Le discours officiel célèbre la fidélité et la famille coloniale, mais le vieillard n'en comprend guère les paroles et se contente de sourire. Le soir, le village festoie en son honneur. Pourtant, au milieu des réjouissances, Meka pressent que cette distinction ne changera rien : sa case, ses champs et son labeur resteront les mêmes. La médaille, inutile et froide, devient ainsi le symbole de l'illusion coloniale. »

\item \textbf{Nombre de mots :} 94 mots (fourchette 90--110 respectée).
\item \textbf{Suppressions opérées :} j'ai supprimé les descriptions secondaires (la voiture noire, les enfants pieds nus, les chèvres et le vin de palme), les répétitions lexicales et les détails pittoresques (la pluie finale), car ils illustrent sans ajouter d'idée nouvelle. Je n'ai gardé que la progression logique : cérémonie $\rightarrow$ discours $\rightarrow$ fête $\rightarrow$ prise de conscience $\rightarrow$ portée symbolique.
\end{enumerate}
\textbf{Barème indicatif (sur 10) :} tableau d'idées : 3 pts ; respect de la longueur et comptage : 2 pts ; reformulation sans copie : 2 pts ; connecteurs et cohérence : 2 pts ; justification des suppressions : 1 pt.
\end{corrige}

\begin{corrige}[Exercice 10 — Langue et stylistique : sujet type Probatoire]
\begin{enumerate}
\item \textbf{Énoncés performatifs (dans l'extrait) :}
\begin{itemize}
\item « Je vous décore au nom de la France » : le verbe « décorer » à la première personne du présent accomplit la décoration par la parole même.
\item « Je déclare la fête terminée » (évoqué plus loin dans le récit / exercice 5) : verbe « déclarer », 1\iere{} personne, présent.
\end{itemize}
\textbf{Énoncés constatifs (dans l'extrait) :} « Meka s'inclina » et « Les villageois applaudirent » (faits vérifiables, passé simple, troisième personne). « Le commandant arriva dans sa voiture noire » (constatif).
\item \textbf{Transformation au discours indirect (constatif) :} Le commandant déclara qu'il décorait Meka au nom de la France. (Le performatif devient constatif : on rapporte la parole, l'action n'est plus accomplie par l'énonciation présente mais rapportée au passé.) \emph{Concordance des temps : présent $\rightarrow$ imparfait.}
\item \textbf{Procédés satiriques (sur le portrait du commandant) :}
\begin{itemize}
\item « rouge de chaleur, s'essuyait le front » : caricature physique qui ridiculise la prétendue supériorité du colonisateur, mal à l'aise sous le climat africain.
\item « songeait déjà à son déjeuner » : ironie qui révèle l'indifférence du commandant : la cérémonie solennelle ne l'intéresse pas, son confort personnel passe avant la « grande amitié » qu'il célèbre.
\item « Son ventre tendu dans son uniforme trop étroit faisait songer à une outre prête à crever » : métaphore caricaturale (outre) qui détruit la dignité du personnage et, par métonymie, celle de l'institution qu'il représente.
\end{itemize}
\item \textbf{Distance entre discours officiel et réalité :} Le commandant proclame solennellement la décoration « au nom de la France » et fait dire par l'interprète que Meka est un « bon et fidèle serviteur » : le discours officiel magnifie la loyauté coloniale. Mais la réalité est tout autre : le vieillard sourit avec ses « dents déchaussées » (image de misère et de décrépitude), tandis que l'officier ne pense qu'à son déjeuner. La grandeur affichée du discours se heurte à la pauvreté du décoré et au désintérêt du décorant : l'écart entre les mots et les choses dénonce l'hypocrisie coloniale.
\item \textbf{Deux reformulations de « Le vieillard souriait, montrant ses dents déchaussées » :}
\begin{itemize}
\item Le vieillard laissa voir, en souriant, ses dents branlantes.
\item Un large sourire éclaira le visage du vieillard aux dents déchaussées.
\end{itemize}
\end{enumerate}
\textbf{Barème indicatif (sur 10) :} question 1 : 2 pts ; question 2 : 1,5 pt ; question 3 : 2 pts ; question 4 : 3 pts ; question 5 : 1,5 pt.
\end{corrige}

\begin{attention}[Points de vigilance — Épreuve de langue]
Pour justifier un performatif, cite toujours la \emph{personne} (1\iere) et le \emph{temps} (présent de l'indicatif) du verbe. Au discours indirect, veille à la concordance des temps (présent $\rightarrow$ imparfait après un verbe introducteur au passé). Pour la satire, nomme le procédé (caricature, ironie, antithèse\ldots) et explique l'effet de sens (dénonciation, ridicule, critique).
\end{attention}

\begin{corrige}[Exercice 11 — Question de synthèse et résumé intégré type Bac]
\begin{enumerate}
\item \textbf{Contexte (4 lignes maximum) :} \emph{Le Vieux Nègre et la Médaille} est un roman de l'écrivain camerounais Ferdinand Oyono, publié en 1956. L'action se déroule au Cameroun, alors sous administration coloniale française, à la veille des indépendances. L'œuvre s'inscrit dans la littérature africaine de dénonciation de la colonisation.
\item \textbf{Développement (15 lignes environ) :}

Dans ce roman, Ferdinand Oyono dénonce l'illusion coloniale en montrant que les honneurs accordés aux colonisés ne sont que des leurres. La médaille remise à Meka, vieillard qui a donné ses terres et ses fils à la France, apparaît comme une récompense vide : elle ne change ni sa pauvreté, ni sa condition, ni son avenir. Premier exemple : lors de la cérémonie, le commandant prononce un discours pompeux sur la fidélité et la grande famille coloniale, mais cette solennité contraste avec la misère réelle de Meka, dont la case reste en torchis et dont les champs demeurent maigres. Le contraste entre les belles paroles officielles et la réalité vécue dénonce l'hypocrisie du système. Deuxième exemple : après la cérémonie, Meka est humilié et arrêté pour s'être égaré en zone européenne, ce qui montre que la médaille ne lui confère aucun droit réel : le « décoré » redevient immédiatement un colonisé méprisé. Enfin, l'ironie du récit, qui présente la médaille comme un honneur suprême alors qu'elle n'est qu'un « petit bout de métal » inutile, invite le lecteur à rire de la supercherie coloniale et à la condamner. Ainsi, par le personnage de Meka, Oyono révèle que la colonisation n'offre aux Africains que des symboles creux en échange de sacrifices bien réels.
\item \textbf{Résumé de la démonstration (50 mots $\pm 10\,\%$) :}

« Oyono dénonce l'illusion coloniale à travers Meka : la médaille, récompense vide, ne change rien à sa misère. Le discours officiel contraste avec la réalité vécue, et l'arrestation du vieillard prouve que l'honneur reçu n'accorde aucun droit. L'ironie du récit achève de ridiculiser la supercherie coloniale. »

(\emph{Comptage : 48 mots. Fourchette 45--55 respectée. Aucune citation, aucun avis nouveau, connecteurs logiques : « à travers », « et », « achève de ».})
\item \textbf{Énoncé performatif de conclusion :} « Je vous invite tous à lire \emph{Le Vieux Nègre et la Médaille} de Ferdinand Oyono. » Cet énoncé est performatif car le verbe « inviter », employé à la première personne du présent, accomplit l'invitation au moment même où la phrase est prononcée : dire, c'est faire.
\end{enumerate}
\textbf{Barème indicatif (sur 20) :} contexte (auteur, date, cadre) : 3 pts ; développement argumenté avec deux exemples précis : 9 pts ; résumé conforme (50 mots, règles respectées) : 5 pts ; énoncé performatif et justification : 3 pts.
\end{corrige}

\begin{attention}[Points de vigilance — Sujet de synthèse]
Dans le résumé de ta propre démonstration (question 3), n'ajoute aucune idée absente du développement : le résumé contracte, il n'enrichit pas. Vérifie le comptage des mots et signale-le explicitement. Pour le contexte (question 1), les trois repères exigibles au Probatoire/Bac sont : \textbf{auteur} (Ferdinand Oyono), \textbf{date} (1956), \textbf{cadre spatio-temporel} (Cameroun, période coloniale française).
\end{attention}

\newpage

\coursec{Fiche bilan}

\coursec{Rédiger un résumé de texte : étude intégrée du \emph{Vieux nègre et la médaille}}

\begin{popfigure}
\begin{tikzpicture}[
  mindmap,
  concept color=popblue,
  level 1 concept/.append style={level distance=3.5cm, sibling angle=60},
  level 2 concept/.append style={level distance=2.5cm},
  every node/.style={font=\small},
  text=white
]
\node[concept, scale=1.2, align=center]{Résumé de texte\\ (\emph{Le Vieux nègre et la médaille})}
  child[concept color=poporange] { node[concept, align=center]{Lecture \& Compréhension\\ (Exposés 1 \& 2)} }
  child[concept color=popgreen] { node[concept, align=center]{Contraction\\ Reformulation idées essentielles} }
  child[concept color=poppurple] { node[concept, align=center]{Langue\\ Constatif / Performatif} }
  child[concept color=poppink] { node[concept, align=center]{Stylistique\\ Tonalités satirique \& polémique} }
  child[concept color=popgold] { node[concept, align=center]{Production\\ Rédaction \& Amélioration} }
  child[concept color=popdark] { node[concept, align=center]{Contexte \& Bilan\\ Intégration} };
\end{tikzpicture}
\legende{Carte mentale de la leçon N22 : Rédiger un résumé de texte}
\end{popfigure}

\courssub{Définitions et notions clés}

\begin{definition}[Résumé (Contraction de texte)]
Exercice de réécriture qui consiste à réduire un texte à ses idées \textbf{essentielles} (environ 1/4 à 1/3 de l'original), en conservant l'ordre logique, le sens et le ton de l'auteur, sans ajouter d'idées personnelles ni de commentaires.
\end{definition}

\begin{definition}[Énoncé constatif]
Énoncé qui \textbf{describe un état de fait} et peut être jugé vrai ou faux (ex: \emph{« Il pleut. »}). Vérité = adéquation au réel.
\end{definition}

\begin{definition}[Énoncé performatif]
Énoncé qui \textbf{réalise l'action} qu'il énonce par le seul fait d'être prononcé dans des conditions appropriées (ex: \emph{« Je vous déclare mari et femme. »}). Conditions de réussite (felicity conditions) : verbe performatif à la 1ère personne du singulier, présent indicatif, contexte institutionnel adéquat, autorité du locuteur.
\end{definition}

\begin{definition}[Tonalité satirique]
Mode d'expression qui \textbf{rit des défauts} (vices, ridicules, travers) pour les corriger par le rire. Procédés principaux : ironie (antiphrase), hyperbole, caricature, parodie, sous-entendu.
\end{definition}

\begin{definition}[Tonalité polémique]
Mode d'expression \textbf{agressif et combatif} qui attaque violemment un adversaire, une idée ou une institution. Procédés : invective, apostrophe, argumentation virulente, champ lexical de la guerre/du combat, accumulation de griefs.
\end{definition}

\courssub{Méthode express : Rédaction du résumé (5 étapes)}

\begin{methode}[Étapes de la contraction]
\textbf{1. Lecture active} : Identifier le thème, la thèse, le plan, le ton, les connecteurs logiques.\\
\textbf{2. Découpage \& Surimposition} : Segmenter en séquences logiques ; surligner les idées forces (mots-clés, phrases-noyaux).\\
\textbf{3. Reformulation (Brouillon)} : Réécrire chaque idée force \emph{avec ses propres mots} (synonymes, changement de structure syntaxique, nominalisation/verbalisation, passage au discours indirect). Supprimer exemples, détails, répétitions, commentaires de l'auteur.\\
\textbf{4. Assemblage \& Liaison} : Relier les phrases reformulées par des connecteurs logiques (donc, cependant, par conséquent, en outre, de plus) pour assurer la cohérence. Respecter l'ordre du texte original.\\
\textbf{5. Vérification \& Comptage} : Contrôler le sens (fidélité), la langue (correction, fluidité), la longueur (nombre de mots $\approx$ 25\text{--}30\% de l'original). Recopier au propre en soignant la présentation.
\end{methode}

\courssub{Repères sur l'œuvre : \emph{Le Vieux nègre et la médaille} (Ferdinand Oyono)}

\begin{savaistu}[Contexte historique et littéraire]
\begin{itemize}
\item \textbf{Auteur / Pays} : Ferdinand Oyono (Cameroun, 1929--2010), figure majeure de la littérature africaine d'expression française.
\item \textbf{Genre} : Roman court / Récit satirique (publié en 1956).
\item \textbf{Contexte} : Cameroun sous administration française (années 1950), période des luttes pour l'indépendance (UPC), critique de la collaboration et de l'assimilation.
\item \textbf{Thèmes majeurs} : Critique du colonialisme, aliénation du colonisé, trahison des élites traditionnelles, illusion des honneurs coloniaux, solitude de la vieillesse, prise de conscience tardive.
\item \textbf{Personnages clés} : Meka (le vieux, figure du collaborateur naïf puis lucide), Kelara (femme lucide), Vandamme (commandant, figure de l'autorité coloniale), le Chef de canton (auxiliaire zélé), le Prêtre (voix de la contestation finale).
\item \textbf{Tonalités dominantes} : \textbf{Satirique} (ironie sur la médaille, naïveté de Meka, description de la cérémonie) \& \textbf{Polémique} (dénonciation de l'ordre colonial, discours final du prêtre, révolte de Meka).
\item \textbf{Structure narrative} : Exposition (attente de la médaille, préparatifs) $\to$ Climax (cérémonie officielle, remise de la médaille) $\to$ Dénouement (vol de la médaille par les siens, désillusion radicale, mort symbolique).
\end{itemize}
\end{savaistu}

\courssub{Distinction Constatif / Performatif (Rappel Bac)}

\begin{propriete}[Tableau comparatif : Constatif vs Performatif]
\begin{center}
\begin{tabularx}{\linewidth}{|l|X|X|}\hline
\rowcolor{popblueL}
\textbf{Critère} & \textbf{Constatif} & \textbf{Performatif} \\ \hline
Fonction & Informer, décrire un état du monde & Agir, transformer le réel par la parole \\ \hline
Évaluation & \textbf{Vrai} / \textbf{Faux} (adéquation au réel) & \textbf{Heureux} / \textbf{Infructueux} (infelicity conditions) \\ \hline
Verbe typique & Verbes d'état, d'action (souvent 3e pers.) & Verbes performatifs (1ère pers. sg, Présent indicatif) : \emph{promettre, jurer, nommer, déclarer, baptiser, ordonner, parier} \\ \hline
Exemple & \emph{« Le ciel est bleu. »} & \emph{« Je promets de venir. »} / \emph{« Je nomme cet enfant Paul. »} \\ \hline
Négation & \emph{« Ce n'est pas vrai. »} & \emph{« Je ne promets pas. »} (refus d'accomplir l'acte) \\ \hline
\end{tabularx}
\end{center}
\end{propriete}

\courssub{Focus stylistique : L'ironie dans \emph{Le Vieux nègre et la médaille}}

\begin{exemplebox}[Procédés ironiques repérés dans l'œuvre]
\begin{itemize}
\item \textbf{Antiphrase} : Dire le contraire de ce que l'on pense (ex: la description flatteuse de la médaille qui vaut « plus que la vie » pour Meka, alors qu'elle symbolise sa servitude).
\item \textbf{Hyperbole} : Exagération grotesque (ex: les préparatifs démesurés pour la cérémonie, l'importance accordée au « costume neuf »).
\item \textbf{Caricature} : Portrait outré des personnages (le Chef de canton servile, le Commandant condescendant).
\item \textbf{Ironie dramatique} : Le lecteur sait (ou devine) la vanité de la médaille avant Meka ; le vol final par ses propres petits-enfants couronne l'absurdité.
\end{itemize}
\end{exemplebox}

\begin{attention}[Pièges à éviter au Probatoire / Bac]
\begin{itemize}
\item \textbf{Confondre résumé et commentaire} : Pas d'analyse, pas de « je pense que », pas d'interprétation personnelle.
\item \textbf{Copier-coller des phrases du texte} : Obligation de reformulation (synonymes, changement de voix, discours indirect).
\item \textbf{Oublier les connecteurs logiques} : Le résumé doit être un texte cohérent, pas une liste d'idées.
\item \textbf{Dépasser la longueur prescrite} : Respecter scrupuleusement le taux de compression (généralement 25\% $\pm$ 5\%).
\item \textbf{Mélanger constatif et performatif} : Un énoncé comme « Je sais qu'il pleut » est constatif (description d'un état mental), pas performatif.
\item \textbf{Confondre satire et polémique} : La satire \emph{rit pour corriger} (distance, humour) ; la polémique \emph{attaque pour détruire} (engagement, violence verbale).
\end{itemize}
\end{attention}

\begin{aretenir}[Checklist avant le Bac - Leçon N22]
$\square$ Je sais définir le résumé et respecter le taux de compression (1/4 -- 1/3).\\
$\square$ Je maîtrise la méthode en 5 étapes (Lecture, Découpage, Reformulation, Assemblage, Vérification).\\
$\square$ Je sais reformuler sans trahir le sens : synonymes, nominalisation/verbalisation, discours indirect.\\
$\square$ Je distingue formellement un énoncé constatif (descriptif, V/F) d'un performatif (actif, 1ère pers. présent).\\
$\square$ Je connais les conditions de réussite d'un performatif (contexte, autorité, formulation).\\
$\square$ J'identifie la tonalité \textbf{satirique} (rire correcteur, ironie) vs \textbf{polémique} (attaque frontale, invective).\\
$\square$ Je relève les procédés stylistiques de l'ironie (antiphrase, hyperboles, caricature) dans \emph{Le Vieux nègre et la médaille}.\\
$\square$ Je situe l'œuvre : contexte colonial camerounais, critique de l'assimilation et de la collaboration.\\
$\square$ Je sais analyser le personnage de Meka (symbole du colonisé aliéné puis lucide).\\
$\square$ Je parviens à produire un résumé cohérent, fidèle, correct linguistiquement et bien présenté.\\
$\square$ Je gère mon temps : 15 min lecture/analyse, 45 min rédaction/brouillon, 15 min copie propre/relecture (épreuve type 1h30).\\
$\square$ Je relis pour éliminer les \textbf{jugements personnels}, les \textbf{citations} inutiles et les \textbf{répétitions}.
\end{aretenir}

\findecours

\end{document}
